\chapter{Optimisation}
\label{vol02:ch15}

\epigraph{To choose well among feasible alternatives is the
engineer's working motion. The mathematics of optimisation is the
discipline of doing it under stated objective and stated
constraint, and of admitting when both have been mis-stated.}{A
working paraphrase of the engineer's register, in the vocabulary
of this volume.}

\chapmeta{Half-life: foundational. Archetypes: optimisation
(formal home; the multivariable apparatus of Volume II Chapter 11
and the one-variable apparatus of Volume II Chapter 5 are the
local cases of what this chapter formalises), scaling.
Exercise target: 40. Project track: simulation with analysis. Hazard
class: none. Reader path default: standard, with sections
explicitly tagged. Named cases: none.}

\noindent
Volume II Chapter 5 introduced optimisation in one variable as
the search for $f'(x) = 0$ on a closed interval. Volume II
Chapter 11 extended it to many variables: $\grad f = \vect{0}$
plus the second-order Hessian test, with equality constraints
handled through Lagrange multipliers and inequality constraints
through Karush-Kuhn-Tucker conditions. This chapter is the
formal home of the optimisation archetype. We define the general
problem, identify the structural class (convex versus nonconvex,
linear versus quadratic versus general nonlinear, continuous
versus integer) that determines what an engineer can hope to
solve, treat duality as the second viewpoint that the apparatus
demands, and develop the working algorithms (gradient descent,
Newton, stochastic gradient) that turn the formal apparatus into
something a finite-precision machine can run on a real dataset.
The chapter closes with the failure modes that even a
well-formulated optimisation can produce: local minima that
look global, objectives that have been gamed by the optimiser,
and the entire genre of optimisation theatre, in which the
appearance of optimisation has been substituted for the
engineering it should have replaced. The canonical reference is
\textcite{text:boyd-vandenberghe2004}; the linear-algebra
machinery beneath every algorithm is in
\textcite{text:strang2016}.

\section{What an optimisation problem is}\pathtag{core}

An optimisation problem has three parts: an objective, a feasible
set, and a decision variable that ranges over the feasible set.
The engineer's job is to find the value of the decision variable
that gives the best objective value, where "best" means smallest
or largest depending on whether the objective is a cost or a
gain.

\subsection{The standard form}

\begin{definition}[Optimisation problem in standard form]
Let $\vect{x} \in \R^{n}$ denote the decision variable. An
\engterm{optimisation problem} in standard form is the search
\[
\min_{\vect{x} \in \R^{n}} \; f_{0}(\vect{x})
\quad \text{subject to} \quad
\begin{cases}
f_{i}(\vect{x}) \leq 0, & i = 1, \ldots, m, \\
h_{j}(\vect{x}) = 0,    & j = 1, \ldots, p.
\end{cases}
\]
The function $f_{0} : \R^{n} \to \R$ is the \engterm{objective}.
The functions $f_{i} : \R^{n} \to \R$ are the \engterm{inequality
constraints}; the functions $h_{j} : \R^{n} \to \R$ are the
\engterm{equality constraints}. The set
\[
\mathcal{F} = \set{\vect{x} \in \R^{n} :
f_{i}(\vect{x}) \leq 0 \;\forall i, \;
h_{j}(\vect{x}) = 0 \;\forall j}
\]
is the \engterm{feasible set}. A point $\vect{x}^{*} \in
\mathcal{F}$ is a \engterm{global optimum} if $f_{0}(\vect{x}^{*})
\leq f_{0}(\vect{x})$ for every $\vect{x} \in \mathcal{F}$, and a
\engterm{local optimum} if the same holds for every $\vect{x} \in
\mathcal{F}$ within some neighbourhood of $\vect{x}^{*}$.
\end{definition}

The standard form fixes a sign convention: minimisation, with
inequalities expressed as "less than or equal to zero." Every
maximisation can be turned into a minimisation by replacing
$f_{0}$ with $-f_{0}$. Every "greater than or equal to" inequality
$g(\vect{x}) \geq 0$ can be rewritten as $-g(\vect{x}) \leq 0$.
The convention has no mathematical content; it lets the rest of
the chapter speak in one voice.

The variable $\vect{x}$ may carry units that do not all cancel:
the design vector for a heat exchanger may contain a tube
diameter in metres and a flow rate in kilograms per second, and
the cost objective may report dollars. The optimiser does not
mind, but the engineer should mind. Mixing units silently is
how a numerical optimiser produces a tube one micrometre wide
and an absurd answer that satisfies every constraint the
formulation actually wrote down.

\subsection{Worked example: the cost-of-pipe problem}

A pipeline of length $L = 100\,\si{\meter}$ must carry a fluid
flow of $Q = 0.05\,\si{\cubic\meter\per\second}$. The pipe
diameter is the design variable $d \in [0.05, 0.30]\,\si{\meter}$.
The material cost per unit length scales as $d^{2}$ (mass of pipe
wall per unit length), and the pumping cost per unit length scales
as $1/d^{5}$ (Darcy-Weisbach pressure loss for fixed $Q$). The
objective is the present-value lifetime cost
\[
f_{0}(d) = c_{1} L d^{2} + c_{2} L Q^{3} d^{-5},
\]
with $c_{1}, c_{2} > 0$ collecting material and pumping
constants. The feasible set is the closed interval $[0.05, 0.30]$.
Volume II Chapter 5's apparatus locates the interior critical
point at $d^{*}$ satisfying $\dv{f_{0}}{d} = 0$, gives the
explicit formula $d^{*} = (5 c_{2} Q^{3} / (2 c_{1}))^{1/7}$, and
verifies $\dv[2]{f_{0}}{d^{2}}(d^{*}) > 0$. If $d^{*}$ falls
inside $[0.05, 0.30]$ the optimum is interior; otherwise the
optimum sits at whichever endpoint gives the smaller cost. The
example is a one-variable, two-bound, smooth-objective problem.
The rest of this chapter generalises it: many variables, many
constraints, possibly nondifferentiable objective, possibly
discrete decision space.

\begin{estimation}
\textbf{Question.} Before working the formula, estimate the optimal
diameter $d^{*}$ for the pipeline above with $c_{1} =
1000\,\text{USD}/\si{\meter\cubed}$ (a representative steel-pipe
material cost, current as of 2025) and $c_{2} = 5 \times 10^{8}
\,\text{USD}\cdot\si{\second\cubed\per\meter\tothe{8}}$ (a
representative pumping cost over a 25-year service life), with $L
= 100\,\si{\meter}$ and $Q = 0.05\,\si{\cubic\meter\per\second}$.
What working diameter would you expect for a low-pressure water
line of this scale, drawn from common sense about plumbing?

\textbf{Calibrated guess.} Domestic-scale pipework runs at
diameters from $25$ to $200\,\si{\milli\meter}$. A $50\,\text{L/s}$
flow is at the upper end of household and the lower end of
commercial scale; a few-hundred-millimetre diameter is the
plausible answer.

\textbf{Calculation.} The first-order condition
$\dv{f_{0}}{d} = 0$ gives $d^{*} = (5 c_{2} Q^{3} / (2 c_{1}))^{1/7}$.
Substituting the numbers, the bracket evaluates to $5 \times 5 \times
10^{8} \times (0.05)^{3} / (2 \times 1000) = 5 \times 5 \times
10^{8} \times 1.25 \times 10^{-4} / 2000 = 3.125 \times 10^{5} /
2000 \approx 156$. Then $d^{*} = 156^{1/7} \approx 2.07
\,\si{\meter}$. The calculation is wildly outside the feasible
range $[0.05, 0.30]\,\si{\meter}$; the constrained optimum is at
the upper endpoint $d = 0.30\,\si{\meter}$. The estimation was
right about the scale (hundreds of millimetres) but the toy
constants $c_{1}, c_{2}$ are not internally calibrated to a real
material-versus-pumping balance at this flow. The exercise is
genuine: the cost coefficients carry units, and an
order-of-magnitude error in either coefficient moves the optimum
by a multiplicative factor in $d^{*}$. The lesson: an
optimisation answer is no better than the cost coefficients it
sits on, and the engineer who reports $d^{*}$ to three decimal
places without checking that the coefficients are consistent with
the actual market and the actual pump curve has reported a
calculation, not an engineering judgment.
\end{estimation}

\subsection{Existence of a minimum}

A minimum may fail to exist. The function $f_{0}(x) = e^{-x}$ on
$\R$ has infimum $0$ but reaches no value: every $x \in \R$ has
$f_{0}(x) > 0$. The function $f_{0}(x) = x$ on $(0, 1)$ has
infimum $0$ but reaches no value because the open interval
excludes its boundary. Two structural conditions guarantee
existence.

\begin{theorem}[Weierstrass]
If the feasible set $\mathcal{F}$ is nonempty, closed, and
bounded, and the objective $f_{0}$ is continuous on $\mathcal{F}$,
then $f_{0}$ attains its minimum and maximum on $\mathcal{F}$.
\end{theorem}

The conditions are tight. A noncompact feasible set, an
unbounded objective, or a discontinuous objective each can
destroy existence. In engineering practice the feasible set is
nearly always closed and bounded by box constraints (every
quantity is bounded by physics, by economics, or by a regulator),
and the objective is nearly always continuous. The Weierstrass
guarantee is the working assumption underneath every numerical
optimiser; it is one of the rare theorems whose hypothesis is
checked by reading the engineering brief.

\subsection{Vocabulary}

\engterm{Active constraint}: a constraint $f_{i}(\vect{x}^{*}) =
0$ at a candidate point. \engterm{Inactive}: $f_{i}(\vect{x}^{*})
< 0$. The active set determines which constraints bind at the
optimum and which are slack. \engterm{Strict feasibility}: a
point $\vect{x}$ satisfying all inequality constraints with
strict inequality, used as the regularity condition for duality.
\engterm{Optimal value}: $p^{*} = f_{0}(\vect{x}^{*})$. The same
problem can have many optima but only one optimal value, by
construction.

\subsection{Structural taxonomy: what determines tractability}

The structural class of an optimisation problem, fixed by the
form of the objective and the constraints, settles in advance
what the engineer can hope to compute. The taxonomy is the first
thing to read off a problem statement, because it selects the
solver, the convergence guarantee, and the runtime scaling
before any number is computed.

\begin{itemize}[leftmargin=*]
  \item \engterm{Linear program} (LP): linear objective, linear
    constraints. Polynomial-time solvable; global optimum at a
    vertex; the workhorse of resource allocation.
  \item \engterm{Quadratic program} (QP): quadratic objective,
    linear constraints. Polynomial-time solvable when the
    objective is convex; least squares, portfolio choice, and
    model predictive control live here.
  \item \engterm{Second-order cone program} (SOCP) and
    \engterm{semidefinite program} (SDP): convex objective over
    cone constraints. Polynomial-time solvable by interior-point
    methods; robust least squares, beamforming, and relaxations
    of hard combinatorial problems live here.
  \item \engterm{General convex program}: convex objective,
    convex feasible set. Polynomial-time solvable to any fixed
    accuracy; every local optimum is global.
  \item \engterm{Smooth nonconvex program}: differentiable but
    not convex. Local optima are computable; the global optimum
    has no efficient guarantee. Neural-network training lives
    here.
  \item \engterm{Integer and combinatorial program}: discrete
    decision variables. NP-hard in general; solved by relaxation,
    branch-and-bound, or approximation. Scheduling, routing, and
    facility location live here.
\end{itemize}

The dividing line that matters most is not linear-versus-nonlinear
but convex-versus-nonconvex. A convex nonlinear problem (a convex
QP, an SOCP, an SDP) is tractable; a nonconvex problem, even one
with a simple-looking objective, is generally not. The chapter
develops the convex case in full because it is the case with a
working theory, and treats the nonconvex case with the honesty
its difficulty demands.

\subsection{Scaling: how cost grows with problem size}

The optimisation archetype carries a scaling archetype inside it.
The engineer who chooses a method must read its per-iteration cost
and its iteration count as functions of the dimension $n$, the
number of constraints $m$, and the data count $N$.

\begin{archetype}[Scaling of optimisation cost]
The total cost of an iterative optimiser is the product of two
factors: the per-iteration cost and the number of iterations.
The per-iteration cost is set by the linear algebra: a
gradient evaluation is $O(n)$ to $O(n^{2})$, a Hessian solve is
$O(n^{3})$ in dense form, a single LP-relaxation solve inside a
branch-and-bound node is polynomial in $(n, m)$. The iteration
count is set by the convergence rate: geometric (gradient descent,
rate governed by the condition number $\kappa$), superlinear
(quasi-Newton), quadratic (Newton near the optimum), or
sublinear $O(1/\sqrt{K})$ (stochastic gradient). The two factors
trade against each other. Newton's method needs few iterations
but pays $O(n^{3})$ per step; gradient descent needs many
iterations but pays $O(n)$ per step. The break-even moves with
$n$: second-order methods dominate at small $n$, first-order
methods dominate at large $n$, and the crossover is the
engineer's working judgment. Every method-selection decision in
this chapter is an instance of this scaling archetype.
\end{archetype}

The scaling archetype recurs in the estimation block of
Section 15.5 (gradient descent versus Newton at $n = 10^{4}$ and
$n = 10^{5}$), in the SGD treatment of Section 15.6 (full gradient
versus minibatch at $N = 10^{6}$), and in the branch-and-bound
estimate of Section 15.7 ($2^{30}$ nodes pruned to $10^{7}$). The
single habit it installs: before running an optimiser, multiply
the per-iteration cost by the expected iteration count, and check
that the product fits the compute budget.

\section{Convex optimisation}\pathtag{core}

The single most important structural distinction in optimisation
is convex versus nonconvex. Convex problems have a working
theory: every local minimum is a global minimum, the optimality
conditions are computable, and a wide family of algorithms
converge under stated step rules. Nonconvex problems have no
such guarantee. Most engineering optimisation problems either
are convex or can be approximated by a convex problem at the
working point.

\subsection{Convex sets}

\begin{definition}[Convex set]
A set $C \subseteq \R^{n}$ is \engterm{convex} if for every pair
$\vect{x}, \vect{y} \in C$ and every $\theta \in [0, 1]$, the
point $\theta \vect{x} + (1 - \theta) \vect{y}$ also lies in $C$.
\end{definition}

Geometrically: every line segment between two points of $C$ stays
inside $C$. The half-space $\set{\vect{x} : \vect{a}^{\transpose}
\vect{x} \leq b}$ is convex; the ball $\set{\vect{x} :
\norm{\vect{x} - \vect{c}} \leq r}$ is convex; the intersection of
convex sets is convex; an affine map of a convex set is convex. A
\engterm{polyhedron} is the intersection of finitely many
half-spaces, hence convex. The feasible set defined by linear
inequalities and equalities is always polyhedral, hence convex.

The standard non-examples: the union of two disjoint balls (the
segment between centres leaves the union); the set $\set{(x, y) :
xy = 1, x > 0}$ (the segment between $(1, 1)$ and $(2, 1/2)$ does
not stay on the curve); the integer lattice (the segment between
two integer points contains non-integer points). The third
non-example is the entire reason Section 15.7 exists.

The convex sets an engineer meets most often are built from a few
primitives. A \engterm{hyperplane} $\set{\vect{x} :
\vect{a}^{\transpose} \vect{x} = b}$ and the two \engterm{half-
spaces} it bounds are convex. A \engterm{norm ball}
$\set{\vect{x} : \norm{\vect{x} - \vect{c}} \leq r}$ is convex in
any norm. An \engterm{ellipsoid} $\set{\vect{x} :
(\vect{x} - \vect{c})^{\transpose} P^{-1} (\vect{x} - \vect{c})
\leq 1}$ with $P \succ 0$ is convex; it is the image of a ball
under an affine map, and the principal axes are the eigenvectors
of $P$ scaled by the square roots of its eigenvalues. The
\engterm{probability simplex} $\set{\vect{x} : \vect{x} \geq
\vect{0}, \vect{1}^{\transpose} \vect{x} = 1}$ is the intersection
of the nonnegative orthant with an affine hyperplane, hence convex,
and it is the feasible set of every problem whose decision variable
is a distribution. The \engterm{positive semidefinite cone}, the
set of symmetric matrices with nonnegative eigenvalues, is convex
and is the feasible set of a semidefinite program. Each is closed
under intersection with the others, and the engineer assembles a
feasible set by intersecting the primitives the constraints
provide.

\paragraph{The separating-hyperplane theorem.}
The structural fact that makes convex sets tractable is that two
disjoint convex sets can always be separated by a hyperplane. If
$C$ and $D$ are nonempty disjoint convex sets, there exist
$\vect{a} \neq \vect{0}$ and $b$ with $\vect{a}^{\transpose}
\vect{x} \leq b$ for all $\vect{x} \in C$ and $\vect{a}^{\transpose}
\vect{x} \geq b$ for all $\vect{x} \in D$. The separating
hyperplane is the geometric content behind duality: the dual
variables are the normal vector of a hyperplane that separates the
feasible set from the sublevel set of the objective at the optimal
value. A supporting hyperplane at a boundary point of a convex set
is the limiting case where $D$ shrinks to that single point, and
it is the geometric content behind the tangent characterisation of
convex functions. Failure of separation is the signature of
nonconvexity: two interleaved nonconvex sets admit no separating
hyperplane, and the duality theory built on separation loses its
guarantee.

\subsection{Convex functions}

\begin{definition}[Convex function]
A function $f : \R^{n} \to \R$, defined on a convex domain $C$, is
\engterm{convex} if for every $\vect{x}, \vect{y} \in C$ and every
$\theta \in [0, 1]$,
\[
f(\theta \vect{x} + (1 - \theta) \vect{y}) \leq
\theta f(\vect{x}) + (1 - \theta) f(\vect{y}).
\]
The function is \engterm{strictly convex} if the inequality is
strict for $\vect{x} \neq \vect{y}$ and $\theta \in (0, 1)$. The
function is \engterm{concave} if $-f$ is convex.
\end{definition}

The graph of a convex function lies below every chord; the chord
is the right-hand side of the inequality, and the curve is the
left. The graph of a concave function lies above every chord.
Linear functions are both convex and concave, the boundary
between the two classes.

For a once-differentiable function, convexity has a tangent
characterisation: $f(\vect{y}) \geq f(\vect{x}) + \grad
f(\vect{x})^{\transpose} (\vect{y} - \vect{x})$ for all
$\vect{x}, \vect{y}$ in the domain. The tangent plane at any
point lies entirely below the graph. For a twice-differentiable
function, convexity has a Hessian characterisation: the Hessian
$\nabla^{2} f$ is positive semidefinite at every point of the
domain. We met the Hessian as the matrix of second partials in
Volume II Chapter 11; we met positive semidefiniteness as a
property of the Gram matrix in Volume II Chapter 9. A
twice-differentiable $f$ is strictly convex when the Hessian is
positive definite (every eigenvalue strictly positive).

The library of convex functions an engineer should recognise on
sight:

\begin{itemize}[leftmargin=*]
  \item Affine functions $f(\vect{x}) = \vect{a}^{\transpose}
    \vect{x} + b$. Both convex and concave.
  \item Quadratic forms $f(\vect{x}) = \vect{x}^{\transpose} P
    \vect{x} + \vect{q}^{\transpose} \vect{x} + r$ with $P$
    positive semidefinite. Convex (strictly so when $P$ is
    positive definite).
  \item Norms $\norm{\vect{x}}$, in any norm. Convex by the
    triangle inequality.
  \item The exponential $e^{a x}$ on $\R$, for any $a$. Convex.
  \item The negative logarithm $-\log x$ on $(0, \infty)$.
    Strictly convex.
  \item The negative entropy $\sum_{i} x_{i} \log x_{i}$ on the
    probability simplex. Strictly convex.
  \item The maximum of a finite collection of convex functions,
    $\max_{i} f_{i}(\vect{x})$. Convex (the pointwise maximum
    of convex functions is convex).
  \item The composition $g(A \vect{x} + \vect{b})$ of a convex
    $g$ with an affine map. Convex.
\end{itemize}

The library is closed under several operations the engineer
relies on: nonnegative weighted sums, pointwise maxima,
compositions with affine maps, perspective. Boyd and
Vandenberghe's Chapter 3 treats this calculus; the working point
is that an engineer who has learned the library can recognise a
convex objective from inspection without computing the Hessian.

\paragraph{Worked example: convexity by Hessian.}
Consider $f(x, y) = e^{x} + e^{y} + (x - y)^{2}$ on $\R^{2}$. The
gradient is $\grad f = (e^{x} + 2(x - y), \; e^{y} - 2(x - y))$.
The Hessian is
\[
\nabla^{2} f =
\begin{pmatrix}
e^{x} + 2 & -2 \\
-2 & e^{y} + 2
\end{pmatrix}.
\]
The diagonal entries are positive. The determinant is
$(e^{x} + 2)(e^{y} + 2) - 4 = e^{x + y} + 2 e^{x} + 2 e^{y} + 4 -
4 = e^{x+y} + 2 e^{x} + 2 e^{y} > 0$ everywhere. A symmetric
$2 \times 2$ matrix with positive diagonal and positive
determinant is positive definite (its leading principal minors
are positive, the Sylvester criterion of Volume II Chapter 9).
The function is strictly convex on all of $\R^{2}$, so the unique
critical point, wherever the gradient vanishes, is the global
minimum. The recognition route reaches the same conclusion
without arithmetic: $e^{x}$ and $e^{y}$ are convex from the
library, $(x - y)^{2}$ is a convex quadratic (it is
$\norm{A \vect{x}}^{2}$ with $A = (1, -1)$), and a nonnegative sum
of convex functions is convex. The Hessian computation confirms
strictness.

\paragraph{Worked example: a function that fails convexity.}
Consider $g(x, y) = x^{2} - y^{2}$, the standard saddle. The
Hessian is $\diag(2, -2)$, with eigenvalues $2$ and $-2$. The
mixed sign of the eigenvalues means the Hessian is indefinite:
positive curvature along $x$, negative along $y$. The function is
neither convex nor concave. Its single critical point $(0, 0)$ is
a saddle, not an extremum; gradient descent initialised off the
$y = 0$ axis runs away to $-\infty$ along $y$. Saddle points are
the characteristic obstruction in nonconvex optimisation, and
high-dimensional nonconvex objectives have far more saddle points
than local minima, a fact that shapes how first- and second-order
methods behave on deep-learning losses.

\begin{mastery}
\textbf{Jensen's inequality and the probabilistic reading of
convexity.} The defining inequality of a convex function extends
from two points to any finite convex combination by induction,
and then to any probability distribution by a limiting argument:
for a convex $f$ and a random variable $X$ with finite mean,
\[
f(\E[X]) \leq \E[f(X)].
\]
This is \engterm{Jensen's inequality}. It is the engine behind a
remarkable number of estimates. The arithmetic-mean / geometric-
mean inequality is Jensen applied to $f = -\log$: $-\log(\E[X])
\leq \E[-\log X]$, i.e.\ $\E[\log X] \leq \log \E[X]$, which
rearranges to the AM-GM bound. The variance is nonnegative because
$x \mapsto x^{2}$ is convex: $(\E[X])^{2} \leq \E[X^{2}]$. The
log-likelihood lower bound used in the expectation-maximisation
algorithm and the evidence lower bound used in variational
inference are both Jensen applied to $\log$. An engineer who reads
convexity probabilistically gains a single tool that produces
order-of-magnitude bounds across estimation, information theory,
and machine learning. The cost is one habit: notice when a
quantity of interest is the expectation of a convex (or concave)
function of a random input, and read off the inequality.
\end{mastery}

\subsection{Convex optimisation problems}

\begin{definition}[Convex optimisation problem]
The problem
\[
\min_{\vect{x}} f_{0}(\vect{x}) \quad \text{subject to} \quad
f_{i}(\vect{x}) \leq 0, \;\; h_{j}(\vect{x}) = 0
\]
is \engterm{convex} when $f_{0}$ and every $f_{i}$ are convex
functions and every $h_{j}$ is affine.
\end{definition}

The asymmetry between inequality and equality constraints is not
arbitrary. A convex set defined by $f_{i}(\vect{x}) \leq 0$ for
convex $f_{i}$ stays convex (the sublevel set of a convex function
is convex). A set defined by $h(\vect{x}) = 0$ for nonlinear $h$
is generally not convex (a curved equality surface fails the
segment test). Equality constraints in a convex problem must
therefore be affine: $A \vect{x} = \vect{b}$.

\begin{archetype}[Optimisation, formal home]
The optimisation archetype was introduced in Volume II Chapter 5
as the one-variable case ($f'(x) = 0$ at an interior optimum,
plus endpoint check), developed in Volume II Chapter 11 as the
multivariable case (gradient zero, Hessian definiteness, Lagrange
multipliers for equality constraints, Karush-Kuhn-Tucker
conditions for inequality constraints), and now finds its formal
home. We have a class of problems (convex programs) for which
"optimum" is a single global object accessible by computable
methods, a duality theory that supplies a second viewpoint and a
certificate of suboptimality, and a working algorithm catalogue
(gradient methods, Newton's method, interior-point methods)
whose convergence is theorem-grade for convex problems and
heuristic for nonconvex problems. The engineer's first
diagnostic question on receiving an optimisation problem is:
\engterm{is it convex}? If yes, the apparatus of this chapter
applies in full. If no, the apparatus applies locally, and the
engineer must additionally consider initialisation, multiple
restarts, and the failure mode of Section 15.8.
\end{archetype}

\subsection{Why convexity matters: the local-implies-global property}

\begin{theorem}[Local optima of convex problems are global]
Let $f_{0}$ be a convex function on a convex feasible set
$\mathcal{F}$. If $\vect{x}^{*}$ is a local optimum of $f_{0}$ on
$\mathcal{F}$, then $\vect{x}^{*}$ is a global optimum.
\end{theorem}

\begin{proof}[Sketch]
Suppose $\vect{x}^{*}$ is a local optimum and $\vect{y} \in
\mathcal{F}$ has $f_{0}(\vect{y}) < f_{0}(\vect{x}^{*})$. Consider
the segment $\vect{x}_{\theta} = \theta \vect{y} + (1 - \theta)
\vect{x}^{*}$ for $\theta \in (0, 1)$, which stays in
$\mathcal{F}$ by convexity of the feasible set. Convexity of
$f_{0}$ gives $f_{0}(\vect{x}_{\theta}) \leq \theta f_{0}(\vect{y})
+ (1 - \theta) f_{0}(\vect{x}^{*}) < f_{0}(\vect{x}^{*})$ for
$\theta > 0$. For $\theta$ small enough $\vect{x}_{\theta}$ lies
inside the local-optimum neighbourhood of $\vect{x}^{*}$, which
contradicts $\vect{x}^{*}$ being a local optimum. Hence no such
$\vect{y}$ exists; $\vect{x}^{*}$ is global.
\end{proof}

The theorem is the load-bearing fact of the chapter. Every
algorithm we treat for convex problems converges to a local
optimum; the theorem upgrades that local guarantee to a global
guarantee. For nonconvex problems the same algorithm finds only
a local optimum, with no statement about the global picture
(\cref{fig:vol02:ch15:convex-nonconvex}).

\begin{figure}[ht]
\centering
\begin{tikzpicture}[x=1cm, y=1cm, font=\small]

% Left panel: convex objective with a single global minimum
\begin{scope}[shift={(0,0)}]
  \node[anchor=south] at (3, 3.4) {convex};
  \draw[->, thick] (-0.2, 0) -- (6.2, 0) node[right, font=\scriptsize] {$x$};
  \draw[->, thick] (0, -0.1) -- (0, 3.4) node[above, font=\scriptsize] {$f(x)$};
  \draw[thick, engblue, smooth, samples=80, domain=0.2:5.8]
    plot (\x, {0.3*(\x-3)*(\x-3) + 0.4});
  \fill[engred] (3, 0.4) circle (2.5pt);
  \node[engred, anchor=west, font=\scriptsize] at (3.15, 0.4)
    {global $= $ local};
  \draw[gray, dashed] (3, 0) -- (3, 0.4);
\end{scope}

% Right panel: nonconvex objective with several local minima
\begin{scope}[shift={(7.5,0)}]
  \node[anchor=south] at (3, 3.4) {nonconvex};
  \draw[->, thick] (-0.2, 0) -- (6.2, 0) node[right, font=\scriptsize] {$x$};
  \draw[->, thick] (0, -0.1) -- (0, 3.4) node[above, font=\scriptsize] {$f(x)$};
  % Three-basin curve, plotted as a sum of shifted gaussians offset by a
  % gentle slope so the basins have different depths.
  \draw[thick, engblue, smooth, samples=120, domain=0.2:5.8]
    plot (\x, {2.6
      - 1.4*exp(-((\x-1.2)*(\x-1.2))/0.3)
      - 2.0*exp(-((\x-3.0)*(\x-3.0))/0.25)
      - 1.1*exp(-((\x-4.7)*(\x-4.7))/0.35)});
  % Three local minima, only the middle one is global
  \fill[engred] (3.0, 0.6) circle (2.5pt);
  \node[engred, anchor=west, font=\scriptsize] at (3.15, 0.6) {global};
  \fill[engblack] (1.2, 1.2) circle (2pt);
  \node[engblack, anchor=south, font=\scriptsize] at (1.2, 1.3) {local};
  \fill[engblack] (4.7, 1.5) circle (2pt);
  \node[engblack, anchor=south, font=\scriptsize] at (4.7, 1.6) {local};
\end{scope}

\end{tikzpicture}
\caption[Convex versus nonconvex objectives]{Convex versus
nonconvex objectives in one variable. The convex case (left) has a
single critical point that is both the local minimum and the global
minimum; gradient descent from any starting point reaches it. The
nonconvex case (right) has three local minima of different depth;
gradient descent converges to whichever basin contains the starting
point, with no local diagnostic distinguishing the global minimum
from the two shallower local minima. The local-implies-global
theorem of Section 15.2 separates the two regimes and is the
reason convex modelling is preferred whenever the engineering
problem admits a convex formulation.}
\label{fig:vol02:ch15:convex-nonconvex}
\end{figure}


\subsection{Recognising convexity in engineering problems}

A least-squares problem $\min_{\vect{x}} \norm{A \vect{x} -
\vect{b}}_{2}^{2}$ is convex (the squared norm of an affine
function is a convex quadratic). A linear program is convex (the
objective is linear, hence convex; the constraints are linear,
hence affine equalities and convex sublevel sets). The problem of
fitting a logistic regression on linearly separable data is
convex. The problem of training a neural network is generally
not convex.

A practical heuristic: if the problem can be written using only
the convex-function library above and the convexity-preserving
operations, the problem is convex; otherwise the engineer
should suspect nonconvexity and apply Section 15.8's diagnostics.

\subsection{Least squares as optimisation}

The most-solved optimisation problem in all of engineering is
least squares, and it is worth seeing it as an optimisation
problem rather than as a linear-algebra recipe. Given $A \in
\R^{m \times n}$ with $m \geq n$ and $\vect{b} \in \R^{m}$, the
problem
\[
\min_{\vect{x} \in \R^{n}} f(\vect{x}) = \tfrac{1}{2}
\norm{A \vect{x} - \vect{b}}_{2}^{2}
\]
fits an overdetermined linear system in the sense that minimises
the sum of squared residuals. The objective is a convex quadratic:
expanding, $f(\vect{x}) = \tfrac{1}{2} \vect{x}^{\transpose}
(A^{\transpose} A) \vect{x} - (A^{\transpose} \vect{b})^{\transpose}
\vect{x} + \tfrac{1}{2} \norm{\vect{b}}^{2}$, with Hessian
$A^{\transpose} A \succeq 0$. The gradient is $\grad f =
A^{\transpose}(A \vect{x} - \vect{b})$, and setting it to zero
gives the \engterm{normal equations}
\[
A^{\transpose} A \vect{x} = A^{\transpose} \vect{b}.
\]
When $A$ has full column rank, $A^{\transpose} A$ is positive
definite, $f$ is strictly convex, and the solution $\vect{x}^{*}
= (A^{\transpose} A)^{-1} A^{\transpose} \vect{b}$ is the unique
global minimum. The optimisation reading and the geometric reading
of Volume II Chapter 9 are the same fact: the residual $A
\vect{x}^{*} - \vect{b}$ is orthogonal to the column space of $A$,
which is exactly the stationarity condition $A^{\transpose}(A
\vect{x}^{*} - \vect{b}) = \vect{0}$.

The conditioning of least squares is the conditioning of
optimisation. The condition number of the Hessian $A^{\transpose}
A$ is $\kappa(A^{\transpose} A) = \kappa(A)^{2}$, the square of the
condition number of $A$. Forming the normal equations squares the
conditioning and can destroy accuracy in finite precision; the
numerically careful route solves the least-squares problem through
the QR factorisation or the singular value decomposition of $A$
directly, never forming $A^{\transpose} A$. The same conditioning
$\kappa(A)^{2}$ governs how slowly gradient descent converges on
the least-squares objective, which is why the estimation block in
Section 15.5 and the ill-conditioned least-squares estimation
exercise both turn on this square.
\Cref{fig:vol02:ch15:contour-descent} draws the gradient-descent
path on a least-squares contour.

\begin{figure}[ht]
\centering
\begin{tikzpicture}[x=1cm, y=1cm, font=\small]

% Mildly elongated quadratic f(x,y) = 1/2 (a x^2 + b y^2) with a=1, b=6.
% Level set f = c is the ellipse x^2/(2c) + y^2/(2c/6) = 1,
% semi-axes (sqrt(2c), sqrt(c/3)). Centre at (3, 2) so axes are visible.
\def\cx{3}
\def\cy{2}

% Axes
\draw[->, thick] (-0.2, 0) -- (6.4, 0) node[right, font=\scriptsize] {$x_{1}$};
\draw[->, thick] (0, -0.2) -- (0, 4.2) node[above, font=\scriptsize] {$x_{2}$};

% Elliptic contours centred at the optimum (cx, cy)
\foreach \c in {0.3, 0.9, 1.8, 3.0, 4.5} {
  \pgfmathsetmacro{\axA}{sqrt(2*\c)}
  \pgfmathsetmacro{\axB}{sqrt(\c/3)}
  \draw[engblue, thin] (\cx, \cy) ellipse ({\axA} and {\axB});
}

% Optimum
\fill[engred] (\cx, \cy) circle (2pt);
\node[engred, anchor=west, font=\scriptsize] at (\cx+0.12, \cy) {$\vect{x}^{*}$};

% Gradient-descent iterates with a modest step size: smooth descent
% from x_0 = (0.4, 3.6) curving into the optimum, with the
% characteristic perpendicular-to-contour first steps.
\fill[engblack] (0.40, 3.60) circle (1.5pt);
\fill[engblack] (1.30, 2.55) circle (1.5pt);
\fill[engblack] (1.95, 2.25) circle (1.5pt);
\fill[engblack] (2.45, 2.10) circle (1.5pt);
\fill[engblack] (2.78, 2.04) circle (1.5pt);
\fill[engblack] (2.95, 2.01) circle (1.5pt);

\draw[engblack, thick]
  (0.40, 3.60) -- (1.30, 2.55) -- (1.95, 2.25) -- (2.45, 2.10)
  -- (2.78, 2.04) -- (2.95, 2.01);

\node[engblack, anchor=south east, font=\scriptsize] at (0.40, 3.62) {$\vect{x}_{0}$};

% Gradient direction at the first point (negative gradient = descent step)
\draw[->, engred, thick] (0.40, 3.60) -- (0.85, 3.07);
\node[engred, anchor=west, font=\scriptsize] at (0.88, 3.20)
  {$-\grad f(\vect{x}_{0})$};

% Caption note
\node[anchor=north, font=\scriptsize, align=center] at (3.1, -0.5)
  {$f(\vect{x}) = \tfrac{1}{2}\norm{A\vect{x} - \vect{b}}_{2}^{2}$,
   $\kappa = 6$};

\end{tikzpicture}
\caption[Gradient descent on a least-squares contour]{Gradient
descent on the convex quadratic of a least-squares objective,
plotted on its elliptic level sets. Each step moves along the
negative gradient, which at every point is orthogonal to the local
level set. With a well-chosen step size on a mildly conditioned
problem ($\kappa = 6$), the iterates curve steadily into the
optimum $\vect{x}^{*}$ without the zigzag of the badly conditioned
case in \cref{fig:vol02:ch15:gradient-zigzag}. The optimum is the
point where the residual $A\vect{x}^{*} - \vect{b}$ is orthogonal
to the column space of $A$, the geometric reading of the
stationarity condition $A^{\transpose}(A\vect{x}^{*} - \vect{b}) =
\vect{0}$.}
\label{fig:vol02:ch15:contour-descent}
\end{figure}


\section{Linear and quadratic programming}\pathtag{standard}

Two structural classes deserve named treatment. They cover most
of what an engineer encounters in resource allocation, design,
and convex modelling.

\subsection{Linear programming}

\begin{definition}[Linear program in standard form]
A \engterm{linear program} (LP) in standard form is
\[
\min_{\vect{x}} \; \vect{c}^{\transpose} \vect{x}
\quad \text{subject to} \quad
A \vect{x} = \vect{b}, \quad \vect{x} \geq \vect{0},
\]
with $\vect{c} \in \R^{n}$, $A \in \R^{m \times n}$, $\vect{b}
\in \R^{m}$, and $\vect{x} \geq \vect{0}$ understood
componentwise.
\end{definition}

An LP has a linear objective and linear (equality plus
nonnegativity) constraints. The feasible set is a polyhedron;
the objective is constant on hyperplanes. The geometry forces
the optimum to occur at a vertex of the polyhedron whenever an
optimum exists and the polyhedron has vertices. There are
finitely many vertices, so the search is, in principle, finite.
The simplex method of \textcite{text:v2ch15-dantzig1947} walks the
vertex graph along edges that decrease the objective; the
interior-point methods of Karmarkar (1984) and successors approach
the optimum through the polyhedron's interior in polynomial time. An
engineer in 2026 does not implement either method by hand; the
working tools are commercial solvers (Gurobi, CPLEX, Mosek) and
open-source solvers (HiGHS, GLPK, CBC, Clp), called from a
modelling layer (CVXPY, JuMP, PuLP, AMPL) that translates the
problem statement into the solver's input format.

The LP formulation covers a wide range of engineering
problems: scheduling, transportation, blending, network flow,
production planning, diet design. Its limitation is the
linearity assumption. Real cost curves are seldom exactly linear;
linear programs are nonetheless the working approximation when
costs are roughly linear in the relevant range, and they are the
relaxation that integer programs use to bound their search
(Section 15.7). The two-variable geometry, with feasible polygon
and parallel objective level sets, generalises directly to $n$
variables with a polyhedron of exponentially many vertices
(\cref{fig:vol02:ch15:feasible-polytope}).

\begin{figure}[ht]
\centering
\begin{tikzpicture}[x=1cm, y=1cm, font=\small]

% A two-variable LP feasible region.
%   x_1, x_2 >= 0,
%   x_1 + 2 x_2 <= 6,
%   3 x_1 + 2 x_2 <= 12.
% Vertices: (0,0), (4,0), (3, 3/2), (0, 3). Pentagon truncated by axes.

% Axes
\draw[->, thick] (-0.3, 0) -- (5.5, 0) node[right, font=\scriptsize] {$x_{1}$};
\draw[->, thick] (0, -0.3) -- (0, 4.0) node[above, font=\scriptsize] {$x_{2}$};

% Feasible region (filled)
\fill[engblue!20]
  (0, 0) -- (4, 0) -- (3, 1.5) -- (0, 3) -- cycle;

% Boundary lines (extended a bit for clarity)
\draw[engblue, thick] (0, 3) -- (6, 0) coordinate (A);
\node[engblue, anchor=west, font=\scriptsize] at (4.2, 1.0)
  {$x_{1} + 2 x_{2} = 6$};
\draw[engred, thick] (0, 6) -- (4, 0) coordinate (B);
\node[engred, anchor=west, font=\scriptsize] at (3.6, 2.5)
  {$3 x_{1} + 2 x_{2} = 12$};

% Vertices
\foreach \v/\lbl/\pos in {%
  (0,0)/$v_{1}$/below left,
  (4,0)/$v_{2}$/below,
  (3,1.5)/$v_{3}$/right,
  (0,3)/$v_{4}$/left%
} {
  \fill[engblack] \v circle (1.8pt);
  \node[font=\scriptsize, \pos] at \v {\lbl};
}

% Objective-direction arrow (maximise c^T x with c = (1, 1))
\draw[->, thick, gray] (3, 3) -- (4.0, 4.0);
\node[gray, anchor=west, font=\scriptsize, align=left] at (4.0, 4.0)
  {direction of\\ increasing $\vect{c}^{\transpose} \vect{x}$};

% Level sets of objective (parallel dashed lines)
\foreach \k in {2, 3.5, 4.5} {
  \draw[gray, dashed]
    ({\k - 0}, 0) -- (0, {\k - 0});
}

% Mark v_3 as the optimum vertex
\draw[engblack, very thick] (3, 1.5) circle (4pt);
\node[engblack, anchor=south west, font=\scriptsize] at (3.05, 1.55)
  {optimum};

\end{tikzpicture}
\caption[Feasible polytope of a two-variable linear
program]{Feasible region of a two-variable linear program with
nonnegativity bounds and two inequality constraints. The feasible
set is a polygon with four vertices. The objective
$\vect{c}^{\transpose} \vect{x}$ is constant on the dashed
parallel lines and increases in the direction of $\vect{c}$. The
optimum lies at a vertex: walking parallel to the level sets in
the direction of $\vect{c}$ and stopping at the last vertex inside
the feasible set picks out $v_{3} = (3, 1.5)$. The simplex method
walks the vertex graph along edges that improve the objective;
interior-point methods approach the optimum through the polytope's
interior. The general $n$-variable case generalises by replacing
the polygon with a polytope and the four vertices with the
exponentially many vertices of the polytope.}
\label{fig:vol02:ch15:feasible-polytope}
\end{figure}


\subsection{Worked LP: the diet problem}

The diet problem (\textcite{text:v2ch15-stigler1945}, posed before
LP existed): choose quantities $x_{i} \geq 0$ of $n$ foods to
minimise total cost $\vect{c}^{\transpose} \vect{x}$ subject to
nutritional requirements $A \vect{x} \geq \vect{b}$, where $A_{ij}$
is the amount of nutrient $i$ in food $j$ and $b_{i}$ is the daily
requirement for nutrient $i$. The standard-form rewrite
introduces slack variables $\vect{s} \geq \vect{0}$ via $A
\vect{x} - \vect{s} = \vect{b}$. The optimum, computed by Dantzig
in 1947 on Stigler's data, gives an annual food cost of $\$39.69$
for adequate nutrition (1939 dollars) on a diet of wheat flour,
cabbage, beef liver, evaporated milk, and dried navy beans, in
specific quantities. The LP solution is mathematically optimal
and gastronomically intolerable; the example is the standard
illustration of the gap between a well-stated objective and the
right objective. A toy five-food, three-nutrient version of the
same LP is implemented as
\repopath{volumes/02-form/ch15-optimisation/code/lp_diet.py};
running it prints the optimal diet, the binding constraints, and
the shadow prices on each nutrient.

\subsection{Worked LP: a two-variable production problem, graphical
and simplex}

A small two-variable LP is solvable by hand in two ways, and
seeing both fixes the geometry that the simplex method automates.
A workshop makes two products. Product 1 yields a profit of
$\$3$ per unit, product 2 a profit of $\$5$ per unit. Three
resources bind: machine A supplies $4$ hours, machine B supplies
$12$ hours, and machine C supplies $18$ hours per shift. Product 1
uses $1$ hour of A per unit; product 2 uses $2$ hours of B and $3$
hours of C per unit. Writing $x_{1}, x_{2} \geq 0$ for the
production quantities, the problem is to maximise $3 x_{1} + 5
x_{2}$ subject to
\[
x_{1} \leq 4, \qquad 2 x_{2} \leq 12, \qquad
3 x_{1} + 2 x_{2} \leq 18, \qquad x_{1}, x_{2} \geq 0.
\]
This is the classic instance from \textcite{text:v2ch15-dantzig1947}'s
successors, sized for hand work.

\paragraph{Graphical solution.}
The feasible set is the polygon bounded by the three resource
lines and the two axes. Its vertices, found by intersecting pairs
of boundary lines, are $(0, 0)$, $(4, 0)$, $(4, 3)$, $(2, 6)$, and
$(0, 6)$. The objective is constant on the parallel lines
$3 x_{1} + 5 x_{2} = \text{const}$; pushing the line outward in
the direction of increasing objective until it last touches the
polygon locates the optimum at a vertex. Evaluating the objective
at each vertex: $(0,0) \mapsto 0$, $(4, 0) \mapsto 12$, $(4, 3)
\mapsto 27$, $(2, 6) \mapsto 36$, $(0, 6) \mapsto 30$. The optimum
is $(x_{1}^{*}, x_{2}^{*}) = (2, 6)$ with profit $p^{*} = 36$. At
this vertex the machine-B constraint ($2 x_{2} = 12$) and the
machine-C constraint ($3 x_{1} + 2 x_{2} = 18$) are both active;
the machine-A constraint ($x_{1} = 2 < 4$) is slack.

\paragraph{Simplex solution.}
The simplex method reaches the same vertex algebraically. Convert
to standard form with slack variables $s_{1}, s_{2}, s_{3} \geq 0$:
\[
x_{1} + s_{1} = 4, \quad 2 x_{2} + s_{2} = 12, \quad
3 x_{1} + 2 x_{2} + s_{3} = 18,
\]
maximise $z = 3 x_{1} + 5 x_{2}$. Start at the origin
$(x_{1}, x_{2}) = (0, 0)$, where the basic variables are the
slacks $(s_{1}, s_{2}, s_{3}) = (4, 12, 18)$ and $z = 0$. The
objective row $z - 3 x_{1} - 5 x_{2} = 0$ has negative reduced
costs on both $x_{1}$ and $x_{2}$; the most negative is $x_{2}$
(coefficient $-5$), so $x_{2}$ enters the basis. The ratio test
$\min(12/2, 18/2) = \min(6, 9) = 6$ selects $s_{2}$ to leave; pivot
on the machine-B row, raising $x_{2}$ to $6$. After the pivot
$z = 30$ and the only remaining negative reduced cost is on $x_{1}$.
The ratio test for $x_{1}$ entering, $\min(4/1, 6/3) = \min(4, 2) =
2$, selects the machine-C slack to leave; pivot, raising $x_{1}$ to
$2$. Now all reduced costs are nonnegative, the optimality test is
met, and the simplex halts at $(x_{1}, x_{2}) = (2, 6)$, $z = 36$.
The two pivots walked the vertex graph $(0,0) \to (0, 6) \to (2,
6)$, each edge strictly increasing the objective, which is the
simplex method in miniature. \Cref{fig:vol02:ch15:lp-graphical}
draws the feasible polygon, the active constraints at the optimum,
and the objective level line.

\begin{figure}[ht]
\centering
\begin{tikzpicture}[x=0.85cm, y=0.85cm, font=\small]

% Constraints: x1 <= 4, 2 x2 <= 12 (x2 <= 6), 3 x1 + 2 x2 <= 18.
% Feasible polygon vertices: (0,0) (4,0) (4,3) (2,6) (0,6).
% Plot region x1 in [0,6], x2 in [0,7].

% Axes
\draw[->, thick] (-0.2, 0) -- (6.6, 0) node[right, font=\scriptsize] {$x_{1}$};
\draw[->, thick] (0, -0.2) -- (0, 7.2) node[above, font=\scriptsize] {$x_{2}$};

% Tick labels
\foreach \t in {2, 4, 6} {
  \draw (\t, 0.08) -- (\t, -0.08) node[below, font=\scriptsize] {\t};
  \draw (0.08, \t) -- (-0.08, \t) node[left, font=\scriptsize] {\t};
}

% Feasible polygon (shaded)
\fill[engblue!12]
  (0,0) -- (4,0) -- (4,3) -- (2,6) -- (0,6) -- cycle;

% Constraint lines
% x1 = 4 (machine A)
\draw[engblack, thin] (4, -0.1) -- (4, 7.0);
\node[engblack, anchor=south, font=\scriptsize, rotate=90] at (4, 6.4)
  {\;$x_{1}=4$ (A)};
% x2 = 6 (machine B)
\draw[engblack, thin] (-0.1, 6) -- (6.4, 6);
\node[engblack, anchor=east, font=\scriptsize] at (6.3, 6.3)
  {$x_{2}=6$ (B)};
% 3 x1 + 2 x2 = 18 (machine C): passes (0,9) and (6,0); in window (2,6)-(6,0)
\draw[engblack, thin] (0, 9) -- (6, 0);
\node[engblack, anchor=west, font=\scriptsize] at (4.6, 2.3)
  {$3x_{1}+2x_{2}=18$ (C)};

% Objective level line through optimum: 3 x1 + 5 x2 = 36 -> x2 = (36 - 3 x1)/5
% at x1=2 -> x2=6; at x1=6 -> x2=3.6; at x1=0 -> 7.2
\draw[engred, thick, dashed] (0, 7.2) -- (6, 3.6);
\node[engred, anchor=west, font=\scriptsize] at (0.2, 6.95)
  {$3x_{1}+5x_{2}=36$};

% Vertices
\foreach \p/\lab/\pos in {%
  {(0,0)}/{}/below left,
  {(4,0)}/{}/below,
  {(4,3)}/{}/right,
  {(0,6)}/{}/left} {
  \fill[engblack] \p circle (1.6pt);
}
% Optimal vertex
\fill[engred] (2,6) circle (2.4pt);
\node[engred, anchor=south west, font=\scriptsize] at (2.1, 6.05)
  {$\vect{x}^{*}=(2,6)$, $z=36$};

\end{tikzpicture}
\caption[Graphical solution of a two-variable LP]{The feasible
polygon of the two-variable production LP, the intersection of the
three resource half-spaces with the first quadrant. The objective
$3x_{1}+5x_{2}$ is constant on the dashed parallel lines; pushing
the line outward until it last touches the polygon locates the
optimum at the vertex $\vect{x}^{*}=(2,6)$, where the machine-B
and machine-C constraints are both active and the machine-A
constraint is slack. The simplex method reaches the same vertex by
walking the edges $(0,0)\to(0,6)\to(2,6)$, each step increasing the
objective. In $n$ dimensions the polygon becomes a polyhedron with
exponentially many vertices, but the optimum is still attained at a
vertex whenever one exists.}
\label{fig:vol02:ch15:lp-graphical}
\end{figure}


\paragraph{The dual and the shadow prices.}
The dual of this maximisation LP assigns a price $y_{1}, y_{2},
y_{3} \geq 0$ to each resource and minimises the total imputed
resource value $4 y_{1} + 12 y_{2} + 18 y_{3}$ subject to the
requirement that no product is underpriced: $y_{1} + 3 y_{3} \geq
3$ (product 1) and $2 y_{2} + 2 y_{3} \geq 5$ (product 2). At the
optimum, complementary slackness fixes the prices. Machine A is
slack, so $y_{1}^{*} = 0$. Products 1 and 2 are both produced, so
their dual constraints are tight: $3 y_{3}^{*} = 3$ gives
$y_{3}^{*} = 1$, and $2 y_{2}^{*} + 2 \cdot 1 = 5$ gives
$y_{2}^{*} = 3/2$. The dual objective $4 \cdot 0 + 12 \cdot 3/2 +
18 \cdot 1 = 18 + 18 = 36$ equals the primal optimum, confirming
strong duality. The shadow prices read directly: one extra hour
of machine B is worth $\$1.50$ of profit, one extra hour of
machine C is worth $\$1.00$, and one extra hour of machine A is
worth nothing because machine A is not the bottleneck. The
engineer should pay up to $\$1.50$ for an additional machine-B
hour and nothing for additional machine-A capacity, a managerial
conclusion that falls out of the dual with no further computation.
The full primal-and-dual solve, including the shadow-price table,
is in
\repopath{volumes/02-form/ch15-optimisation/code/lp_production_simplex.py}.

\subsection{Quadratic programming}

\begin{definition}[Quadratic program]
A \engterm{quadratic program} (QP) is
\[
\min_{\vect{x}} \; \tfrac{1}{2} \vect{x}^{\transpose} P \vect{x}
+ \vect{q}^{\transpose} \vect{x} + r
\quad \text{subject to} \quad
G \vect{x} \leq \vect{h}, \quad A \vect{x} = \vect{b},
\]
with $P \in \R^{n \times n}$ symmetric. The QP is \engterm{convex}
if $P$ is positive semidefinite.
\end{definition}

QPs have a quadratic objective and linear constraints. Convex QPs
inherit the global-optimality guarantee of convex programs.
Engineering instances of QPs:

\begin{itemize}[leftmargin=*]
  \item Least-squares with linear constraints. The objective
    $\tfrac{1}{2} \norm{A \vect{x} - \vect{b}}_{2}^{2}$ expands
    to $\tfrac{1}{2} \vect{x}^{\transpose} (A^{\transpose} A)
    \vect{x} - \vect{b}^{\transpose} A \vect{x} +
    \tfrac{1}{2}\norm{\vect{b}}^{2}$, a convex QP whenever
    $A^{\transpose} A$ is positive semidefinite (always; positive
    definite when $A$ has full column rank).
  \item Portfolio optimisation in the Markowitz form. Maximise
    expected return minus a risk-aversion-weighted variance:
    $\max_{\vect{w}} \boldsymbol{\mu}^{\transpose} \vect{w} -
    \gamma \vect{w}^{\transpose} \Sigma \vect{w}$ subject to
    $\vect{1}^{\transpose} \vect{w} = 1$ and $\vect{w} \geq
    \vect{0}$. The covariance matrix $\Sigma$ is positive
    semidefinite by construction; the QP is convex.
  \item Support vector machine training. The dual SVM problem
    is a QP with a linear equality constraint and box
    constraints on the dual variables.
  \item Model predictive control. At each time step the
    controller solves a QP whose decision variable is the
    sequence of future control inputs over a planning horizon;
    the constraints encode actuator limits and state safety
    bounds.
\end{itemize}

The QP class extends to second-order cone programs (SOCPs) and
semidefinite programs (SDPs); see Boyd and Vandenberghe Chapters
4 and 5.

\subsection{Worked QP: the Markowitz portfolio}

The Markowitz mean-variance portfolio is the QP an engineer is
most likely to meet outside least squares, and it rewards being
worked end to end. An investor allocates a unit of capital across
$n$ assets. The allocation is a weight vector $\vect{w} \in \R^{n}$
with $\vect{1}^{\transpose} \vect{w} = 1$; each $w_{i}$ is the
fraction of capital in asset $i$. The asset returns have mean
vector $\boldsymbol{\mu} \in \R^{n}$ (expected return per unit
capital) and covariance matrix $\Sigma \in \R^{n \times n}$,
positive semidefinite by construction. The portfolio return is
$\boldsymbol{\mu}^{\transpose} \vect{w}$ and the portfolio variance
is $\vect{w}^{\transpose} \Sigma \vect{w}$. The investor wants high
return and low variance, two objectives that pull against each
other. The standard scalarisation trades them with a risk-aversion
parameter $\gamma > 0$:
\[
\min_{\vect{w}} \;\; -\boldsymbol{\mu}^{\transpose} \vect{w}
+ \gamma \, \vect{w}^{\transpose} \Sigma \vect{w}
\quad \text{subject to} \quad
\vect{1}^{\transpose} \vect{w} = 1.
\]
The objective is a convex quadratic (its Hessian is $2 \gamma
\Sigma \succeq 0$), the single constraint is affine, and the
problem is a convex QP. Large $\gamma$ punishes variance and
produces a cautious portfolio; small $\gamma$ chases return and
produces an aggressive one.

\paragraph{Closed form for the equality-only case.}
With only the budget constraint (short sales allowed, so no
$\vect{w} \geq \vect{0}$), the QP has a closed-form solution by the
Lagrange apparatus of Volume II Chapter 11. The Lagrangian is
$\mathcal{L} = -\boldsymbol{\mu}^{\transpose} \vect{w} + \gamma
\vect{w}^{\transpose} \Sigma \vect{w} + \nu (\vect{1}^{\transpose}
\vect{w} - 1)$. Stationarity in $\vect{w}$ gives $-\boldsymbol{\mu}
+ 2 \gamma \Sigma \vect{w} + \nu \vect{1} = \vect{0}$, so
\[
\vect{w}(\nu) = \frac{1}{2 \gamma} \Sigma^{-1}
(\boldsymbol{\mu} - \nu \vect{1}),
\]
and the budget constraint $\vect{1}^{\transpose} \vect{w} = 1$
fixes the multiplier:
\[
\nu^{*} = \frac{\vect{1}^{\transpose} \Sigma^{-1} \boldsymbol{\mu}
- 2 \gamma}{\vect{1}^{\transpose} \Sigma^{-1} \vect{1}}.
\]
The optimal portfolio is the sum of two fixed vectors, one
independent of $\gamma$ and one scaling as $1/\gamma$: a
\engterm{minimum-variance} component $\Sigma^{-1} \vect{1} /
(\vect{1}^{\transpose} \Sigma^{-1} \vect{1})$ that every solution
shares, plus a \engterm{return-tilt} component proportional to
$\Sigma^{-1} \boldsymbol{\mu}$ whose weight grows as the investor
becomes less risk-averse. This two-fund structure is the
mathematical content of the mutual-fund separation theorem: every
mean-variance-optimal portfolio is a blend of the same two funds,
and the investor's risk aversion sets only the blend ratio.

\paragraph{A three-asset instance.}
Take three assets with annual returns $\boldsymbol{\mu} = (0.10,
0.12, 0.15)$ and covariance (in squared-return units)
\[
\Sigma = \begin{pmatrix}
0.040 & 0.006 & 0.000 \\
0.006 & 0.090 & 0.010 \\
0.000 & 0.010 & 0.160
\end{pmatrix},
\]
so asset 1 is low-return and low-variance, asset 3 is high-return
and high-variance, and asset 2 sits between. At a moderate
$\gamma = 5$ the closed form returns weights near $\vect{w}^{*}
\approx (0.60, 0.27, 0.13)$: the cautious investor parks most
capital in the low-variance asset and tilts modestly toward the
two higher-return assets. Lowering $\gamma$ to $1$ shifts weight
toward asset 3 (the return tilt grows by a factor of five),
raising both the expected return and the variance. The locus of
optima as $\gamma$ ranges over $(0, \infty)$ is the
\engterm{efficient frontier} of \cref{fig:vol02:ch15:efficient-frontier},
the upper-left boundary of the achievable risk-return region; every
portfolio strictly inside the region is dominated by a frontier
portfolio that offers more return at the same risk or the same
return at less risk.

\begin{figure}[ht]
\centering
\begin{tikzpicture}[x=1.0cm, y=0.9cm, font=\small]

% Axes: horizontal = risk (std dev), vertical = return.
\draw[->, thick] (-0.2, 0) -- (8.6, 0)
  node[right, font=\scriptsize] {risk $\sqrt{\vect{w}^{\transpose}\Sigma\vect{w}}$};
\draw[->, thick] (0, -0.2) -- (0, 6.4)
  node[above, font=\scriptsize] {return $\boldsymbol{\mu}^{\transpose}\vect{w}$};

% Cloud of random feasible portfolios (interior of the achievable region).
\foreach \x/\y in {%
  3.1/1.4, 3.6/1.1, 4.2/2.0, 4.8/1.6, 5.3/2.4, 5.9/2.0,
  3.4/2.2, 4.0/1.3, 4.6/2.7, 5.1/1.9, 5.7/3.0, 6.3/2.6,
  4.3/3.1, 4.9/2.3, 5.5/3.5, 6.0/3.1, 6.6/3.6, 7.0/3.2,
  3.8/2.6, 5.0/3.4, 5.8/2.7, 6.4/4.0, 7.2/4.1, 6.8/3.0} {
  \fill[enggray!55] (\x, \y) circle (1.0pt);
}

% The efficient frontier: upper-left convex boundary (parabola opening right).
% Parametrise risk r in [2.6, 7.6], return = 1.4 + 1.55*sqrt(r-2.6).
\draw[engblue, very thick, smooth, samples=40, domain=2.6:7.6]
  plot ({\x}, {1.4 + 1.55*sqrt(\x - 2.6)});
\node[engblue, anchor=west, font=\scriptsize] at (5.6, 5.4)
  {efficient frontier};

% The lower (inefficient) boundary for completeness.
\draw[enggray, thin, dashed, smooth, samples=40, domain=2.6:7.6]
  plot ({\x}, {1.4 - 0.55*sqrt(\x - 2.6)});

% Minimum-variance portfolio (leftmost point of the region).
\fill[engred] (2.6, 1.4) circle (2.2pt);
\node[engred, anchor=east, font=\scriptsize] at (2.45, 1.4)
  {min-variance};

% Three marked portfolios along the frontier for gamma = 10, 1, 0.1.
\fill[engorange] (3.6, 2.95) circle (2.0pt);
\node[engorange, anchor=north west, font=\scriptsize] at (3.65, 2.85)
  {$\gamma=10$};
\fill[engorange] (5.0, 3.92) circle (2.0pt);
\node[engorange, anchor=north west, font=\scriptsize] at (5.05, 3.82)
  {$\gamma=1$};
\fill[engorange] (7.0, 4.86) circle (2.0pt);
\node[engorange, anchor=north east, font=\scriptsize] at (6.9, 4.96)
  {$\gamma=0.1$};

\end{tikzpicture}
\caption[The Markowitz efficient frontier]{The achievable
risk-return region for a portfolio of risky assets. Each grey dot is
a feasible portfolio $\vect{w}$ on the simplex; the cloud fills a
region whose upper-left boundary is the efficient frontier (solid).
A portfolio below the frontier is dominated: another portfolio
offers the same return at lower risk, or more return at the same
risk. The minimum-variance portfolio (red) is the leftmost point.
Sweeping the risk-aversion parameter $\gamma$ from large to small
traces the frontier upward and to the right: a risk-averse investor
($\gamma=10$) sits near the minimum-variance point; a
risk-tolerant investor ($\gamma=0.1$) concentrates in
high-return, high-variance assets. Each marked point is the
solution of one convex quadratic program.}
\label{fig:vol02:ch15:efficient-frontier}
\end{figure}


The realistic version adds a no-short-selling constraint
$\vect{w} \geq \vect{0}$, which turns the equality-only QP into an
inequality-constrained QP with no closed form; the solution is the
projection of the unconstrained optimum onto the simplex, computed
by a QP solver in microseconds. The structure survives: the
frontier is still convex, the two-fund intuition still guides the
qualitative reading, and the active set at the optimum tells the
investor which assets the no-short constraint has forced to zero.
The runnable solve over a grid of $\gamma$ values, producing the
frontier and the active-set table, is in
\repopath{volumes/02-form/ch15-optimisation/code/markowitz_portfolio.py}.

\begin{mastery}
\textbf{Interior-point methods and the barrier.} The simplex
method walks the boundary of the feasible polyhedron from vertex
to vertex. Interior-point methods take the opposite route, moving
through the interior toward the optimum, and they are the engine
behind every modern convex solver for LP, QP, SOCP, and SDP. The
construction replaces the inequality constraints $f_{i}(\vect{x})
\leq 0$ with a \engterm{logarithmic barrier} added to the
objective:
\[
\min_{\vect{x}} \; t f_{0}(\vect{x}) - \sum_{i=1}^{m} \log(-f_{i}
(\vect{x})),
\]
where the barrier term $-\log(-f_{i})$ rises to $+\infty$ as
$\vect{x}$ approaches the boundary of the $i$-th constraint,
keeping the iterate strictly feasible. The parameter $t > 0$ sets
the trade-off: for small $t$ the barrier dominates and the
minimiser sits near the analytic centre of the feasible set; as
$t \to \infty$ the original objective dominates and the minimiser
approaches the true optimum. The \engterm{central path} is the
curve of barrier minimisers traced as $t$ increases. The method
solves a sequence of smooth unconstrained problems by Newton's
method, raising $t$ by a constant factor at each outer step, and
reaches an $\epsilon$-accurate optimum in $O(\sqrt{m} \log(1/
\epsilon))$ Newton steps for self-concordant barriers, the result
that made convex optimisation polynomial-time across the LP, QP,
SOCP, and SDP classes (Nesterov and Nemirovski 1994). The
engineer's working knowledge: the barrier converts a constrained
problem into a sequence of unconstrained Newton solves, the
duality gap at each $t$ is $m/t$ (which gives the solver a
principled stopping rule), and the off-the-shelf solver doing
this is the reason a hundred-thousand-variable LP solves in
seconds. \Cref{fig:vol02:ch15:central-path} draws the central path
on a polyhedral feasible region.
\end{mastery}

\begin{figure}[ht]
\centering
\begin{tikzpicture}[x=0.95cm, y=0.95cm, font=\small]

% A polygonal feasible region with the central path curving from the
% analytic centre to the optimal vertex.
% Pentagon-ish feasible polygon.
\coordinate (P1) at (0.4, 0.5);
\coordinate (P2) at (4.8, 0.2);
\coordinate (P3) at (6.2, 2.6);
\coordinate (P4) at (3.4, 5.0);
\coordinate (P5) at (0.0, 3.2);

\fill[engblue!10] (P1) -- (P2) -- (P3) -- (P4) -- (P5) -- cycle;
\draw[engblack, thin] (P1) -- (P2) -- (P3) -- (P4) -- (P5) -- cycle;

% Objective direction (minimise c^T x): arrow pointing toward optimal vertex P2.
\draw[->, thick, enggray] (5.6, 4.6) -- (6.3, 3.9);
\node[enggray, anchor=west, font=\scriptsize] at (5.5, 4.7)
  {$-\vect{c}$ (decrease)};

% Analytic centre (large t small): roughly centroid.
\coordinate (CC) at (2.96, 2.3);
\fill[engred] (CC) circle (2.0pt);
\node[engred, anchor=south, font=\scriptsize] at (2.96, 2.45)
  {analytic centre ($t\!\to\!0$)};

% Optimal vertex.
\fill[engorange] (P3) circle (2.4pt);
\node[engorange, anchor=west, font=\scriptsize] at ($(P3)+(0.12,0.0)$)
  {$\vect{x}^{*}$ ($t\!\to\!\infty$)};

% Central path: smooth curve from analytic centre to optimal vertex.
\draw[engblue, very thick, smooth]
  plot coordinates {(2.96,2.3) (3.6,2.45) (4.4,2.55) (5.2,2.6) (5.9,2.62) (6.2,2.6)};
\node[engblue, anchor=north, font=\scriptsize] at (4.4, 2.35)
  {central path};

% A few intermediate central-path points marked.
\foreach \pt in {(3.6,2.45), (4.4,2.55), (5.2,2.6)} {
  \fill[engblue] \pt circle (1.4pt);
}

\end{tikzpicture}
\caption[The central path of an interior-point method]{The central
path of a barrier interior-point method on a polyhedral feasible
region. For each barrier weight $t$, the minimiser of
$t\,\vect{c}^{\transpose}\vect{x} - \sum_{i}\log(-f_{i}(\vect{x}))$
sits strictly inside the feasible set. At $t\to 0$ the barrier
dominates and the minimiser is the analytic centre (red); as $t$
increases the objective pulls the minimiser along the central path
(blue) toward the optimal vertex $\vect{x}^{*}$ (orange). The solver
takes one or a few Newton steps at each $t$, then raises $t$ by a
constant factor, reaching an $\epsilon$-accurate optimum in
$O(\sqrt{m}\log(1/\epsilon))$ steps. The duality gap at barrier
weight $t$ is exactly $m/t$, which gives the solver a principled
stopping rule and the engineer a certificate of how far the current
iterate sits from optimality.}
\label{fig:vol02:ch15:central-path}
\end{figure}


\subsection{Primal and dual: the LP case}

Every LP has a paired LP, called its dual. The original is the
\engterm{primal}; the paired is the \engterm{dual}. We state the
pairing for the standard form above and develop the general
duality theory in Section 15.4.

\begin{definition}[LP dual, standard form]
The LP $\min_{\vect{x}} \vect{c}^{\transpose} \vect{x}$ subject to
$A \vect{x} = \vect{b}$, $\vect{x} \geq \vect{0}$ has dual
\[
\max_{\boldsymbol{\nu}} \; \vect{b}^{\transpose}
\boldsymbol{\nu} \quad \text{subject to} \quad
A^{\transpose} \boldsymbol{\nu} \leq \vect{c}.
\]
The dual variable $\boldsymbol{\nu} \in \R^{m}$ has one component
per primal equality constraint.
\end{definition}

The dual variable carries a price interpretation. Each $\nu_{i}$
is the marginal value, in objective units, of relaxing the
$i$-th primal constraint by one unit. For the diet problem, the
dual variable $\nu_{i}$ is the implicit price of nutrient $i$ in
the optimal diet: the rate at which the cost would fall if the
$i$-th nutritional requirement were lowered. The dual problem is
the same problem viewed from the constraint side instead of the
decision-variable side.

The dual objective $\vect{b}^{\transpose} \boldsymbol{\nu}$ is a
lower bound on the primal optimal value $\vect{c}^{\transpose}
\vect{x}^{*}$ for any dual-feasible $\boldsymbol{\nu}$. At the
optimum the bound is tight: $\vect{b}^{\transpose}
\boldsymbol{\nu}^{*} = \vect{c}^{\transpose} \vect{x}^{*}$. The
tightness, called \engterm{strong duality}, is the LP version of
the more general theorem in Section 15.4.

\section{Lagrangian duality}\pathtag{standard}

Duality is the discipline of associating to every constrained
problem a paired problem whose optimal value bounds the original.
For a convex problem the bound is tight. For a nonconvex problem
the bound is a valid weak lower bound that an optimiser can use
as a stopping certificate. Duality is also the apparatus that
generates the Karush-Kuhn-Tucker conditions and gives the
Lagrange multipliers their interpretation as constraint shadow
prices.

\subsection{The Lagrangian}

\begin{definition}[Lagrangian]
For the standard-form problem
\[
\min_{\vect{x}} f_{0}(\vect{x}) \;\;\text{s.t.}\;\; f_{i}(\vect{x})
\leq 0, \;\; h_{j}(\vect{x}) = 0,
\]
the \engterm{Lagrangian} is the function
\[
\mathcal{L}(\vect{x}, \boldsymbol{\lambda}, \boldsymbol{\nu}) =
f_{0}(\vect{x})
+ \sum_{i=1}^{m} \lambda_{i} f_{i}(\vect{x})
+ \sum_{j=1}^{p} \nu_{j} h_{j}(\vect{x}),
\]
with $\boldsymbol{\lambda} \in \R^{m}_{\geq 0}$ and
$\boldsymbol{\nu} \in \R^{p}$. The variables $\boldsymbol{\lambda}$
and $\boldsymbol{\nu}$ are the \engterm{dual variables} or
\engterm{Lagrange multipliers}.
\end{definition}

The Lagrangian is the objective plus a weighted sum of the
constraints, with weights chosen freely (subject to nonnegativity
for inequality multipliers). Volume II Chapter 11's Lagrangian was
the equality-only version of this object; the present definition
extends it to inequalities by using nonnegative multipliers,
because relaxing $f_{i}(\vect{x}) \leq 0$ to $f_{i}(\vect{x}) \leq
\epsilon$ for $\epsilon > 0$ enlarges the feasible set, so the
optimum can only decrease, so the multiplier $\lambda_{i} \geq 0$
is the rate of that decrease.

\subsection{The dual function and the dual problem}

\begin{definition}[Dual function, dual problem]
The \engterm{Lagrange dual function} is
\[
g(\boldsymbol{\lambda}, \boldsymbol{\nu}) =
\inf_{\vect{x}} \mathcal{L}(\vect{x}, \boldsymbol{\lambda},
\boldsymbol{\nu}).
\]
The infimum is over all $\vect{x} \in \R^{n}$, with no
constraints. The \engterm{dual problem} is
\[
\max_{\boldsymbol{\lambda} \geq \vect{0}, \boldsymbol{\nu}}
g(\boldsymbol{\lambda}, \boldsymbol{\nu}).
\]
\end{definition}

The dual function is a pointwise infimum of affine functions of
$(\boldsymbol{\lambda}, \boldsymbol{\nu})$ (when the Lagrangian is
viewed for fixed $\vect{x}$ as an affine function of the dual
variables). A pointwise infimum of affine functions is concave.
The dual problem is therefore the maximisation of a concave
function over a convex set, which is itself a convex problem.
This is true regardless of whether the primal is convex.

\subsection{Weak and strong duality}

\begin{theorem}[Weak duality]
For any primal-feasible $\vect{x}$ and any dual-feasible
$(\boldsymbol{\lambda}, \boldsymbol{\nu})$ with
$\boldsymbol{\lambda} \geq \vect{0}$,
\[
g(\boldsymbol{\lambda}, \boldsymbol{\nu}) \leq f_{0}(\vect{x}).
\]
In particular, $d^{*} \leq p^{*}$, where $p^{*}$ is the primal
optimal value and $d^{*}$ is the dual optimal value.
\end{theorem}

The proof is one inequality. For primal-feasible $\vect{x}$ we
have $f_{i}(\vect{x}) \leq 0$ and $h_{j}(\vect{x}) = 0$, so
$\sum_{i} \lambda_{i} f_{i}(\vect{x}) \leq 0$ and $\sum_{j} \nu_{j}
h_{j}(\vect{x}) = 0$, hence $\mathcal{L}(\vect{x},
\boldsymbol{\lambda}, \boldsymbol{\nu}) \leq f_{0}(\vect{x})$.
Taking the infimum over $\vect{x}$ on the left preserves the
inequality.

The difference $p^{*} - d^{*} \geq 0$ is the \engterm{duality gap}.
A zero duality gap means the primal and dual share their optimal
value.

\begin{theorem}[Strong duality, Slater's condition]
If the primal problem is convex and there exists a strictly
feasible point (a point $\vect{x}_{0}$ with $f_{i}(\vect{x}_{0}) <
0$ for all $i$ and $h_{j}(\vect{x}_{0}) = 0$ for all $j$, with the
strict inequality replaced by linearity-and-feasibility for any
affine $f_{i}$), then strong duality holds: $p^{*} = d^{*}$.
\end{theorem}

Slater's condition is the standard regularity condition. For a
convex problem with at least one interior point of the feasible
set, the duality gap closes. For the LP and convex QP cases of
Section 15.3, Slater's condition is automatic if the feasible
set has positive Lebesgue measure, and strong duality is the
working assumption.

For nonconvex problems strong duality typically fails. The
duality gap can be very large, and the dual lower bound can be
loose. The dual is still useful, both as a relaxation that
provides a weak lower bound and as an interpretive object that
gives multipliers their shadow-price reading.

\subsection{KKT conditions, formal version}

Volume II Chapter 11 stated the KKT conditions for a convex
problem with equality and inequality constraints. We restate
them in the duality framework.

\begin{theorem}[Karush-Kuhn-Tucker]
Suppose $f_{0}, f_{i}$ are differentiable and $h_{j}$ is
differentiable. If $\vect{x}^{*}$ and $(\boldsymbol{\lambda}^{*},
\boldsymbol{\nu}^{*})$ are primal- and dual-optimal with zero
duality gap, then they satisfy
\begin{align*}
f_{i}(\vect{x}^{*}) &\leq 0, & i &= 1, \ldots, m
&& \text{(primal feasibility)} \\
h_{j}(\vect{x}^{*}) &= 0, & j &= 1, \ldots, p \\
\lambda_{i}^{*} &\geq 0, & i &= 1, \ldots, m
&& \text{(dual feasibility)} \\
\lambda_{i}^{*} f_{i}(\vect{x}^{*}) &= 0, & i &= 1, \ldots, m
&& \text{(complementary slackness)} \\
\grad f_{0}(\vect{x}^{*}) + \sum_{i} \lambda_{i}^{*} \grad
f_{i}(\vect{x}^{*}) &+ \sum_{j} \nu_{j}^{*} \grad
h_{j}(\vect{x}^{*}) = \vect{0}
&& \text{(stationarity)}
\end{align*}
For a convex problem with Slater's condition, the KKT conditions
are also sufficient: any $(\vect{x}^{*}, \boldsymbol{\lambda}^{*},
\boldsymbol{\nu}^{*})$ satisfying them is primal- and dual-optimal.
\end{theorem}

Complementary slackness is the most useful KKT condition for
hand work. It says that at the optimum, for each inequality
constraint, either the constraint is active ($f_{i}(\vect{x}^{*})
= 0$) or its multiplier vanishes ($\lambda_{i}^{*} = 0$). One of
the two is zero, possibly both, never neither nonzero. The
condition partitions the inequality constraints into the active
set (with positive multipliers) and the inactive set (with zero
multipliers), and the engineer's analysis runs case by case on
this partition.

\subsection{Worked KKT example: a linearly constrained QP}

Minimise $f_{0}(\vect{x}) = \tfrac{1}{2} \norm{\vect{x}}_{2}^{2}$
subject to $\vect{a}^{\transpose} \vect{x} \geq b$, with
$\vect{a} \in \R^{n}$ and $b > 0$. Rewrite the inequality as
$f_{1}(\vect{x}) = b - \vect{a}^{\transpose} \vect{x} \leq 0$. The
Lagrangian is
\[
\mathcal{L} = \tfrac{1}{2} \vect{x}^{\transpose} \vect{x} +
\lambda (b - \vect{a}^{\transpose} \vect{x}).
\]
Stationarity: $\vect{x}^{*} = \lambda^{*} \vect{a}$.
Complementary slackness: $\lambda^{*} (b - \vect{a}^{\transpose}
\vect{x}^{*}) = 0$. Two cases. Case A: $\lambda^{*} = 0$, hence
$\vect{x}^{*} = \vect{0}$. Primal feasibility requires $b - 0
\leq 0$, i.e.\ $b \leq 0$, contrary to assumption. Case B: $b -
\vect{a}^{\transpose} \vect{x}^{*} = 0$. Substituting
$\vect{x}^{*} = \lambda^{*} \vect{a}$ gives $\lambda^{*} = b /
\norm{\vect{a}}_{2}^{2} > 0$. The optimum is $\vect{x}^{*} = b
\vect{a} / \norm{\vect{a}}_{2}^{2}$, the orthogonal projection of
the origin onto the half-space's boundary hyperplane, with
optimal value $p^{*} = b^{2} / (2 \norm{\vect{a}}_{2}^{2})$ and
multiplier $\lambda^{*} = b / \norm{\vect{a}}_{2}^{2}$. The
multiplier reports the rate at which $p^{*}$ decreases per unit
relaxation of $b$: $\dv{p^{*}}{b} = b / \norm{\vect{a}}_{2}^{2} =
\lambda^{*}$, in agreement with the shadow-price reading.
\Cref{fig:vol02:ch15:kkt-geometry} draws this inequality-constraint
projection, with the active and inactive cases side by side;
\cref{fig:vol02:ch15:lagrange-geometry} draws the
equality-constraint version of the same picture.

\begin{figure}[ht]
\centering
\begin{tikzpicture}[x=1cm, y=1cm, font=\small]

% LEFT PANEL: active constraint. Minimise (1/2)||x||^2 subject to
% a^T x >= b with b > 0 so the origin is infeasible; optimum on boundary.
\begin{scope}[shift={(0,0)}]
  \node[anchor=south, font=\small] at (0, 3.2) {constraint active};
  \draw[->, thick] (-3.0, 0) -- (3.0, 0) node[right, font=\scriptsize] {$x_{1}$};
  \draw[->, thick] (0, -2.4) -- (0, 2.9) node[above, font=\scriptsize] {$x_{2}$};
  % Objective level circles
  \foreach \r in {0.5, 1.0, 1.5} { \draw[engblue, thin] (0,0) circle (\r); }
  % Constraint line a^T x = b, a = (1,1), b = 1.4: x_2 = 1.4 - x_1
  % Feasible side a^T x >= b is the upper-right half-plane.
  \draw[engred, very thick] (-1.3, 2.7) -- (2.7, -1.3);
  % Shade feasible half-plane lightly
  \fill[engred!8] (-1.3, 2.7) -- (2.7, -1.3) -- (2.9, 2.7) -- cycle;
  \node[engred, anchor=west, font=\scriptsize] at (1.6, 0.5)
    {$\vect{a}^{\transpose}\vect{x} = b$};
  % Optimum = projection of origin onto line = (b/||a||^2) a = (0.7,0.7)
  \fill[engblack] (0.7, 0.7) circle (2.2pt);
  \node[engblack, anchor=south west, font=\scriptsize] at (0.78, 0.74)
    {$\vect{x}^{*}$};
  % gradient of f points outward along x^*, lambda* grad f1 = -lambda* a inward
  \draw[->, engblue, thick] (0.7, 0.7) -- (1.2, 1.2);
  \node[engblue, anchor=west, font=\scriptsize] at (1.2, 1.25)
    {$\grad f$};
  \fill[engblack] (0,0) circle (1.3pt);
  \node[anchor=north, font=\scriptsize] at (0, -2.05) {$\lambda^{*} > 0$};
\end{scope}

% RIGHT PANEL: inactive constraint. Same objective but b < 0 so the
% unconstrained optimum (origin) is already feasible.
\begin{scope}[shift={(7.5,0)}]
  \node[anchor=south, font=\small] at (0, 3.2) {constraint inactive};
  \draw[->, thick] (-3.0, 0) -- (3.0, 0) node[right, font=\scriptsize] {$x_{1}$};
  \draw[->, thick] (0, -2.4) -- (0, 2.9) node[above, font=\scriptsize] {$x_{2}$};
  \foreach \r in {0.5, 1.0, 1.5} { \draw[engblue, thin] (0,0) circle (\r); }
  % Constraint line pushed away from the origin: a^T x = -1.2,
  % feasible side a^T x >= -1.2 contains the origin.
  \draw[engred, very thick] (-2.8, 1.6) -- (1.6, -2.8);
  \fill[engred!8] (-2.8, 1.6) -- (1.6, -2.8) -- (2.9, 2.9) -- (-2.8, 2.9) -- cycle;
  \node[engred, anchor=west, font=\scriptsize] at (-2.7, 0.9)
    {$\vect{a}^{\transpose}\vect{x} = b$};
  % Optimum = unconstrained minimum at the origin
  \fill[engblack] (0,0) circle (2.2pt);
  \node[engblack, anchor=south west, font=\scriptsize] at (0.08, 0.08)
    {$\vect{x}^{*}=\vect{0}$};
  \node[anchor=north, font=\scriptsize] at (0, -2.05) {$\lambda^{*} = 0$};
\end{scope}

\end{tikzpicture}
\caption[KKT geometry: active and inactive inequality
constraints]{The two cases of complementary slackness for the
projection problem $\min \tfrac{1}{2}\norm{\vect{x}}^{2}$ subject
to $\vect{a}^{\transpose}\vect{x} \geq b$. When $b > 0$ (left) the
unconstrained optimum (the origin) is infeasible, the constraint
binds, the optimum sits on the boundary hyperplane, and the
multiplier $\lambda^{*} = b/\norm{\vect{a}}^{2}$ is strictly
positive; the objective gradient at $\vect{x}^{*}$ is parallel to
the constraint normal. When $b < 0$ (right) the origin is already
feasible, the constraint is slack, the optimum is the
unconstrained minimum, and the multiplier $\lambda^{*} = 0$.
Complementary slackness $\lambda^{*}(b - \vect{a}^{\transpose}
\vect{x}^{*}) = 0$ holds in both cases: one factor is always zero.}
\label{fig:vol02:ch15:kkt-geometry}
\end{figure}


\begin{figure}[ht]
\centering
\begin{tikzpicture}[x=1cm, y=1cm, font=\small]

% Axes
\draw[->, thick] (-3.6, 0) -- (3.6, 0) node[right, font=\scriptsize] {$x_{1}$};
\draw[->, thick] (0, -0.4) -- (0, 4.5) node[above, font=\scriptsize] {$x_{2}$};

% Objective: f(x) = 1/2 || x ||^2, level circles
\foreach \r in {0.6, 1.2, 1.8, 2.4} {
  \draw[engblue, thin] (0,0) circle (\r);
}
\node[engblue, font=\scriptsize] at (1.05, 0.30)
  {$f(\vect{x}) = \tfrac{1}{2}\|\vect{x}\|^{2}$};

% Constraint: a^T x = b, with a = (1, 2), b = 4. Line: x_2 = (4 - x_1)/2.
\draw[engred, very thick] (-3.0, 3.5) -- (3.0, 0.5);
\node[engred, anchor=west, font=\scriptsize] at (3.0, 0.5)
  {$\vect{a}^{\transpose} \vect{x} = b$};

% Optimum: projection of origin onto the line.
%   x^* = (b / ||a||^2) a = (4/5)(1, 2) = (0.8, 1.6)
\fill[engblack] (0.8, 1.6) circle (2.5pt);
\node[engblack, anchor=west, font=\scriptsize] at (0.95, 1.6)
  {$\vect{x}^{*}$};

% Gradient of f at x^*: equals x^*
\draw[->, thick, engblue] (0.8, 1.6) -- (1.6, 3.2);
\node[engblue, anchor=south west, font=\scriptsize] at (1.6, 3.2)
  {$\grad f(\vect{x}^{*}) = \vect{x}^{*}$};

% Gradient of constraint h(x) = a^T x - b: equals a = (1, 2)
\draw[->, thick, engred] (0.8, 1.6) -- (1.8, 3.6);
\node[engred, anchor=east, font=\scriptsize] at (1.7, 3.5)
  {$\nu^{*} \, \nabla h$};

% Annotation: the two gradients are parallel at the optimum,
% with proportionality constant nu^*
\node[anchor=west, align=left, font=\scriptsize] at (-3.4, -0.9)
  {Stationarity at $\vect{x}^{*}$: \\
   $\grad f(\vect{x}^{*}) + \nu^{*}\, \nabla h(\vect{x}^{*}) = \vect{0}$};

% Mark origin
\fill[engblack] (0, 0) circle (1.5pt);

\end{tikzpicture}
\caption[Lagrange-multiplier geometry for an equality
constraint]{Geometry of a Lagrange multiplier for a single linear
equality constraint $\vect{a}^{\transpose} \vect{x} = b$ with
quadratic objective $f(\vect{x}) = \tfrac{1}{2}\|\vect{x}\|^{2}$.
The level sets of $f$ are circles centred on the origin; the
constraint set is the red line. At the optimum $\vect{x}^{*}$, the
constraint line is tangent to a level circle of $f$, which forces
$\grad f$ and $\nabla h$ to be parallel: $\grad f(\vect{x}^{*}) +
\nu^{*} \nabla h(\vect{x}^{*}) = \vect{0}$. The multiplier
$\nu^{*}$ is the proportionality constant; geometrically it is the
rate at which the optimal value changes when the constraint level
$b$ is perturbed by one unit. The inequality case
$f_{i}(\vect{x}) \leq 0$ replaces the tangency picture by
complementary slackness: either the constraint is active (the
optimum lies on the boundary, gradients align) or its multiplier
vanishes (the constraint does not bind).}
\label{fig:vol02:ch15:lagrange-geometry}
\end{figure}


\subsection{Worked example: the dual function of a simple problem}

To see the dual function as an explicit object, take the
one-variable problem $\min_{x} x^{2}$ subject to $x \geq 1$, i.e.\
$f_{1}(x) = 1 - x \leq 0$. The Lagrangian is $\mathcal{L}(x,
\lambda) = x^{2} + \lambda(1 - x)$. The dual function is the
unconstrained infimum over $x$:
\[
g(\lambda) = \inf_{x} \left( x^{2} - \lambda x + \lambda \right).
\]
The inner objective is a convex parabola in $x$; its minimiser is
$x = \lambda/2$, and substituting gives $g(\lambda) = -\lambda^{2}/4
+ \lambda$. This is a concave parabola in $\lambda$, as the general
theory promises. The dual problem maximises $g$ over $\lambda \geq
0$: $g'(\lambda) = -\lambda/2 + 1 = 0$ at $\lambda^{*} = 2$, giving
$d^{*} = g(2) = -1 + 2 = 1$. The primal optimum is $x^{*} = 1$
(the constraint binds) with $p^{*} = 1$. The two values agree,
$p^{*} = d^{*} = 1$, so the duality gap is zero, as Slater's
condition guarantees for this convex problem with a strictly
feasible point ($x = 2$ satisfies $1 - x < 0$). The multiplier
$\lambda^{*} = 2$ is the shadow price: relaxing the constraint to
$x \geq 1 - \epsilon$ lowers the optimum by approximately
$2 \epsilon$, since $\dv{p^{*}}{b}$ at the binding right-hand side
$b = 1$ is $\dv{}{b}(b^{2}) = 2b = 2$.

\subsection{The min-max viewpoint and the saddle point}

Strong duality has a geometric reading as a saddle point of the
Lagrangian. Define the primal as $p^{*} = \min_{\vect{x}}
\max_{\boldsymbol{\lambda} \geq \vect{0}, \boldsymbol{\nu}}
\mathcal{L}(\vect{x}, \boldsymbol{\lambda}, \boldsymbol{\nu})$, the
inner maximisation forcing feasibility (an infeasible $\vect{x}$
sends the inner max to $+\infty$). The dual is $d^{*} =
\max_{\boldsymbol{\lambda} \geq \vect{0}, \boldsymbol{\nu}}
\min_{\vect{x}} \mathcal{L}$. Weak duality is the statement that
swapping a max-min for a min-max can only lower the value:
$\max \min \leq \min \max$, so $d^{*} \leq p^{*}$. Strong duality
is the statement that, for a convex problem with Slater's
condition, the order can be swapped without loss: the Lagrangian
has a saddle point $(\vect{x}^{*}, \boldsymbol{\lambda}^{*},
\boldsymbol{\nu}^{*})$ at which it is simultaneously minimised over
$\vect{x}$ and maximised over the multipliers. The KKT conditions
are exactly the stationarity statement at this saddle point. The
min-max reading recurs throughout engineering: in game theory as
the value of a zero-sum game, in robust optimisation as the
worst-case design, and in the training of generative adversarial
networks as the equilibrium of two competing objectives.

\begin{mastery}
\textbf{Sensitivity analysis: the multiplier as a derivative.} The
shadow-price reading of a Lagrange multiplier can be made exact.
Perturb the right-hand sides of the constraints: replace
$f_{i}(\vect{x}) \leq 0$ with $f_{i}(\vect{x}) \leq u_{i}$ and
$h_{j}(\vect{x}) = 0$ with $h_{j}(\vect{x}) = v_{j}$, and let
$p^{*}(\vect{u}, \vect{v})$ denote the perturbed optimal value.
For a convex problem with strong duality at the unperturbed point,
the optimal multipliers are the gradients of the optimal value
with respect to the perturbations:
\[
\pdv{p^{*}}{u_{i}}\bigg|_{\vect{u} = \vect{0}} = -\lambda_{i}^{*},
\qquad
\pdv{p^{*}}{v_{j}}\bigg|_{\vect{v} = \vect{0}} = -\nu_{j}^{*}.
\]
A large multiplier flags a constraint that is expensive to
tighten and rewarding to relax; a zero multiplier flags a
constraint that does not currently bind and whose small
perturbation costs nothing. This is the quantitative core of every
shadow-price argument: the LP resource prices of Section 15.3, the
diet problem's nutrient prices, and the production LP's machine-
hour prices are all instances of $\partial p^{*} / \partial u$.
The engineer who reports an optimum without its multipliers has
reported half the answer; the multipliers say which constraints to
spend the next dollar relaxing.
\end{mastery}

\begin{mastery}
\textbf{Constraint qualifications: when KKT actually holds.} The
KKT theorem as stated requires zero duality gap, which a convex
problem with Slater's condition supplies. For a nonconvex problem
the KKT conditions are necessary at a local optimum only under a
\engterm{constraint qualification}, a regularity condition on the
constraint geometry at the candidate point. The weakest commonly
checked version is the \engterm{linear independence constraint
qualification} (LICQ): the gradients of the active constraints
(the $\grad f_{i}(\vect{x}^{*})$ for active $i$ together with all
$\grad h_{j}(\vect{x}^{*})$) are linearly independent. Under LICQ
the multipliers exist and are unique. A weaker condition, the
\engterm{Mangasarian-Fromovitz constraint qualification} (MFCQ),
asks only that the equality-constraint gradients be independent and
that a direction exist that strictly decreases every active
inequality while staying tangent to the equalities; MFCQ gives
existence of multipliers but not uniqueness, and it is equivalent
to the boundedness of the multiplier set. The failure mode the
qualifications guard against is concrete. Minimise $x$ subject to
$x^{2} \leq 0$: the only feasible point is $x = 0$, the optimum,
but the active-constraint gradient $\dv{}{x}(x^{2}) = 2x$ vanishes
there, LICQ fails, and the stationarity equation $1 + \lambda \cdot
0 = 0$ has no solution in $\lambda$. No multiplier exists; the KKT
stationarity condition is simply false at the optimum. The
engineering lesson: a solver that returns a point with no valid
multipliers has not necessarily failed, and the first diagnostic is
to check whether a constraint qualification holds, because a
tangent or degenerate constraint can break the optimality
certificate without breaking the optimum.
\end{mastery}

\begin{mastery}
\textbf{Second-order conditions: when KKT is not enough.} For a convex
problem the KKT conditions are sufficient: any point satisfying them is
globally optimal. For a nonconvex problem they are only necessary, and
a point can satisfy every KKT condition while being a saddle or even a
constrained maximum. The discriminator is the second-order condition,
the constrained analogue of the Hessian test. At a KKT point
$(\vect{x}^{*}, \boldsymbol{\lambda}^{*}, \boldsymbol{\nu}^{*})$, form
the Hessian of the Lagrangian in the primal variable, $\nabla^{2}_{
\vect{x}\vect{x}} \mathcal{L} = \nabla^{2} f_{0} + \sum_{i}
\lambda_{i}^{*} \nabla^{2} f_{i} + \sum_{j} \nu_{j}^{*} \nabla^{2}
h_{j}$, and restrict it to the \engterm{critical cone}, the directions
$\vect{d}$ that stay tangent to the active constraints (the $\vect{d}$
with $\grad h_{j}^{\transpose} \vect{d} = 0$ for all $j$ and $\grad
f_{i}^{\transpose} \vect{d} = 0$ for the active inequalities with
positive multiplier). The \engterm{second-order sufficient condition}
for a strict local minimum is that $\vect{d}^{\transpose} \nabla^{2}_{
\vect{x}\vect{x}} \mathcal{L} \, \vect{d} > 0$ for every nonzero
$\vect{d}$ in the critical cone: the Lagrangian curves upward in every
feasible direction. The restriction to the critical cone is the entire
subtlety. The full Lagrangian Hessian can be indefinite, with negative
curvature in directions that point out of the feasible set, and the
point is still a local minimum because those descent directions are
blocked by the constraints. A worked sign of this: minimising
$-x^{2} - y^{2}$ on the circle $x^{2} + y^{2} = 1$ has a Lagrangian
$-x^{2} - y^{2} + \nu(x^{2} + y^{2} - 1)$ whose unconstrained Hessian
$\diag(2\nu - 2, 2\nu - 2)$ is negative definite at the optimal $\nu^{*}
= 1$ would suggest a maximum, yet restricted to the tangent direction
of the circle the second-order term is positive and the KKT point is a
genuine constrained minimum. The engineering lesson: a solver that
returns a KKT point on a nonconvex problem has returned a candidate,
not a verified minimum, and the second-order check on the critical cone
is what upgrades the candidate to a certified local optimum.
\end{mastery}

\begin{mastery}
\textbf{The duality gap and certificates of suboptimality.} Weak
duality gives more than a theorem; it gives the working engineer a
running certificate of how close an iterate sits to the optimum.
At any iteration of a primal-dual method the engineer holds a
primal-feasible $\vect{x}$ and a dual-feasible $(\boldsymbol{\lambda},
\boldsymbol{\nu})$, and the difference $f_{0}(\vect{x}) -
g(\boldsymbol{\lambda}, \boldsymbol{\nu}) \geq 0$ is the
\engterm{surrogate duality gap}. Because $g(\boldsymbol{\lambda},
\boldsymbol{\nu}) \leq p^{*} \leq f_{0}(\vect{x})$, the true
optimum is trapped in the interval $[g, f_{0}]$, and the surrogate
gap is an upper bound on the suboptimality of the current primal
point with no knowledge of $p^{*}$ required. This is the basis of
every principled stopping rule in convex optimisation: iterate
until the surrogate gap falls below a tolerance, and report the
gap alongside the answer so the downstream user knows the certified
accuracy. For an interior-point method the gap is exactly $m/t$ at
barrier weight $t$, so the solver knows its own accuracy at every
step. For a nonconvex problem strong duality fails and the gap does
not close, but the dual value $g$ is still a valid lower bound,
which is exactly what branch-and-bound uses to prune: a subtree
whose dual bound exceeds the incumbent cannot contain a better
solution. The duality gap is the single object that unifies
stopping rules, optimality certificates, and combinatorial pruning.
\end{mastery}

\subsection{Augmented Lagrangian and a worked SQP step}

Two methods turn the duality apparatus into a constrained solver
the engineer actually runs, and both deserve to be seen in motion
rather than only named.

The \engterm{augmented Lagrangian} (method of multipliers) repairs
a defect of the plain quadratic penalty. To solve $\min f_{0}
(\vect{x})$ subject to $h(\vect{x}) = 0$, the penalty method
minimises $f_{0}(\vect{x}) + \tfrac{\rho}{2} h(\vect{x})^{2}$ and
drives $\rho \to \infty$; the trouble is that the Hessian
conditioning degrades as $\rho$ grows, and the unconstrained
subproblems become numerically vicious. The augmented Lagrangian
adds an explicit multiplier estimate to the penalty:
\[
\mathcal{L}_{\rho}(\vect{x}, \nu) = f_{0}(\vect{x}) + \nu \,
h(\vect{x}) + \tfrac{\rho}{2} h(\vect{x})^{2}.
\]
The algorithm alternates a minimisation over $\vect{x}$ at fixed
$(\nu, \rho)$ with a multiplier update $\nu \leftarrow \nu + \rho \,
h(\vect{x})$, which is a gradient ascent step on the dual. The
multiplier update lets the constraint be satisfied at a finite,
moderate $\rho$, sidestepping the ill-conditioning of pure penalty.
A one-step illustration: minimise $f_{0}(x) = \tfrac{1}{2} x^{2}$
subject to $h(x) = x - 1 = 0$, whose answer is $x^{*} = 1$, $\nu^{*}
= -1$. Start with $\nu_{0} = 0$, $\rho = 1$. The $\vect{x}$-step
minimises $\tfrac{1}{2} x^{2} + 0 \cdot (x - 1) + \tfrac{1}{2}(x -
1)^{2}$; setting the derivative to zero, $x + (x - 1) = 0$, gives
$x_{1} = 1/2$. The multiplier update gives $\nu_{1} = 0 + 1 \cdot
(1/2 - 1) = -1/2$. The next $\vect{x}$-step minimises $\tfrac{1}{2}
x^{2} - \tfrac{1}{2}(x - 1) + \tfrac{1}{2}(x - 1)^{2}$; the
derivative $x - 1/2 + (x - 1) = 0$ gives $x_{2} = 3/4$. The
multiplier update gives $\nu_{2} = -1/2 + (3/4 - 1) = -3/4$. The
iterates $x_{k} = 1 - 2^{-k}$ and $\nu_{k} = -(1 - 2^{-k})$ march
geometrically to the optimum at a fixed $\rho = 1$, where a pure
penalty would need $\rho \to \infty$ to reach the same accuracy.
The same construction with vector constraints and inequalities is
the workhorse behind the alternating-direction method of
multipliers (ADMM) that scales convex problems across many machines.

\engterm{Sequential quadratic programming} (SQP) is the constrained
analogue of Newton's method. At the current iterate $(\vect{x}_{k},
\boldsymbol{\nu}_{k})$ it forms a quadratic model of the Lagrangian
and a linear model of the constraints, and solves the resulting
equality-constrained QP for a step. For the problem $\min f_{0}
(\vect{x})$ subject to $h(\vect{x}) = 0$, the QP subproblem is
\[
\min_{\vect{p}} \;\; \grad f_{0}(\vect{x}_{k})^{\transpose} \vect{p}
+ \tfrac{1}{2} \vect{p}^{\transpose} W_{k} \vect{p}
\quad \text{subject to} \quad
\grad h(\vect{x}_{k})^{\transpose} \vect{p} + h(\vect{x}_{k}) = 0,
\]
where $W_{k} = \nabla^{2}_{\vect{x}\vect{x}} \mathcal{L}(\vect{x}_{k},
\boldsymbol{\nu}_{k})$ is the Hessian of the Lagrangian. The KKT
system of this QP is a single linear solve, and its solution
$(\vect{p}, \boldsymbol{\nu}_{k+1})$ supplies both the step and the
new multiplier. A worked step: minimise $f_{0}(x_{1}, x_{2}) =
x_{1}^{2} + x_{2}^{2}$ subject to $h = x_{1} + x_{2} - 2 = 0$, with
optimum $(1, 1)$. From $\vect{x}_{0} = (2, 0)$, the gradient is
$\grad f_{0} = (4, 0)$, the Lagrangian Hessian is $W_{0} = 2 I$ (the
constraint is linear, so it contributes nothing), and the constraint
gradient is $(1, 1)$ with residual $h(\vect{x}_{0}) = 0$. The QP
subproblem $\min_{\vect{p}} (4, 0)^{\transpose} \vect{p} +
\vect{p}^{\transpose} \vect{p}$ subject to $p_{1} + p_{2} = 0$ has
KKT stationarity $2 p_{1} + 4 + \nu = 0$, $2 p_{2} + \nu = 0$, and
$p_{1} + p_{2} = 0$; solving, $p_{1} = -1$, $p_{2} = 1$, $\nu = -2$.
The step takes $\vect{x}_{1} = (2, 0) + (-1, 1) = (1, 1)$, the exact
optimum in one move, because the objective is quadratic and the
constraint linear, exactly the case where SQP inherits Newton's
one-step finish. On a nonlinear constraint SQP converges
quadratically near the optimum, and the working solvers (SNOPT,
filterSQP, IPOPT in its SQP mode) wrap this step in a line search
or trust region and a quasi-Newton approximation to $W_{k}$ so the
engineer never forms the Lagrangian Hessian by hand.

\section{Gradient methods, Newton's method}\pathtag{core}

Outside the closed-form cases (a few small LPs and QPs solvable
by hand) the engineer reaches an optimum by iteration. We treat
the two foundational iterative families: gradient-based methods
that use first derivatives, and Newton's method which uses second
derivatives. Both apply to unconstrained problems on the working
domain. Constrained problems are reduced to unconstrained
problems by penalty methods, barrier methods, or the
projected-gradient construction.

\subsection{Gradient descent}

Volume II Chapter 11 introduced the gradient as the direction of
steepest ascent. Gradient descent moves against it.

\begin{definition}[Gradient descent]
The \engterm{gradient descent} iteration on a differentiable
$f : \R^{n} \to \R$, with initial point $\vect{x}_{0}$ and step
sizes $\eta_{k} > 0$, is
\[
\vect{x}_{k+1} = \vect{x}_{k} - \eta_{k} \grad f(\vect{x}_{k}),
\quad k = 0, 1, 2, \ldots
\]
The step size $\eta_{k}$ is the \engterm{learning rate} when used
in machine learning practice; the step size $\eta_{k}$ is
selected by a fixed schedule, by line search, or by adaptive
rules.
\end{definition}

The iteration's geometry: at $\vect{x}_{k}$, the gradient
$\grad f(\vect{x}_{k})$ points uphill; the negative gradient
points downhill; we step in the downhill direction by an amount
proportional to the gradient's magnitude. As $\vect{x}_{k}$
approaches a critical point the gradient shrinks to zero, the
step shrinks to zero, and the iteration freezes at the critical
point.

\begin{theorem}[Convergence of gradient descent on a smooth
strongly convex objective]
Suppose $f$ is differentiable, $\mu$-strongly convex (the Hessian
is at least $\mu \cdot I$ in semidefinite order), and $L$-smooth
(the gradient is $L$-Lipschitz: $\norm{\grad f(\vect{x}) - \grad
f(\vect{y})} \leq L \norm{\vect{x} - \vect{y}}$). With constant
step size $\eta = 1/L$, gradient descent satisfies
\[
\norm{\vect{x}_{k} - \vect{x}^{*}}_{2}^{2} \leq
(1 - \mu/L)^{k} \norm{\vect{x}_{0} - \vect{x}^{*}}_{2}^{2}.
\]
\end{theorem}

The convergence is geometric, with rate $1 - \mu/L \in [0, 1)$.
The ratio $\kappa = L/\mu$ is the \engterm{condition number} of
the optimisation problem; the iteration count to reach error
$\epsilon$ scales as $\kappa \log(1/\epsilon)$. A
well-conditioned problem (small $\kappa$) converges fast; an
ill-conditioned problem (large $\kappa$) converges slowly. The
canonical bad case is the elongated quadratic
$f(x_{1}, x_{2}) = \tfrac{1}{2}(x_{1}^{2} + 100 x_{2}^{2})$, with
condition number $\kappa = 100$; gradient descent zigzags down
the long axis and converges only after many iterations
(\cref{fig:vol02:ch15:gradient-zigzag}).

\begin{figure}[ht]
\centering
\begin{tikzpicture}[x=1cm, y=1cm, font=\small]

% Axes
\draw[->, thick] (-3.6, 0) -- (3.6, 0) node[right, font=\scriptsize] {$x_{1}$};
\draw[->, thick] (0, -1.5) -- (0, 1.5) node[above, font=\scriptsize] {$x_{2}$};

% Elliptic contours of f(x_1, x_2) = 1/2 (x_1^2 + 100 x_2^2)
% level set f = c: x_1^2/(2c) + x_2^2/(c/50) = 1
% semi-axes (sqrt(2c), sqrt(c/50)). Choose c values that give visible ellipses.
\foreach \c in {0.5, 1.5, 4.5, 9, 13.5} {
  \pgfmathsetmacro{\axA}{sqrt(2*\c)}
  \pgfmathsetmacro{\axB}{sqrt(\c/50)}
  \draw[engblue, thin] (0,0) ellipse ({\axA} and {\axB});
}

% Gradient-descent zigzag trajectory from (3.0, 1.0) with eta = 0.018
% chosen so the iterates oscillate visibly across the long axis.
% Hand-computed iterates of (x_{k+1}) = (1 - eta) x_1, (1 - 100 eta) x_2
% with eta = 0.018: (1 - eta) = 0.982, (1 - 100 eta) = -0.8
\fill[engred] (3.00, 1.00) circle (1.6pt);
\fill[engred] (2.95, -0.80) circle (1.6pt);
\fill[engred] (2.89, 0.64) circle (1.6pt);
\fill[engred] (2.84, -0.51) circle (1.6pt);
\fill[engred] (2.79, 0.41) circle (1.6pt);
\fill[engred] (2.74, -0.33) circle (1.6pt);
\fill[engred] (2.69, 0.26) circle (1.6pt);
\fill[engred] (2.64, -0.21) circle (1.6pt);

\draw[engred, thick]
  (3.00, 1.00) -- (2.95, -0.80) -- (2.89, 0.64) -- (2.84, -0.51)
  -- (2.79, 0.41) -- (2.74, -0.33) -- (2.69, 0.26) -- (2.64, -0.21);

% Annotate starting and current points
\node[engred, anchor=west, font=\scriptsize] at (3.05, 1.10) {$\vect{x}_{0}$};
\node[engred, anchor=west, font=\scriptsize] at (2.65, -0.30) {$\vect{x}_{7}$};

% Annotate the long and short axes
\node[anchor=west, font=\scriptsize] at (3.5, 0.05) {long axis};
\node[anchor=west, font=\scriptsize] at (0.15, 1.2)
  {short axis (steep curvature)};

% Title under the figure body
\node[anchor=north, font=\scriptsize, align=center] at (0, -1.9)
  {$f(x_{1}, x_{2}) = \tfrac{1}{2}(x_{1}^{2} + 100 x_{2}^{2})$,
   $\kappa = 100$};

\end{tikzpicture}
\caption[Gradient-descent zigzag on an elongated quadratic]{Gradient
descent with a fixed step size $\eta$ on the elongated quadratic
$f(x_{1}, x_{2}) = \tfrac{1}{2}(x_{1}^{2} + 100 x_{2}^{2})$, plotted
on its elliptic level sets. The step size large enough to make
useful progress along the long axis ($x_{1}$, eigenvalue $1$)
overshoots the short axis ($x_{2}$, eigenvalue $100$) and the
iteration zigzags. The iteration count to reach a given accuracy
scales with the condition number $\kappa = L/\mu = 100$, the ratio
of the largest to the smallest Hessian eigenvalue. Newton's method
on the same problem reaches the optimum in one step; preconditioning
the gradient by an approximate inverse Hessian achieves the same
effect for general smooth objectives.}
\label{fig:vol02:ch15:gradient-zigzag}
\end{figure}


\begin{estimation}
\textbf{Question.} A medium-scale least-squares problem has $n =
10^{4}$ variables and condition number $\kappa \approx 10^{3}$
(the ratio of the largest to the smallest singular value of the
design matrix). Estimate the iteration count gradient descent
needs to reach relative error $10^{-6}$ from a unit-norm initial
error, and compare against Newton's method on the same problem.

\textbf{Calibrated guess.} The convergence rate is geometric with
factor $1 - 1/\kappa = 0.999$. Each iteration shrinks the error by
$0.1\%$; reaching $10^{-6}$ should take a few times $\kappa
\log(10^{6})$ iterations.

\textbf{Calculation.} The error bound $\norm{\vect{x}_{k} -
\vect{x}^{*}}^{2} \leq (1 - 1/\kappa)^{k}$ reaches $10^{-12}$ in
the norm-squared (so $10^{-6}$ in the norm) when $(1 - 10^{-3})^{k}
= 10^{-12}$, i.e.\ $k = 12 \log 10 / \log(1/(1 - 10^{-3})) \approx
12 \times 2.303 / (10^{-3}) \approx 2.8 \times 10^{4}$ iterations.
Each iteration costs one matrix-vector multiply at
$O(n^{2}) = 10^{8}$ flops, totalling $\approx 3 \times 10^{12}$
flops, roughly an hour on a single 2026-era CPU core at $10^{9}$
flops per second of sustained throughput.

Newton's method takes $\sim 5$ iterations on a quadratic (it
converges in one step exactly for a quadratic, and a few steps for
a near-quadratic), each costing one Hessian solve at $O(n^{3}) =
10^{12}$ flops. Five iterations totals $\approx 5 \times 10^{12}$
flops, comparable to gradient descent's total at this scale but
trading sustained matrix-vector throughput for a single very large
linear solve that must fit in memory. The break-even moves
sharply: at $n = 10^{5}$ the Newton cost rises by $10^{3}$ while
the gradient-descent cost rises by $10^{2}$, and gradient methods
dominate. The principle of method selection (see the box at the
end of this section) is the formal statement of this scaling
analysis.
\end{estimation}

\paragraph{Worked example: gradient descent on a two-variable
quadratic.}
Take $f(x_{1}, x_{2}) = \tfrac{1}{2}(x_{1}^{2} + 9 x_{2}^{2})$,
so $\grad f = (x_{1}, 9 x_{2})$ and the Hessian is $\diag(1, 9)$
with $\mu = 1$, $L = 9$, $\kappa = 9$. The iteration with constant
step $\eta$ is $x_{1}^{(k+1)} = (1 - \eta) x_{1}^{(k)}$ and
$x_{2}^{(k+1)} = (1 - 9 \eta) x_{2}^{(k)}$, decoupled because the
Hessian is diagonal. The $x_{1}$ coordinate contracts by $1 -
\eta$ per step, the $x_{2}$ coordinate by $|1 - 9 \eta|$. With
$\eta = 1/L = 1/9$ the rates are $1 - 1/9 = 8/9 \approx 0.889$ for
$x_{1}$ and $|1 - 1| = 0$ for $x_{2}$; the $x_{2}$ coordinate
collapses in one step and $x_{1}$ crawls. The optimal symmetric
step (Exercise on the gradient-descent rate) is $\eta^{*} = 2/(\mu
+ L) = 2/10 = 0.2$, giving worst-case rate $(\kappa - 1)/(\kappa +
1) = 8/10 = 0.8$ in both coordinates and the fastest plain-gradient
convergence available. From $\vect{x}_{0} = (1, 1)$ at $\eta^{*} =
0.2$: $x_{1}$ follows $1, 0.8, 0.64, 0.512, \ldots$ and $x_{2}$
follows $1, -0.8, 0.64, -0.512, \ldots$, the alternating sign in
$x_{2}$ being the overshoot that produces the zigzag of
\cref{fig:vol02:ch15:gradient-zigzag} at higher $\kappa$. Reaching
$\norm{\vect{x}_{k}} < 10^{-6}$ takes $\log(10^{-6})/\log(0.8)
\approx 62$ iterations, exactly the $\kappa \log(1/\epsilon)$
scaling the convergence theorem predicts. The example is the
whole story of first-order convergence in two coordinates: the
condition number sets the rate, and the best step balances the
two extreme curvatures against each other.

\subsection{Step-size selection}

A constant step size requires knowledge of $L$, which the
engineer typically does not have in closed form. Three working
strategies cover most cases.

\engterm{Backtracking line search}: at each step, start with a
generous step $\eta_{0}$ and shrink $\eta \leftarrow \beta \eta$
(with $\beta \in (0, 1)$) until the Armijo condition $f(\vect{x}_{k}
- \eta \grad f(\vect{x}_{k})) \leq f(\vect{x}_{k}) - \alpha \eta
\norm{\grad f(\vect{x}_{k})}_{2}^{2}$ holds for a fixed $\alpha
\in (0, 1/2)$.

\engterm{Decreasing schedule}: choose $\eta_{k} = \eta_{0} /
(1 + k/\tau)$ for some time scale $\tau$. The schedule is
heuristic but widely used in stochastic settings.

\engterm{Adaptive rules}: AdaGrad, RMSProp, Adam, and their
relatives adjust the per-coordinate step size from running
gradient statistics. The full treatment is in Volume VII
(machine learning); the recognition that they exist and that
they replace a hand-tuned constant step is the engineering
working knowledge for this chapter.

\subsection{Newton's method}

Newton's method uses the second derivative. The motivating
identity: at $\vect{x}_{k}$, the second-order Taylor expansion of
$f$ is
\[
f(\vect{x}_{k} + \vect{p}) \approx f(\vect{x}_{k}) +
\grad f(\vect{x}_{k})^{\transpose} \vect{p} +
\tfrac{1}{2} \vect{p}^{\transpose} H_{k} \vect{p},
\]
where $H_{k} = \nabla^{2} f(\vect{x}_{k})$ is the Hessian.
Minimising the quadratic in $\vect{p}$ gives the Newton step
$\vect{p} = -H_{k}^{-1} \grad f(\vect{x}_{k})$.

\begin{definition}[Newton's method]
The \engterm{Newton iteration} on a twice-differentiable $f$ with
positive-definite Hessian is
\[
\vect{x}_{k+1} = \vect{x}_{k} - H_{k}^{-1} \grad f(\vect{x}_{k}).
\]
\end{definition}

For a strictly convex quadratic $f(\vect{x}) = \tfrac{1}{2}
\vect{x}^{\transpose} P \vect{x} + \vect{q}^{\transpose}
\vect{x}$, Newton's method finds the optimum in one step from
any starting point: the Newton step at $\vect{x}_{0}$ is $-P^{-1}
(P \vect{x}_{0} + \vect{q}) = -\vect{x}_{0} - P^{-1} \vect{q}$,
giving $\vect{x}_{1} = -P^{-1} \vect{q}$, the analytical optimum.
For a non-quadratic objective near a strict local minimum,
Newton's method converges quadratically: the error squares at
each step. The price is a Hessian solve at every iteration,
costing $O(n^{3})$ in dense form, which is prohibitive for
large $n$.

Pure Newton has two failure modes that the working version
repairs. Far from the optimum the full Newton step can overshoot
or even increase the objective, because the quadratic model is a
poor fit over a long step; the repair is \engterm{damped Newton},
which inserts a line-search step length $\eta_{k} \in (0, 1]$ so
that $\vect{x}_{k+1} = \vect{x}_{k} - \eta_{k} H_{k}^{-1} \grad
f(\vect{x}_{k})$ and accepts only steps that satisfy the Armijo
condition, switching to the undamped step ($\eta_{k} = 1$) once
the iterate is close enough for quadratic convergence. The second
failure mode is an indefinite Hessian at a nonconvex point, where
$-H_{k}^{-1} \grad f$ need not be a descent direction at all; the
repair is to modify the Hessian, replacing $H_{k}$ with $H_{k} +
\tau I$ for $\tau > 0$ large enough to make the modified Hessian
positive definite, which interpolates between the Newton step
($\tau = 0$) and a scaled gradient step ($\tau \to \infty$). The
modification is the same idea that appears in Levenberg-Marquardt
and in trust-region subproblems, and it is the reason a robust
Newton implementation does more than invert a Hessian.

\paragraph{Worked example: two Newton steps on a non-quadratic
objective.}
Take $f(x) = x^{4} - 3 x^{2} + x$, with $f'(x) = 4 x^{3} - 6 x +
1$ and $f''(x) = 12 x^{2} - 6$. From $x_{0} = 1.5$: $f'(1.5) =
13.5 - 9 + 1 = 5.5$ and $f''(1.5) = 27 - 6 = 21$, so $x_{1} = 1.5
- 5.5/21 \approx 1.2381$. At $x_{1}$: $f'(1.2381) \approx 4(1.898)
- 7.429 + 1 \approx 1.163$ and $f''(1.2381) \approx 12(1.533) - 6
\approx 12.39$, so $x_{2} \approx 1.2381 - 1.163/12.39 \approx
1.1443$. The true root of $f'$ near here is $x^{*} \approx 1.1309$.
The errors are $|x_{0} - x^{*}| \approx 0.369$, $|x_{1} - x^{*}|
\approx 0.107$, $|x_{2} - x^{*}| \approx 0.0134$. The ratio of
each error to the previous one shrinks ($0.29$ then $0.13$), the
signature of better-than-linear convergence sharpening toward the
quadratic rate as the iterate nears the root, and three or four
steps suffice for machine precision. The same objective started
from $x_{0} = 0$, where $f''(0) = -6 < 0$, would take an ascent
step toward a local maximum, which is exactly the indefinite-
Hessian failure the modification above guards against.

\subsection{Quasi-Newton, recognition level}

The Broyden-Fletcher-Goldfarb-Shanno method (BFGS) and its
limited-memory variant (L-BFGS) maintain a low-rank approximation
to the inverse Hessian, updated from gradient differences across
iterations. The per-iteration cost is $O(n)$ in L-BFGS, against
$O(n^{3})$ for full Newton. The convergence is superlinear,
between gradient descent's geometric rate and Newton's quadratic
rate. L-BFGS is the working off-the-shelf method for medium-scale
unconstrained smooth optimisation; the engineer uses it through
\repopath{scipy.optimize.minimize} with \verb|method='L-BFGS-B'|
or equivalent in another environment, knowing the cost model and
the convergence regime without rederiving the update formula.

\subsection{Line search and trust region}

The two foundational strategies for turning a search direction
into a step are line search and trust region. Both guard against
the failure of a raw Newton or gradient step to decrease the
objective, which happens whenever the local quadratic model is a
poor fit over the length of the step.

\engterm{Line search} fixes the direction $\vect{p}_{k}$ (the
negative gradient, the Newton step, or a quasi-Newton step) and
searches along it for an acceptable step length $\eta_{k} > 0$. A
purely greedy exact line search, $\eta_{k} = \argmin_{\eta} f(
\vect{x}_{k} + \eta \vect{p}_{k})$, is rarely worth its cost; the
working practice accepts any step that decreases the objective by
enough. The \engterm{Armijo} (sufficient-decrease) condition
\[
f(\vect{x}_{k} + \eta \vect{p}_{k}) \leq f(\vect{x}_{k}) + c_{1}
\eta \, \grad f(\vect{x}_{k})^{\transpose} \vect{p}_{k},
\qquad c_{1} \in (0, 1),
\]
demands that the realised decrease be at least a fraction $c_{1}$
of the decrease the linear model predicts. The \engterm{curvature}
(Wolfe) condition adds a lower bound on the step so the iteration
does not stall with tiny steps. Backtracking from a generous
initial step, multiplying by $\beta \in (0, 1)$ until Armijo holds,
is the standard cheap line search and is the rule the gradient-
descent estimation block earlier in this section implicitly
assumed when it set $\eta = 1/L$.

\engterm{Trust region} inverts the order: fix a radius $\Delta_{k}$
within which the local quadratic model
\[
m_{k}(\vect{p}) = f(\vect{x}_{k}) + \grad f(\vect{x}_{k})^{\transpose}
\vect{p} + \tfrac{1}{2} \vect{p}^{\transpose} H_{k} \vect{p}
\]
is trusted, and minimise $m_{k}$ over $\norm{\vect{p}} \leq
\Delta_{k}$. The radius adapts: if the realised decrease matches
the model's predicted decrease (a ratio near $1$), the model is
trustworthy and $\Delta_{k}$ grows; if the realised decrease is
much smaller than predicted, the model is locally wrong and
$\Delta_{k}$ shrinks. Trust-region methods handle indefinite
Hessians gracefully, which is why they are the method of choice
for nonconvex smooth problems where a raw Newton step can point
uphill. The Levenberg-Marquardt method for nonlinear least
squares is a trust-region method in disguise: its damping
parameter is the Lagrange multiplier of the trust-region
constraint, interpolating between a Gauss-Newton step (small
damping, large trust region) and a gradient step (large damping,
small trust region). The engineer meets Levenberg-Marquardt in
every curve-fitting library and now knows what its damping
parameter controls.

\paragraph{Worked example: a trust-region step on an indefinite
model.}
Take a single trust-region step at an iterate where the Hessian is
indefinite, the case a raw Newton step cannot handle. Let the local
model be $m(\vect{p}) = \grad^{\transpose} \vect{p} + \tfrac{1}{2}
\vect{p}^{\transpose} H \vect{p}$ with gradient $\grad = (1, 0)^{
\transpose}$ and Hessian $H = \diag(2, -1)$, indefinite because of the
negative curvature in the second coordinate. The unconstrained
minimiser would solve $H \vect{p} = -\grad$, which fixes the first
coordinate at $-1/2$ but slides to $-\infty$ along the second, since
$m$ has no minimum in the negative-curvature direction. The Newton
step is meaningless here. The trust-region subproblem
$\min_{\norm{\vect{p}} \leq \Delta} m(\vect{p})$ has a finite answer
for any radius. Its solution satisfies $(H + \lambda I) \vect{p} =
-\grad$ for a multiplier $\lambda \geq 0$ chosen so that $H + \lambda
I \succeq 0$ and the step reaches the boundary $\norm{\vect{p}} =
\Delta$. With $\Delta = 1$, the smallest $\lambda$ making $H + \lambda
I$ positive semidefinite is $\lambda = 1$, which lifts the $-1$
eigenvalue to zero. At $\lambda = 1$ the matrix $H + I = \diag(3, 0)$
is singular: the first coordinate is pinned at $p_{1} = -1/3$ by the
stationarity equation, and the second coordinate is free to run along
the zero-curvature direction out to the boundary, $p_{2} = \pm
\sqrt{\Delta^{2} - p_{1}^{2}} = \pm\sqrt{1 - 1/9} \approx \pm 0.943$.
The step rides the negative-curvature direction to the trust-region
boundary, where the model keeps decreasing. The example shows the
mechanism \cref{fig:vol02:ch15:trust-region} draws: when curvature is
negative, the step runs to the boundary along the descent direction the
model offers, the graceful handling of indefinite Hessians that makes
trust-region methods the default for nonconvex smooth problems.

\begin{figure}[ht]
\centering
\begin{tikzpicture}[x=1cm, y=1cm, font=\small]

% Contours of a quadratic model centred off to the upper right,
% with the trust region (a disc of radius Delta) around the current
% iterate at the origin. The constrained minimiser sits on the
% boundary of the disc; the unconstrained model minimiser sits outside.

\def\cx{3.4}
\def\cy{2.3}

% Axes through the current iterate.
\draw[->, thick] (-2.6, 0) -- (5.2, 0) node[right, font=\scriptsize] {$p_{1}$};
\draw[->, thick] (0, -2.4) -- (0, 3.6) node[above, font=\scriptsize] {$p_{2}$};

% Model contours: ellipses centred at the model minimiser (cx, cy).
\foreach \c in {0.5, 1.3, 2.4, 3.8} {
  \pgfmathsetmacro{\axA}{sqrt(2*\c)*1.05}
  \pgfmathsetmacro{\axB}{sqrt(\c)*0.75}
  \draw[engblue, thin, rotate around={20:(\cx,\cy)}]
    (\cx, \cy) ellipse ({\axA} and {\axB});
}

% Unconstrained model minimiser (Newton step), outside the trust region.
\fill[enggray] (\cx, \cy) circle (2pt);
\node[enggray, anchor=west, font=\scriptsize] at (\cx+0.12, \cy)
  {$\vect{p}^{\mathrm{N}}$ (full step)};

% Trust region: disc of radius Delta around the current iterate (origin).
\def\Del{1.9}
\draw[engred, thick, dashed] (0,0) circle (\Del);
\node[engred, anchor=north east, font=\scriptsize] at (-0.05, -0.05)
  {$\norm{\vect{p}} \le \Delta_{k}$};

% Current iterate.
\fill[engblack] (0,0) circle (2pt);
\node[engblack, anchor=north east, font=\scriptsize] at (0.0, 0.15)
  {$\vect{x}_{k}$};

% Constrained step: the model minimiser restricted to the disc lies on
% the boundary, on the dogleg toward p^N. Direction (cx,cy) normalised
% to radius Delta: (3.4,2.3)/|.| * 1.9. |.| = sqrt(3.4^2+2.3^2)=4.105.
% scaled: (1.574, 1.065).
\fill[engred] (1.574, 1.065) circle (2pt);
\node[engred, anchor=south west, font=\scriptsize] at (1.50, 1.12)
  {$\vect{p}_{k}^{*}$};
\draw[->, engred, very thick] (0,0) -- (1.574, 1.065);

\end{tikzpicture}
\caption[The trust-region subproblem]{The trust-region step. At the
current iterate $\vect{x}_{k}$ the algorithm builds a quadratic model
$m_{k}(\vect{p})$ of the objective, whose contours are the ellipses
and whose unconstrained minimiser is the full Newton step
$\vect{p}^{\mathrm{N}}$. The model is trusted only inside the disc
$\norm{\vect{p}} \le \Delta_{k}$. When the full step falls outside the
disc, the constrained minimiser $\vect{p}_{k}^{*}$ lies on the
boundary, in a direction that interpolates between the steepest-descent
direction (for small $\Delta_{k}$) and the Newton direction (for large
$\Delta_{k}$). The radius adapts from the agreement between the model's
predicted decrease and the realised decrease, growing when the model is
trustworthy and shrinking when it is not. The interpolation between
gradient and Newton steps is exactly the mechanism the
Levenberg-Marquardt damping parameter controls.}
\label{fig:vol02:ch15:trust-region}
\end{figure}


\begin{mastery}
\textbf{Conjugate gradient: the optimal first-order method for a
quadratic.} Between gradient descent's geometric crawl and Newton's
expensive solve sits the \engterm{conjugate gradient} (CG) method,
which minimises the quadratic $f(\vect{x}) = \tfrac{1}{2}
\vect{x}^{\transpose} A \vect{x} - \vect{b}^{\transpose} \vect{x}$
(with $A \succ 0$) using only matrix-vector products with $A$ and no
Hessian solve. CG builds its search directions to be mutually
\engterm{$A$-conjugate}, $\vect{p}_{i}^{\transpose} A \vect{p}_{j} = 0$
for $i \neq j$, and this conjugacy makes each one-dimensional line
minimisation permanent: a step along $\vect{p}_{k}$ never spoils the
minimisation already achieved along $\vect{p}_{0}, \ldots,
\vect{p}_{k-1}$. The consequence is exact convergence in at most $n$
steps for an $n$-dimensional quadratic, against gradient descent's
infinite geometric tail. The convergence rate before step $n$ is the
load-bearing improvement: CG reaches error $\epsilon$ in $O(\sqrt{
\kappa} \log(1/\epsilon))$ iterations, the square root of gradient
descent's $O(\kappa \log(1/\epsilon))$, the same
square-root-of-condition-number speedup that momentum and Nesterov
acceleration deliver and that no first-order method can beat for a
worst-case quadratic. CG is the workhorse inside every large sparse
linear solver and inside the truncated-Newton methods that approximate
a Newton step by running a few CG iterations on the Hessian system
without ever factorising it, which is how second-order information is
exploited at the scale where a dense Hessian solve is impossible. The
working knowledge: when the Hessian is available only as a
matrix-vector product (a Hessian-vector product, cheap by automatic
differentiation), CG turns that product into a near-second-order step,
and the $\sqrt{\kappa}$ rate is the reason it earns its keep over plain
gradient descent.
\end{mastery}

\subsection{Nonsmooth objectives: subgradient and proximal methods}

The methods above assume a differentiable objective. Many
engineering objectives are not differentiable everywhere. The
absolute value $\abs{x}$, the $\ell_{1}$ norm $\norm{\vect{x}}_{1}$,
the hinge loss $\max(0, 1 - y \, \vect{w}^{\transpose} \vect{x})$,
and the pointwise maximum $\max_{i} f_{i}(\vect{x})$ all have kinks
where the gradient does not exist. At a kink the gradient is
replaced by the \engterm{subgradient}, the set of all slopes of
supporting lines: a vector $\vect{g}$ is a subgradient of a convex
$f$ at $\vect{x}$ if $f(\vect{y}) \geq f(\vect{x}) +
\vect{g}^{\transpose}(\vect{y} - \vect{x})$ for all $\vect{y}$. At a
smooth point the subgradient is the single gradient; at the kink of
$\abs{x}$ at the origin the subgradient is the whole interval
$[-1, 1]$. Optimality for a nonsmooth convex problem reads
$\vect{0} \in \partial f(\vect{x}^{*})$: zero is a valid
subgradient at the optimum, the nonsmooth generalisation of
$\grad f = \vect{0}$.

\engterm{Subgradient descent} runs the gradient iteration with any
subgradient in place of the gradient, $\vect{x}_{k+1} =
\vect{x}_{k} - \eta_{k} \vect{g}_{k}$ with $\vect{g}_{k} \in
\partial f(\vect{x}_{k})$. It converges, but slowly, at the
$O(1/\sqrt{K})$ rate even for strongly convex problems, because the
subgradient does not vanish near the optimum and a diminishing step
schedule $\eta_{k} \to 0$ is required to settle. The slow rate
makes subgradient descent the method of last resort, used when no
structure beyond convexity is available.

The structure that most nonsmooth engineering objectives do have is
a split into a smooth part and a simple nonsmooth part, $f(\vect{x})
= s(\vect{x}) + r(\vect{x})$ with $s$ smooth (a data-fit term) and
$r$ nonsmooth but simple (a regulariser). The \engterm{proximal
gradient method} exploits the split: take a gradient step on the
smooth part, then apply the \engterm{proximal operator} of the
nonsmooth part,
\[
\vect{x}_{k+1} = \operatorname{prox}_{\eta r}\!\left(\vect{x}_{k} -
\eta \, \grad s(\vect{x}_{k})\right),
\qquad
\operatorname{prox}_{\eta r}(\vect{z}) = \argmin_{\vect{u}}
\left( r(\vect{u}) + \tfrac{1}{2 \eta} \norm{\vect{u} -
\vect{z}}_{2}^{2} \right).
\]
The proximal operator is a small optimisation that pulls toward
$\vect{z}$ while respecting $r$; for the simple regularisers that
arise in practice it has a closed form, and the method recovers the
$O(1/K)$ rate of smooth gradient descent despite the nonsmoothness.

\paragraph{Worked example: LASSO and soft-thresholding.}
The LASSO regression problem is the canonical proximal target:
\[
\min_{\vect{x}} \;\; \tfrac{1}{2} \norm{A \vect{x} - \vect{b}}_{2}^{2}
+ \lambda \norm{\vect{x}}_{1},
\]
a smooth least-squares term plus the nonsmooth $\ell_{1}$
regulariser that drives sparsity. The smooth part has gradient
$A^{\transpose}(A \vect{x} - \vect{b})$; the nonsmooth part is
$r(\vect{x}) = \lambda \norm{\vect{x}}_{1}$, whose proximal
operator separates coordinate by coordinate. For the scalar problem
$\argmin_{u} (\lambda \abs{u} + \tfrac{1}{2 \eta}(u - z)^{2})$, the
optimality condition $0 \in \lambda \, \partial \abs{u} +
(u - z)/\eta$ splits into three cases by the subgradient of
$\abs{u}$, and the solution is the \engterm{soft-thresholding}
operator
\[
\operatorname{prox}_{\eta \lambda \abs{\cdot}}(z) =
\sgn(z) \max(\abs{z} - \eta \lambda, \, 0),
\]
which shrinks $z$ toward zero by $\eta \lambda$ and sets it exactly
to zero when $\abs{z} \leq \eta \lambda$
(\cref{fig:vol02:ch15:soft-threshold}). The dead zone around the
origin is what makes LASSO produce exactly-zero coefficients and
hence sparse, interpretable models, where the smooth $\ell_{2}$
ridge penalty of Section 15.8 only shrinks coefficients toward zero
without ever zeroing them. The full proximal-gradient iteration for
LASSO, a gradient step on the least-squares term followed by
coordinate-wise soft-thresholding, is the \engterm{iterative
shrinkage-thresholding algorithm} (ISTA); its accelerated cousin
FISTA reaches the optimal $O(1/K^{2})$ rate by adding a Nesterov
momentum term. A single worked iteration: with $A = I$, $\vect{b} =
(3, 0.5, -2)$, $\lambda = 1$, $\eta = 1$, from $\vect{x}_{0} =
\vect{0}$ the gradient step gives $\vect{z} = \vect{b} = (3, 0.5,
-2)$, and soft-thresholding by $\eta \lambda = 1$ yields
$\vect{x}_{1} = (2, 0, -1)$: the small coefficient $0.5$ is killed,
the large ones are shrunk by one. The middle coordinate is sparse
after a single step, which is the behaviour the $\ell_{1}$ penalty
was chosen to produce.

\begin{figure}[ht]
\centering
\begin{tikzpicture}[x=0.85cm, y=0.85cm, font=\small]

% Axes: input z (horizontal), output prox (vertical). Threshold tau = 1.5.
\draw[->, thick] (-4.4, 0) -- (4.6, 0) node[right, font=\scriptsize] {$z$};
\draw[->, thick] (0, -3.4) -- (0, 3.6)
  node[above, font=\scriptsize] {$\operatorname{prox}_{\tau\norm{\cdot}_{1}}(z)$};

% Identity line y = z (dashed grey reference).
\draw[enggray, thin, dashed] (-3.0, -3.0) -- (3.0, 3.0);
\node[enggray, anchor=south east, font=\scriptsize] at (3.0, 3.0) {$y=z$};

% Tick marks at +/- tau.
\draw (1.5, 0.08) -- (1.5, -0.08) node[below, font=\scriptsize] {$\tau$};
\draw (-1.5, 0.08) -- (-1.5, -0.08) node[below, font=\scriptsize] {$-\tau$};

% Soft-threshold curve, tau = 1.5:
%   z <= -tau : z + tau   (from (-3,-1.5) to (-1.5,0))
%   |z| <= tau : 0        (from (-1.5,0) to (1.5,0))
%   z >= tau  : z - tau   (from (1.5,0) to (3,1.5))
\draw[engblue, very thick] (-3.0, -1.5) -- (-1.5, 0.0)
  -- (1.5, 0.0) -- (3.0, 1.5);

% Highlight the dead zone.
\fill[engred!12] (-1.5, -0.18) rectangle (1.5, 0.18);
\node[engred, anchor=north, font=\scriptsize] at (0, -0.25)
  {dead zone: output exactly $0$};

\end{tikzpicture}
\caption[The soft-thresholding operator]{The proximal operator of
the scaled $\ell_{1}$ norm, the soft-thresholding map
$\operatorname{prox}_{\tau\norm{\cdot}_{1}}(z) =
\sgn(z)\max(\abs{z}-\tau, 0)$. Inputs of magnitude below the
threshold $\tau$ are set exactly to zero (the red dead zone); inputs
above the threshold are shrunk toward zero by $\tau$. The flat
central segment is the source of sparsity in LASSO and in every
$\ell_{1}$-regularised estimator: the nonsmooth kink of the absolute
value at the origin produces a whole interval of inputs that the
operator collapses to zero, which a smooth penalty such as
$\ell_{2}$ never does. The dashed line $y=z$ is the identity for
reference; the operator tracks it with a constant offset outside the
dead zone.}
\label{fig:vol02:ch15:soft-threshold}
\end{figure}


\begin{principle}[Method selection for smooth unconstrained
optimisation]
For a smooth convex objective in $n$ variables, choose by scale:
$n$ small ($\lesssim 10^{2}$), use Newton's method or interior
point. $n$ medium ($10^{2}$ to $10^{5}$), use L-BFGS or
nonlinear conjugate gradient. $n$ large ($> 10^{5}$, or stochastic
objective), use gradient descent with momentum or Adam. For
nonsmooth objectives, use proximal methods or subgradient
descent. For constrained problems, use projected gradient,
augmented Lagrangian, or interior point.
\end{principle}

\Cref{fig:vol02:ch15:method-comparison} draws the convergence
rates of these methods schematically; the asymptotic separation
between gradient descent's $1 - 1/\kappa$ rate, Nesterov's
$1 - 1/\sqrt{\kappa}$ rate, and Newton's quadratic rate is the
load-bearing reason a working engineer chooses Newton or L-BFGS on
medium-scale smooth problems when the second-order information is
available, and falls back on gradient methods only at the scales
where the Hessian solve becomes infeasible. The standard
references for these rates are
\textcite{text:v2ch15-nocedal-wright2006} for the numerical-analysis
viewpoint and \textcite{text:v2ch15-bottou-curtis-nocedal2018} for
the machine-learning viewpoint.

\begin{figure}[p]
\centering
\begin{tikzpicture}[x=1cm, y=1cm, font=\small]

% Log-scale plot, y axis: log10 of suboptimality
% x axis: iteration count
\draw[->, thick] (0, 0) -- (12.5, 0) node[right, font=\scriptsize] {iteration $k$};
\draw[->, thick] (0, 0) -- (0, 8.5) node[above, font=\scriptsize]
  {$\log_{10}(f - f^{*})$};

% x-axis tick marks (k = 0, 50, 100, 150, 200, 250, 300, 350, 400)
\foreach \xk/\lbl in {0/0, 1.5/50, 3.0/100, 4.5/150, 6.0/200, 7.5/250, 9.0/300, 10.5/350, 12.0/400} {
  \draw (\xk, -0.05) -- (\xk, 0.05);
  \node[font=\tiny, anchor=north] at (\xk, -0.05) {\lbl};
}

% y-axis tick marks (log10 from 0 down to -8)
\foreach \yk/\lbl in {0/0, 1/-1, 2/-2, 3/-3, 4/-4, 5/-5, 6/-6, 7/-7, 8/-8} {
  \draw (-0.05, \yk) -- (0.05, \yk);
  \node[font=\tiny, anchor=east] at (-0.1, \yk) {\lbl};
}

% Gradient descent on a strongly convex quadratic with kappa = 100.
% rate = 1 - mu/L = 0.99 per step, so log10 of error decreases
% by log10(0.99) = -0.00436 per step. Reaches 1e-8 at about 1830 steps.
% Plot from k = 0..400, error goes from 1 to about 0.0178 (log10 = -1.75)
\draw[engblue, thick]
  (0, 0)
  -- (1.5, {0.00436*50})
  -- (3.0, {0.00436*100})
  -- (4.5, {0.00436*150})
  -- (6.0, {0.00436*200})
  -- (7.5, {0.00436*250})
  -- (9.0, {0.00436*300})
  -- (10.5, {0.00436*350})
  -- (12.0, {0.00436*400});
\node[engblue, anchor=west, font=\scriptsize] at (12.0, {0.00436*400 - 0.1})
  {gradient descent ($\kappa = 100$)};

% Accelerated gradient on the same problem.
% Rate ~ 1 - sqrt(mu/L) = 1 - 0.1 = 0.9 per step, log10(0.9) = -0.0458
% so reaches 1e-8 at about 175 steps.
\draw[engred, thick]
  (0, 0)
  -- (1.5, {0.0458*50})
  -- (3.0, {0.0458*100})
  -- (4.5, {0.0458*150})
  -- (5.4, 8.0);
\node[engred, anchor=west, font=\scriptsize] at (5.4, 8.1)
  {Nesterov accelerated};

% Newton's method: quadratic convergence, error squares.
% Starting from 1, error after k steps is 1 / (2^(2^k)). Reaches 1e-8
% by k = 3. Plot as a steep curve to the floor.
\draw[engblack, thick]
  (0, 0)
  -- (0.45, 1)
  -- (0.9, 4)
  -- (1.35, 8.0);
\fill[engblack] (0, 0) circle (1.6pt);
\fill[engblack] (0.45, 1) circle (1.6pt);
\fill[engblack] (0.9, 4) circle (1.6pt);
\fill[engblack] (1.35, 8.0) circle (1.6pt);
\node[engblack, anchor=west, font=\scriptsize] at (1.4, 8.2)
  {Newton};

% L-BFGS, superlinear (between geometric and quadratic). Plot in between.
\draw[gray, thick]
  (0, 0)
  -- (1.5, 0.7)
  -- (3.0, 2.0)
  -- (4.5, 4.0)
  -- (6.0, 6.5)
  -- (7.5, 8.0);
\node[gray, anchor=west, font=\scriptsize] at (7.5, 8.2)
  {L-BFGS};

% SGD on the noisy version, O(1/sqrt(k)) rate; log10 error = -0.5 log10(k+1) + const
% At k=50: -0.5 log10(51) = -0.85. At k=400: -0.5 log10(401) = -1.30.
% Add noise (jitter).
\draw[engblue, thick, dash pattern=on 2pt off 2pt]
  (0, 0)
  -- (0.5, 0.10) -- (1.0, 0.40) -- (1.5, 0.85)
  -- (2.0, 0.95) -- (3.0, 1.10) -- (4.0, 1.18)
  -- (5.0, 1.25) -- (6.0, 1.27) -- (7.5, 1.30)
  -- (9.0, 1.32) -- (10.5, 1.33) -- (12.0, 1.30);
\node[engblue, anchor=west, font=\scriptsize] at (12.0, 1.20)
  {SGD ($\sigma^{2} = 1$)};

% Legend / annotation block under the body
\node[anchor=north west, align=left, font=\scriptsize] at (0, -1.0) {
  Convergence rates on a smooth strongly convex objective with $\kappa = 100$.\\
  Newton: quadratic (error squares each step). \\
  Nesterov accelerated gradient: geometric with rate $1 - \sqrt{\mu/L}$. \\
  L-BFGS: superlinear in practice on smooth medium-scale problems. \\
  Gradient descent: geometric with rate $1 - \mu/L = 1 - 1/\kappa$. \\
  SGD on the noisy version: $O(1/\sqrt{k})$ to a noise floor. \\
  The right method depends on $n$, noise, smoothness, and per-iteration cost.
};

\end{tikzpicture}
\caption[Convergence rates of optimisation methods]{Schematic
convergence rates of the chapter's algorithms on a smooth strongly
convex problem with condition number $\kappa = 100$. Newton's
method squares the error each step and reaches double precision in
under five iterations, at the cost of an $O(n^{3})$ Hessian solve
per step. Nesterov's accelerated gradient achieves a geometric
rate of $1 - \sqrt{\mu/L} = 1 - 1/\sqrt{\kappa}$, a $\sqrt{\kappa}$
improvement over plain gradient descent's $1 - 1/\kappa$. L-BFGS
sits between, with superlinear convergence on smooth medium-scale
problems at $O(n)$ per step. Stochastic gradient descent on the
noisy version of the same problem converges at $O(1/\sqrt{k})$ to
a noise floor set by the per-sample gradient variance; the $1/k$
rate of the strongly convex case requires variance reduction or a
schedule $\eta_{k} \propto 1/k$. The choice of method is governed
by the dimension $n$, the access pattern to the gradient (full
batch versus stochastic), the smoothness regime, and the
per-iteration cost. The principle of Section 15.5 codifies the
selection rule.}
\label{fig:vol02:ch15:method-comparison}
\end{figure}


\section{Stochastic optimisation, SGD}\pathtag{standard}

The objective in many engineering settings is a sum or expectation
over data:
\[
f(\boldsymbol{\theta}) = \frac{1}{N} \sum_{i=1}^{N}
\ell(\boldsymbol{\theta}, \vect{z}_{i}),
\]
with $N$ data points $\vect{z}_{i}$, parameter
$\boldsymbol{\theta} \in \R^{p}$, and per-sample loss $\ell$. A
linear regression's least-squares objective, a logistic regression's
log-loss, and a neural network's training loss all fit this form.
For $N$ large (a million data points, or a hundred million), each
gradient evaluation $\grad f$ costs $N$ per-sample gradient
computations, which is prohibitive for an iterative method that
needs many gradient evaluations.

\subsection{Stochastic gradient descent}

\begin{definition}[Stochastic gradient descent]
\engterm{Stochastic gradient descent} (SGD) replaces the full
gradient $\grad f(\boldsymbol{\theta}_{k}) = (1/N) \sum_{i} \grad
\ell(\boldsymbol{\theta}_{k}, \vect{z}_{i})$ with an unbiased
estimate built from a randomly sampled subset (a \engterm{minibatch}
$B_{k} \subseteq \set{1, \ldots, N}$ of size $|B_{k}| = b$):
\[
\boldsymbol{\theta}_{k+1} = \boldsymbol{\theta}_{k} - \eta_{k}
\cdot \frac{1}{b} \sum_{i \in B_{k}} \grad
\ell(\boldsymbol{\theta}_{k}, \vect{z}_{i}).
\]
The minibatch is resampled at every iteration.
\end{definition}

The per-iteration cost falls from $N$ gradient evaluations to $b$,
a factor-$N/b$ saving. The price is gradient noise: the
minibatch gradient is $\E[\grad f]$ in expectation but has variance
$\sigma^{2}/b$ where $\sigma^{2}$ is the per-sample gradient
variance. The variance shrinks as the minibatch grows, but the
per-iteration cost grows in proportion; the engineer's choice of
$b$ balances per-iteration cost against per-iteration progress.

\subsection{Convergence}

\begin{theorem}[SGD convergence on a smooth convex objective]
Suppose $f$ is convex and $L$-smooth, and the per-sample gradient
has variance bounded by $\sigma^{2}$. With step sizes $\eta_{k} =
\eta_{0} / \sqrt{k+1}$, SGD achieves
\[
\E[f(\bar{\boldsymbol{\theta}}_{K})] - f(\boldsymbol{\theta}^{*})
\leq O(1/\sqrt{K}),
\]
where $\bar{\boldsymbol{\theta}}_{K}$ is the running average of
the iterates. For $\mu$-strongly convex $f$, the rate improves to
$O(1/K)$.
\end{theorem}

The rate $O(1/\sqrt{K})$ is slower than gradient descent's
geometric rate; SGD trades per-iteration progress for
per-iteration cost. For modern problem sizes (millions of data
points, billions of parameters) the trade is overwhelmingly in
SGD's favour: a single full-gradient evaluation costs more than
thousands of SGD iterations, and the SGD iterations make
non-trivial progress at every step.

\paragraph{Worked example: the bias-variance budget of a
minibatch.}
A minibatch gradient is an unbiased estimate of the full gradient,
so the question is entirely one of variance. Suppose the per-sample
gradient at the current parameter has component-wise variance
$\sigma^{2} = 4$ (in squared gradient units) across the $N = 10^{6}$
training samples. A minibatch of size $b$ drawn without replacement
has gradient variance $(\sigma^{2}/b)(1 - b/N) \approx \sigma^{2}/b$
for $b \ll N$. At $b = 1$ the variance is $4$; at $b = 256$ it is
$4/256 \approx 0.0156$, a factor-$256$ reduction; at $b = 10^{4}$
it is $4 \times 10^{-4}$. The variance falls as $1/b$ while the
per-iteration cost rises as $b$, so the variance reduction per unit
of compute is constant: doubling the batch halves the variance and
doubles the cost, a wash on the variance-per-flop ledger. The
reason to use a moderate batch ($b = 32$ to $b = 512$ in 2026
practice) is not the variance-per-flop ratio, which is flat, but
hardware throughput: a single matrix multiply over a batch
saturates a modern accelerator's arithmetic units far better than
$b$ separate vector operations, so the wall-clock cost of a
size-$b$ batch is sublinear in $b$ over the relevant range. The
batch size is a hardware-utilisation decision wearing the costume
of a statistics decision, and the engineer who tunes it should
profile the throughput curve, not just the convergence curve.

\subsection{Variance reduction, recognition level}

The $1/\sqrt{K}$ rate of SGD has been improved to $1/K$ (matching
full gradient descent on strongly convex problems) by methods
that reduce the gradient variance: SAG (stochastic average
gradient), SAGA, SVRG (stochastic variance-reduced gradient), and
control-variate constructions. These methods maintain auxiliary
state (cached per-sample gradients or full-gradient checkpoints)
to subtract a learned mean from the minibatch gradient. The cost
of the state is $O(N)$ memory. The recognition for an engineer:
when SGD has converged but the engineer wants the geometric rate
of full gradient descent without the full-gradient cost, the
variance-reduction family is the working tool. The full
treatment is in modern optimisation references (Bubeck 2015,
Bottou-Curtis-Nocedal 2018); we leave this at recognition level.

\subsection{Momentum and Adam, recognition level}

Three modifications to SGD recur in practice:

\engterm{Momentum} (Polyak): maintain a velocity $\vect{v}_{k}$
that accumulates past gradients with a decay $\beta \in (0, 1)$.
The update becomes $\vect{v}_{k+1} = \beta \vect{v}_{k} +
\hat{\grad} f$, $\boldsymbol{\theta}_{k+1} =
\boldsymbol{\theta}_{k} - \eta \vect{v}_{k+1}$. Momentum smooths
gradient noise and accelerates convergence in elongated valleys
where vanilla SGD zigzags.

\engterm{Nesterov accelerated gradient}: a refinement of momentum
that evaluates the gradient at the look-ahead point
$\boldsymbol{\theta}_{k} + \beta \vect{v}_{k}$ rather than at
$\boldsymbol{\theta}_{k}$. Nesterov's method achieves the optimal
$O(1/K^{2})$ rate for smooth convex problems, an improvement over
gradient descent's $O(1/K)$ rate.

\engterm{Adam} (\textcite{text:v2ch15-kingma-ba2015-adam}): combines
momentum with per-coordinate adaptive step sizes built from running
second-moment estimates of the gradient. Adam is the default first
attempt for training a modern neural network. Its convergence
theory is weaker than SGD's (some convex problems on which Adam
fails to converge are known); its practical performance on deep
nonconvex objectives is the empirical reason for its dominance,
current as of 2026.

\begin{mastery}
\textbf{Momentum as a damped dynamical system.} Polyak's momentum
has a mechanical reading that explains why it accelerates. Write the
momentum update as a single second-order recurrence in the
parameter: $\boldsymbol{\theta}_{k+1} = \boldsymbol{\theta}_{k} -
\eta \grad f(\boldsymbol{\theta}_{k}) + \beta(\boldsymbol{\theta}_{k}
- \boldsymbol{\theta}_{k-1})$. The term $\beta(\boldsymbol{\theta}_{k}
- \boldsymbol{\theta}_{k-1})$ carries the previous step forward, so
the iterate behaves like a particle with mass rolling on the loss
surface under a force $-\grad f$ and a friction set by $\beta$. The
continuous limit is the equation of a damped harmonic oscillator,
$\ddot{\boldsymbol{\theta}} + a \dot{\boldsymbol{\theta}} + \grad
f(\boldsymbol{\theta}) = 0$, with the friction coefficient $a$ tied
to $1 - \beta$. On the elongated quadratic that makes plain
gradient descent zigzag, the heavy-ball particle builds speed along
the shallow valley floor and the friction damps the oscillation
across the steep walls, which is exactly the acceleration the
method delivers. Choosing $\beta$ is choosing the damping: too
little friction (small $\beta$) and the particle barely accelerates;
too much (large $\beta$) and it overshoots and rings. The critical
damping that reaches the minimum fastest corresponds to $\beta =
((\sqrt{\kappa} - 1)/(\sqrt{\kappa} + 1))^{2}$, and it converts
gradient descent's $1 - 1/\kappa$ rate into the $1 - 1/\sqrt{\kappa}$
rate that Nesterov's method attains, a square-root improvement in
the condition-number dependence that is the whole reason momentum
exists. The dynamical reading also explains why momentum tolerates
the gradient noise of SGD: the velocity term is a running average
of past gradients, and averaging reduces variance, so momentum
smooths the stochastic gradient as a side effect of its mechanics.
\end{mastery}

\begin{mastery}
\textbf{Online learning and the regret bound.} The convergence theorems
above assume a fixed objective. Many engineering systems instead face a
sequence of objectives revealed one at a time: a pricing system sees a
new demand curve each day, a control loop sees a new disturbance each
step, an ad system sees a new query each request. The framework is
\engterm{online convex optimisation}: at each round $t$ the engineer
commits a decision $\vect{x}_{t}$, then an adversary reveals a convex
loss $f_{t}$, and the loss $f_{t}(\vect{x}_{t})$ is incurred. The
performance measure is not distance to an optimum (there is no fixed
optimum) but \engterm{regret}, the gap between the cumulative loss
incurred and the cumulative loss of the best single fixed decision in
hindsight:
\[
\mathrm{Regret}_{T} = \sum_{t=1}^{T} f_{t}(\vect{x}_{t}) -
\min_{\vect{x}} \sum_{t=1}^{T} f_{t}(\vect{x}).
\]
Online gradient descent, $\vect{x}_{t+1} = \vect{x}_{t} - \eta_{t}
\grad f_{t}(\vect{x}_{t})$, achieves $\mathrm{Regret}_{T} =
O(\sqrt{T})$ for convex losses with a diminishing step
$\eta_{t} \propto 1/\sqrt{t}$, and $O(\log T)$ for strongly convex
losses. The $\sqrt{T}$ bound means the per-round regret
$\mathrm{Regret}_{T}/T \to 0$: the online learner's average
performance approaches that of the best fixed decision, even against an
adversary who chooses the losses knowing the algorithm. The connection
to stochastic optimisation is exact: running online gradient descent on
losses $f_{t}$ drawn i.i.d.\ from a distribution, and averaging the
iterates, recovers the SGD guarantee of this section, so the regret
bound is the adversarial generalisation of the SGD convergence rate.
The engineer's working takeaway: when the objective drifts or arrives
in a stream, the right guarantee is a regret bound, not a convergence
rate, and online gradient descent is the first method to reach for.
\end{mastery}

\begin{mastery}
\textbf{Derivative-free optimisation.} Every method so far has assumed
a computable gradient. Some engineering objectives offer none: the
output of a legacy simulator with no adjoint, a physical experiment
where each evaluation is a wind-tunnel run, a black-box hyperparameter
tuning where the objective is a full training run. \engterm{Derivative-
free optimisation} (DFO) works from function values alone, and the
engineer should know the three families and their cost. Pattern-search
and Nelder-Mead simplex methods maintain a geometric pattern of trial
points and move it downhill by reflecting the worst vertex; they are
simple and tolerant of noise on low-dimensional smooth problems but scale poorly
beyond a few dozen variables. Model-based DFO (the trust-region method
of the NEWUOA and BOBYQA family) fits a quadratic surrogate to sampled
function values and minimises it inside a trust region, recovering much
of Newton's efficiency without a gradient, at the cost of $O(n^{2})$
samples to fit the surrogate. Bayesian optimisation fits a Gaussian-
process surrogate with calibrated uncertainty and chooses the next
evaluation to balance exploiting the surrogate's predicted optimum
against exploring its uncertain regions, which makes it the method of
choice when each evaluation is expensive (minutes to days) and the
dimension is modest. The unifying cost model: DFO spends function
evaluations to reconstruct the gradient information that gradient-based
methods get for free, so its iteration count scales worse with
dimension, and the engineer reaches for it only when a gradient is
genuinely unavailable. The first question on any black-box problem is
whether automatic differentiation or an adjoint can supply a gradient
after all, because a gradient-based method will almost always beat DFO
when a gradient exists.
\end{mastery}

The engineer's working catalogue: vanilla SGD is the baseline;
SGD with momentum is the next step up; Adam is the modern
default for deep nonconvex problems; L-BFGS is the smooth
medium-scale workhorse. The reader implements the first three in
the chapter project (Section 15.9). A reference implementation of
vanilla gradient descent on the elongated quadratic of this
section is at
\repopath{volumes/02-form/ch15-optimisation/code/gradient_descent.py},
and a comparison of gradient descent, momentum, and Adam on the
Rosenbrock function (the canonical nonconvex benchmark) is at
\repopath{volumes/02-form/ch15-optimisation/code/rosenbrock_methods.py}.
Reference convergence counts from these scripts are tabulated in
\repopath{volumes/02-form/ch15-optimisation/data/rosenbrock_convergence.csv}.

\section{Integer and combinatorial optimisation preview}\pathtag{standard}

Some decision variables take values only in a discrete set:
binary on-off switches, integer counts, categorical assignments,
graph structures. Optimisation over such variables is integer
programming or combinatorial optimisation, the family of problems
in which convexity has, structurally, no direct analogue.

\subsection{Integer linear programming}

\begin{definition}[Integer linear program]
An \engterm{integer linear program} (ILP) is an LP with the
additional constraint that each component of $\vect{x}$ is
integer-valued: $\vect{x} \in \Z^{n}$. A \engterm{mixed-integer
linear program} (MILP) requires only a subset of components to be
integer.
\end{definition}

ILP captures a wide range of engineering problems: facility
location (build the warehouse or do not), scheduling (assign each
job to a machine), routing (a vehicle either traverses an edge or
does not), production planning with set-up costs (produce a
positive batch with a fixed cost or produce zero), graph
colouring, the travelling salesman problem.

\subsection{Why integer programming is hard}

The integer constraint destroys convexity: the integer lattice is
not convex, so the LP relaxation's vertex enumeration argument
does not apply. The decision space grows exponentially in the
number of integer variables: a problem with $n$ binary variables
has $2^{n}$ feasible candidates in the worst case.

\begin{theorem}[Cook 1971; Karp 1972]
The integer linear programming feasibility problem is NP-complete.
Many of its specialisations (the travelling salesman problem,
graph colouring, set cover, knapsack with general profits) are
NP-hard.
\end{theorem}

NP-hardness, treated formally in Volume II Chapter 17, is the
working signal that the engineer cannot expect a polynomial-time
algorithm for the worst case of an integer problem. The engineer
nonetheless solves integer problems daily; the working strategies
are branch-and-bound, cutting planes, and approximation algorithms
that trade an optimality gap for tractable runtime.

\subsection{Worked branch-and-bound: a small knapsack}

A four-item 0/1 knapsack makes the branch-and-bound mechanics
concrete. Items have values $\vect{v} = (60, 100, 120, 40)$ and
weights $\vect{w} = (10, 20, 30, 10)$; the capacity is $W = 50$.
Maximise $\sum_{i} v_{i} x_{i}$ over $x_{i} \in \set{0, 1}$ subject
to $\sum_{i} w_{i} x_{i} \leq 50$. The value-to-weight ratios are
$6, 5, 4, 4$, so the LP relaxation (fractional knapsack) fills
greedily by ratio: take item 1 ($w = 10$), item 2 ($w = 20$,
running weight $30$), then $20/30$ of item 3 to fill the remaining
$20$ of capacity, for a relaxed value $60 + 100 + (2/3)(120) = 240$.
The relaxation gives the root upper bound $240$ and a fractional
variable $x_{3} = 2/3$ to branch on.

\emph{Branch $x_{3} = 0$.} Capacity $50$ over items $\set{1, 2, 4}$.
Greedy by ratio takes items 1, 2, 4 (total weight $40 \leq 50$),
all integer, value $60 + 100 + 40 = 200$. This is an integer-
feasible solution; the incumbent becomes $200$.

\emph{Branch $x_{3} = 1$.} Commit item 3 ($w = 30$, $v = 120$),
remaining capacity $20$ over items $\set{1, 2, 4}$. The LP
relaxation takes item 1 ($w = 10$), then $10/20$ of item 2, for a
relaxed value $120 + 60 + (1/2)(100) = 230$. The bound $230$
exceeds the incumbent $200$, so this subtree may contain a better
solution; branch on the fractional $x_{2}$.

\emph{Subbranch $x_{3} = 1, x_{2} = 0$.} Remaining capacity $20$
over items $\set{1, 4}$: take both ($w = 20$), value $120 + 60 +
40 = 220$, integer-feasible. The incumbent improves to $220$.

\emph{Subbranch $x_{3} = 1, x_{2} = 1$.} Committed weight $30 + 20
= 50 = W$ exactly, capacity exhausted, value $120 + 100 = 220$.
Integer-feasible, value $220$, tying the incumbent.

All open subtrees now have LP bounds at or below the incumbent
$220$ and are pruned. The optimum is $220$, achieved by
$\set{1, 3, 4}$ or $\set{2, 3}$. The tree explored a handful of nodes
instead of the $2^{4} = 16$ leaves of brute force; the LP
relaxation supplied the bound that pruned the rest. The pattern
scales: the tighter the relaxation, the more aggressively the tree
prunes, which is the practical lesson restated in the
branch-and-bound estimate of the estimation exercise set. A
runnable implementation that prints the visited-node count against
the brute-force leaf count is in
\repopath{volumes/02-form/ch15-optimisation/code/knapsack_branch_and_bound.py}.

\subsection{Worked assignment problem: integrality for free}

Not every integer program is hard. The \engterm{assignment
problem} asks for a one-to-one matching of $n$ agents to $n$ tasks
that minimises total cost: choose $x_{ij} \in \set{0, 1}$, with
$x_{ij} = 1$ when agent $i$ is assigned task $j$, to minimise
$\sum_{ij} c_{ij} x_{ij}$ subject to one task per agent
($\sum_{j} x_{ij} = 1$ for each $i$) and one agent per task
($\sum_{i} x_{ij} = 1$ for each $j$). Written this way it is a
$0/1$ integer program with $n^{2}$ binary variables and $n!$
feasible matchings. The structure rescues it. The constraint matrix
of the assignment problem is \engterm{totally unimodular} (every
square submatrix has determinant $0$, $+1$, or $-1$), and a theorem
of integer programming guarantees that an LP over a totally
unimodular constraint matrix with integer right-hand sides has
integer vertices. The LP relaxation, which drops the integrality
and asks only $x_{ij} \geq 0$, therefore returns an integer optimum
on its own; no branching is needed.

A three-agent instance makes the point concrete. The cost matrix
\[
C = \begin{pmatrix}
9 & 2 & 7 \\
6 & 4 & 3 \\
5 & 8 & 1
\end{pmatrix}
\]
gives agent $i$ a cost $c_{ij}$ for task $j$. Brute force checks
$3! = 6$ matchings: $(1{\to}1, 2{\to}2, 3{\to}3)$ costs $9 + 4 + 1 =
14$; $(1{\to}2, 2{\to}3, 3{\to}1)$ costs $2 + 3 + 5 = 10$;
$(1{\to}2, 2{\to}1, 3{\to}3)$ costs $2 + 6 + 1 = 9$;
$(1{\to}3, 2{\to}2, 3{\to}1)$ costs $7 + 4 + 5 = 16$; the remaining
two cost $12$ and $19$. The optimum is $9$, the matching
$1{\to}2$, $2{\to}1$, $3{\to}3$. The Hungarian algorithm of
Kuhn (1955) reaches the same answer in $O(n^{3})$ time without
enumerating matchings: it subtracts the row minima and column
minima to expose zero-cost assignments, then covers the zeros with
a minimal set of lines and adjusts until a complete zero-cost
matching appears. Solving the LP relaxation with a generic LP
solver returns the same integer matching, because total
unimodularity guarantees the relaxation's vertex is integral. The
lesson generalises: before reaching for branch-and-bound, check
whether the constraint matrix has special structure (network flow,
assignment, transportation, interval scheduling) that makes the LP
relaxation integral, in which case the polynomial-time LP solver
is also the integer solver and the NP-hardness of the general case
never bites.

\subsection{Branch-and-bound, recognition level}

\engterm{Branch-and-bound} explores the integer feasible set as a
tree. At the root, solve the LP relaxation (relax integer
variables to continuous). If the LP optimum is integer-feasible,
done. Otherwise pick a fractional variable $x_{i}$ with relaxed
value $x_{i}^{*} \notin \Z$ and branch: the left subtree adds the
constraint $x_{i} \leq \lfloor x_{i}^{*} \rfloor$, the right
subtree adds $x_{i} \geq \lceil x_{i}^{*} \rceil$. Solve the LP
relaxation in each subtree, recurse. Any integer-feasible
solution found anywhere in the tree provides an upper bound on
the optimum (for a minimisation); any LP relaxation provides a
lower bound on the subtree's optimum. A subtree whose LP
relaxation has objective value worse than the best known integer
solution can be \engterm{pruned}: the subtree cannot improve on
the incumbent, so the search skips it.

The art of branch-and-bound is in the bounds. A tight LP
relaxation prunes aggressively and the search runs in seconds; a
loose LP relaxation prunes weakly and the search runs for days
without progress. \engterm{Cutting planes} (Gomory cuts, lift-
and-project, Chvátal-Gomory) tighten the LP relaxation by adding
linear inequalities valid for the integer feasible set but
violated by the LP optimum. The combination of branch-and-bound
with cutting planes and primal heuristics is the working
\engterm{branch-and-cut} engine inside every modern MIP solver.

\subsection{Relaxations}

The general working strategy: replace the hard problem with an
easier problem whose optimum bounds the hard problem's optimum.
LP relaxation drops integrality; Lagrangian relaxation moves
hard constraints into the objective with multipliers; semidefinite
relaxation lifts a quadratic problem to a matrix problem (Goemans
and Williamson 1995, MAX-CUT). Each relaxation gives a bound; the
engineer chooses the relaxation whose bound is tight enough for
the application.

\subsection{Approximation algorithms, recognition level}

For NP-hard problems, an algorithm that returns a solution
within a known factor of optimal in polynomial time is called an
\engterm{approximation algorithm}. The greedy algorithm for set
cover achieves a $\ln n$ approximation factor (Chvátal 1979); the
Christofides algorithm for the metric travelling salesman
achieves a $3/2$ factor (Christofides 1976). The engineer who
needs guaranteed performance on a NP-hard problem at scale uses
an approximation algorithm and reports the approximation factor
along with the solution. The full treatment is in
\textcite{text:cormen-clrs} and \textcite{text:kleinberg-tardos};
this chapter establishes the recognition.

\section{Case study: sizing a cantilever truss}\pathtag{standard}

The chapter's apparatus earns its place when it carries a real
engineering problem from a plain-English brief through a formal
program, a solution, and an interpretation that changes a design
decision. We work a structural-sizing problem end to end, choosing
a truss small enough to solve by hand and rich enough to exercise
the formulation, the constraint reading, the active-set logic, and
the failure check.

\subsection{The brief}

A six-member cantilever truss carries a vertical tip load $P =
40\,\si{\kilo\newton}$ at its lower free node
(\cref{fig:vol02:ch15:truss-sizing}). The bay is square with side
$\ell = 3\,\si{\meter}$, so the two diagonals have length
$\ell \sqrt{2} \approx 4.24\,\si{\meter}$ and the four
chord-and-vertical members have length $\ell$. The members are
steel with allowable stress $\sigma_{\mathrm{allow}} =
160\,\si{\mega\pascal}$ (a representative working value for
structural steel with a safety factor already applied, current as
of 2025) and density $\rho = 7850\,\si{\kilogram\per\meter\cubed}$.
The fabricator can supply any cross-sectional area above a
manufacturing floor $A_{\min} = 2\,\si{\centi\meter\squared}$ (below
this the members are too slender to handle and weld reliably). The
design question in plain English: choose the six member areas to
build the lightest truss that does not overstress any member.

\begin{figure}[ht]
\centering
\begin{tikzpicture}[x=1.0cm, y=1.0cm, font=\small]

% A statically determinate cantilever truss: two top nodes, two bottom nodes,
% fixed wall at left, downward load at the free right end.
% Node coordinates (metres scaled): A=(0,2) B=(3,2) wall-top, C=(0,0) D=(3,0).
\coordinate (A) at (0, 2);
\coordinate (B) at (3, 2);
\coordinate (C) at (0, 0);
\coordinate (D) at (3, 0);

% Wall hatch on the left.
\draw[thick] (-0.25, -0.4) -- (-0.25, 2.4);
\foreach \yy in {-0.3, 0.0, 0.3, 0.6, 0.9, 1.2, 1.5, 1.8, 2.1} {
  \draw[thin] (-0.25, \yy) -- (-0.55, {\yy + 0.25});
}

% Members.
\draw[engblue, very thick] (A) -- (B) node[midway, above, font=\scriptsize] {$1$};
\draw[engblue, very thick] (C) -- (D) node[midway, below, font=\scriptsize] {$2$};
\draw[engblue, very thick] (A) -- (C) node[midway, left, font=\scriptsize] {$3$};
\draw[engblue, very thick] (B) -- (D) node[midway, right, font=\scriptsize] {$4$};
\draw[engblue, very thick] (A) -- (D) node[pos=0.62, above right, font=\scriptsize] {$5$};
\draw[engblue, very thick] (C) -- (B) node[pos=0.40, below right, font=\scriptsize] {$6$};

% Pinned supports at the wall (A and C).
\foreach \n in {A, C} {
  \fill[engblack] (\n) circle (2.2pt);
}
\node[anchor=east, font=\scriptsize] at ($(A) + (-0.15, 0.18)$) {pin};
\node[anchor=east, font=\scriptsize] at ($(C) + (-0.15, -0.18)$) {pin};

% Free nodes B, D.
\foreach \n in {B, D} {
  \fill[engred] (\n) circle (2.2pt);
}

% Applied load at D (downward).
\draw[->, very thick, engred] (D) -- ($(D) + (0, -1.1)$)
  node[anchor=north, font=\scriptsize] {$P$};

% Span dimension.
\draw[<->, thin, enggray] (0, -0.9) -- (3, -0.9);
\node[enggray, anchor=north, font=\scriptsize] at (1.5, -0.95) {$\ell$};

\end{tikzpicture}
\caption[A six-member cantilever truss for sizing]{The six-member
cantilever truss of the sizing case study. Two pinned supports at
the wall (left) carry a tip load $P$ applied at the lower free node.
Each member $j$ carries an axial force $N_{j}$ fixed by static
equilibrium; the design variables are the six cross-sectional areas
$A_{j}$. The optimisation minimises total material volume
$\sum_{j} A_{j} L_{j}$ subject to a stress limit
$\abs{N_{j}}/A_{j} \leq \sigma_{\mathrm{allow}}$ on each member and a
minimum-area floor $A_{j} \geq A_{\min}$. Because the forces are
independent of the areas in a statically determinate truss, the
stress constraints decouple and the optimal areas are read off one
member at a time, which makes the example a clean bridge from
formulation to closed-form solution to interpretation.}
\label{fig:vol02:ch15:truss-sizing}
\end{figure}


\subsection{Formulation}

The decision variable is the area vector $\vect{A} = (A_{1},
\ldots, A_{6})$, with $A_{j}$ the cross-sectional area of member
$j$. The objective is total material volume (mass divided by the
common density, so minimising volume minimises mass):
\[
\min_{\vect{A}} \;\; V(\vect{A}) = \sum_{j=1}^{6} A_{j} L_{j},
\]
with $L_{j}$ the length of member $j$. The constraints are a stress
limit on each member and the manufacturing floor:
\[
\frac{\abs{N_{j}}}{A_{j}} \leq \sigma_{\mathrm{allow}}, \qquad
A_{j} \geq A_{\min}, \qquad j = 1, \ldots, 6,
\]
where $N_{j}$ is the axial force in member $j$ (tension positive,
compression negative). The stress constraint rearranges to the
linear form $A_{j} \geq \abs{N_{j}} / \sigma_{\mathrm{allow}}$, so
the whole problem is a linear program in $\vect{A}$ once the forces
$N_{j}$ are known.

The structural fact that makes this clean: the truss is statically
determinate, so the member forces are fixed by equilibrium alone
and do not depend on the areas. We can solve for the forces once,
up front, and the optimisation then sees them as constants. In a
statically indeterminate structure the forces would redistribute as
the areas change, coupling the equilibrium to the design variables
and turning the problem nonlinear; that coupling is the subject of
structural-optimisation courses and Volume IV. Here the
determinacy decouples force from area and hands us a transparent
LP.

\subsection{The forces from equilibrium}

Resolve the tip load through the truss by method of joints. The
load $P$ at the lower free node is carried by the lower chord
(member 2) in tension and the rising diagonal (member 5) in
compression. Taking the geometry of \cref{fig:vol02:ch15:truss-sizing}
with the standard joint equations, the member forces work out to
the following magnitudes for $P = 40\,\si{\kilo\newton}$ and a
square bay (so each diagonal makes $45^{\circ}$):
\[
\begin{aligned}
\abs{N_{1}} &= P = 40\,\si{\kilo\newton} &\text{(top chord)}, \\
\abs{N_{2}} &= P = 40\,\si{\kilo\newton} &\text{(bottom chord)}, \\
\abs{N_{3}} &= P = 40\,\si{\kilo\newton} &\text{(rear vertical)}, \\
\abs{N_{4}} &= 0 &\text{(front vertical)}, \\
\abs{N_{5}} &= P\sqrt{2} \approx 56.6\,\si{\kilo\newton}
  &\text{(loaded diagonal)}, \\
\abs{N_{6}} &= 0 &\text{(idle diagonal)}.
\end{aligned}
\]
Two members carry no force. Member 4, the front vertical, and
member 6, the idle diagonal, are \engterm{zero-force members} for
this single load case: the equilibrium of their end joints is
satisfied without them. Zero-force members are a standard feature
of statically determinate trusses, and they pose the first design
judgment the optimisation will surface.

\subsection{Solving the program}

Because the stress constraints decouple (each member's area appears
in exactly one stress constraint and in the objective), the LP
separates into six one-variable problems. For each member the
objective $A_{j} L_{j}$ increases in $A_{j}$, so the optimiser
drives each area to the smallest value its constraints allow:
\[
A_{j}^{*} = \max\!\left(
\frac{\abs{N_{j}}}{\sigma_{\mathrm{allow}}}, \; A_{\min}
\right).
\]
The stress-limited area for a member carrying $40\,\si{\kilo\newton}$
is $\abs{N_{j}} / \sigma_{\mathrm{allow}} = 40 \times 10^{3} /
(160 \times 10^{6}) = 2.5 \times 10^{-4}\,\si{\meter\squared} =
2.5\,\si{\centi\meter\squared}$, comfortably above the
$2\,\si{\centi\meter\squared}$ floor, so the stress constraint
binds. The loaded diagonal needs $56.6 \times 10^{3} / (160 \times
10^{6}) = 3.54 \times 10^{-4}\,\si{\meter\squared} =
3.54\,\si{\centi\meter\squared}$. The two zero-force members carry
nothing, so their stress-limited area is zero and the
manufacturing floor binds instead: $A_{4}^{*} = A_{6}^{*} =
A_{\min} = 2\,\si{\centi\meter\squared}$. The optimal areas are
\[
\vect{A}^{*} = (2.5, \; 2.5, \; 2.5, \; 2.0, \; 3.54, \; 2.0)
\;\si{\centi\meter\squared}.
\]
The optimal volume is $V^{*} = \sum_{j} A_{j}^{*} L_{j} =
(2.5 + 2.5 + 2.5)(3)(10^{-4}) + (2.0)(3)(10^{-4}) +
(3.54)(4.24)(10^{-4}) + (2.0)(4.24)(10^{-4})$, which evaluates to
about $5.2 \times 10^{-3}\,\si{\meter\cubed}$, a mass of $\rho V^{*}
\approx 41\,\si{\kilogram}$.

\subsection{The active set and the shadow prices}

The optimum partitions the members by which constraint binds, and
the partition is the engineering content. Members 1, 2, 3, and 5
are \engterm{stress-limited}: their stress constraint is active,
their area is set by the force they carry, and their shadow price
on the stress limit is positive. Members 4 and 6 are
\engterm{floor-limited}: the stress constraint is slack (they carry
no force), the manufacturing floor is active, and the shadow price
on their stress limit is zero. The shadow price on the
stress-limit constraint of a loaded member is the volume saved per
unit relaxation of the allowable stress: $\partial V^{*} /
\partial \sigma_{\mathrm{allow}}$ for member $j$ equals
$-\abs{N_{j}} L_{j} / \sigma_{\mathrm{allow}}^{2}$, so a stronger
steel grade buys the most weight saving on the members carrying the
most force over the longest length, which is the loaded diagonal.
The duality reading turns directly into a procurement decision:
spending on higher-strength steel pays back on members 1, 2, 3,
and 5 and pays back nothing on members 4 and 6, whose size is fixed
by manufacturability rather than by strength.

\subsection{Interpretation and the design decision}

The optimisation has surfaced a decision the brief did not state.
Members 4 and 6 carry no load for this single tip-load case and sit
at the manufacturing floor purely because the fabricator cannot
make them thinner. Three readings compete. First, the optimiser is
correct and the members should be built at the floor, because a
real structure sees more than one load case (wind, a load at the
upper node, a construction load) and the zero-force members carry
force under those cases; sizing them to the floor is cheap
insurance. Second, if the single tip load is genuinely the only
case, members 4 and 6 are candidates for removal, which would
change the topology of the truss and is a \engterm{topology
optimisation} question one level above this sizing problem.
Third, the zero-force result is an artifact of the idealised
pin-jointed model, and a real welded joint carries moment, so the
members are not truly idle and the floor area is doing real work.
The optimisation cannot choose among these readings; it can only
report that the constraint binding on members 4 and 6 is the
manufacturing floor, not strength, which is exactly the flag the
engineer needs to ask the right follow-up question. This is the
chapter's thesis in miniature: the optimiser solved the LP that was
written down, and the engineering judgment lives in reading the
active set and deciding whether the formulation captured every load
case that matters. The full multi-load extension, in which each of
several load cases contributes its own stress constraint per member
and the area must satisfy the worst of them, stays an LP of the same
shape with more constraint rows, and a runnable LP solve of the
single-case problem with its shadow-price table is in
\repopath{volumes/02-form/ch15-optimisation/code/truss_sizing_lp.py}.

\section{Nonlinear least squares: Gauss-Newton and
Levenberg-Marquardt}\pathtag{standard}

The truss case was a linear program because the constraints rearranged
to linear inequalities in the design variable. Most model calibration
is not so kind. When the quantity being fit depends nonlinearly on the
parameters, the least-squares objective is nonconvex in general, and
the linear-algebra recipe of the normal equations no longer reaches the
optimum in one solve. The structure that survives is that the objective
is still a sum of squared residuals, and that structure buys a method
faster than plain gradient descent and cheaper than full Newton. This
section develops the method, and the case study that follows runs it
end to end on a real calibration brief.

\subsection{The nonlinear least-squares problem}

A model $\hat{y}(t; \boldsymbol{\theta})$ predicts an observable from an
input $t$ and a parameter vector $\boldsymbol{\theta} \in \R^{p}$.
Given measurements $(t_{i}, y_{i})$ for $i = 1, \ldots, N$, the
\engterm{residual} of the $i$-th datum is $r_{i}(\boldsymbol{\theta}) =
y_{i} - \hat{y}(t_{i}; \boldsymbol{\theta})$, and the calibration
minimises the sum of squared residuals
\[
\min_{\boldsymbol{\theta}} \;\; F(\boldsymbol{\theta}) =
\tfrac{1}{2} \sum_{i=1}^{N} r_{i}(\boldsymbol{\theta})^{2}
= \tfrac{1}{2} \norm{\vect{r}(\boldsymbol{\theta})}_{2}^{2},
\]
with $\vect{r} = (r_{1}, \ldots, r_{N})^{\transpose}$. When
$\hat{y}$ is linear in $\boldsymbol{\theta}$ the residual is affine,
$F$ is a convex quadratic, and the problem is the ordinary least
squares of Section 15.2. When $\hat{y}$ is nonlinear in
$\boldsymbol{\theta}$, the residual is nonlinear, $F$ is generally
nonconvex, and the problem can have several local minima.

The gradient and Hessian of $F$ follow from the chain rule. Let $J \in
\R^{N \times p}$ be the \engterm{Jacobian} of the residual,
$J_{ij} = \partial r_{i} / \partial \theta_{j}$. Then
\[
\grad F(\boldsymbol{\theta}) = J^{\transpose} \vect{r},
\qquad
\nabla^{2} F(\boldsymbol{\theta}) = J^{\transpose} J +
\sum_{i=1}^{N} r_{i} \, \nabla^{2} r_{i}.
\]
The Hessian has two parts: the \engterm{Gauss-Newton term}
$J^{\transpose} J$, built from first derivatives alone, and the
\engterm{second-order term} $\sum_{i} r_{i} \nabla^{2} r_{i}$, which
weights each residual's curvature by the residual itself. The
observation that makes the method: near a good fit the residuals
$r_{i}$ are small, so the second-order term is small, and
$J^{\transpose} J$ approximates the full Hessian well. This is the
\engterm{small-residual regime}, and it is the regime of every
well-posed calibration where the model can actually represent the data.

\subsection{The Gauss-Newton step}

\begin{definition}[Gauss-Newton iteration]
The \engterm{Gauss-Newton} step at $\boldsymbol{\theta}_{k}$ solves the
linear least-squares subproblem obtained by linearising the residual,
$\vect{r}(\boldsymbol{\theta}_{k} + \vect{p}) \approx \vect{r}_{k} +
J_{k} \vect{p}$, and minimising $\tfrac{1}{2} \norm{\vect{r}_{k} +
J_{k} \vect{p}}_{2}^{2}$ over the step $\vect{p}$:
\[
\vect{p}_{k} = -\,(J_{k}^{\transpose} J_{k})^{-1} J_{k}^{\transpose}
\vect{r}_{k},
\qquad
\boldsymbol{\theta}_{k+1} = \boldsymbol{\theta}_{k} + \vect{p}_{k}.
\]
\end{definition}

The Gauss-Newton step is the Newton step with the second-order Hessian
term dropped. It costs one Jacobian evaluation and one linear
least-squares solve of size $N \times p$ per iteration, never a second
derivative of the model. In the small-residual regime it converges
nearly quadratically, the same fast tail as Newton, because the dropped
term vanishes as the residuals shrink. The step is the solution of the
normal equations $J_{k}^{\transpose} J_{k} \vect{p}_{k} = -
J_{k}^{\transpose} \vect{r}_{k}$, and as in linear least squares the
numerically careful route solves the linearised subproblem through a QR
or SVD factorisation of $J_{k}$ rather than forming $J_{k}^{\transpose}
J_{k}$ and squaring the conditioning.

Gauss-Newton has one failure mode that bites in practice: far from the
optimum, or when the Jacobian is rank-deficient, $J_{k}^{\transpose}
J_{k}$ is singular or nearly so, the step blows up, and the iteration
diverges. The repair is the same Hessian modification that Section 15.5
introduced for Newton's method.

\subsection{Levenberg-Marquardt}

\begin{definition}[Levenberg-Marquardt iteration]
The \engterm{Levenberg-Marquardt} step damps the Gauss-Newton solve
with a nonnegative multiple of the identity,
\[
\vect{p}_{k} = -\,(J_{k}^{\transpose} J_{k} + \mu_{k} I)^{-1}
J_{k}^{\transpose} \vect{r}_{k},
\qquad \mu_{k} \geq 0,
\]
and adapts $\mu_{k}$: the damping shrinks (toward the Gauss-Newton step)
when a step succeeds and grows (toward a short gradient step) when a
step fails.
\end{definition}

The damping term is exactly the trust-region modification of
\cref{fig:vol02:ch15:trust-region}: $\mu_{k}$ is the Lagrange
multiplier of a trust-region constraint $\norm{\vect{p}} \leq
\Delta_{k}$ on the Gauss-Newton subproblem. Small $\mu_{k}$ gives a
large trust region and the Gauss-Newton step; large $\mu_{k}$ gives a
small trust region and a step of length $\approx \norm{J_{k}^{\transpose}
\vect{r}_{k}} / \mu_{k}$ in the steepest-descent direction. The
adaptive damping interpolates between the two regimes automatically,
which lets Levenberg-Marquardt converge from a poor starting point where
pure Gauss-Newton diverges, while retaining the fast tail near the
optimum where the damping has shrunk to near zero. Levenberg-Marquardt
is the default method inside every curve-fitting routine an engineer
reaches for, from \repopath{scipy.optimize.curve_fit} to the fitting
dialog of a spreadsheet, and the engineer now knows what its single
tuning knob, the damping, controls.

\subsection{Case study: calibrating a battery-discharge decay model}

A battery test rig records the open-circuit voltage of a cell as it
relaxes after a load is removed. The physics of the relaxation is a
diffusion-and-relaxation process whose voltage recovery is well
described by a single-exponential model with an offset,
\[
\hat{V}(t; A, k, C) = A \, e^{-k t} + C,
\]
where $t$ is time since the load was removed (in hours), $C$ is the
settled rest voltage, $A$ is the recovery amplitude, and $k$ is the
relaxation rate. The engineer needs all three parameters: $C$ for the
state-of-charge lookup, $A$ and $k$ for the cell's internal-resistance
and diffusion characterisation. The data are thirty voltage samples
over five hours, each carrying measurement noise of about
$50\,\si{\milli\volt}$ standard deviation from the data-acquisition
hardware. The calibration question in plain English: find the
$(A, k, C)$ whose decay curve best fits the recorded relaxation in the
least-squares sense, and report whether the fit is trustworthy.

\paragraph{The brief becomes a program.}
The decision variable is $\boldsymbol{\theta} = (A, k, C)$. The
residual of the $i$-th sample is $r_{i} = V_{i} - \hat{V}(t_{i};
\boldsymbol{\theta})$, and the calibration is the nonlinear
least-squares program $\min_{\boldsymbol{\theta}} \tfrac{1}{2}
\sum_{i} r_{i}^{2}$. The model is nonlinear in $k$ (the parameter sits
in an exponent), so the problem is genuinely nonlinear least squares,
not the linear case of Section 15.2.

\paragraph{The Jacobian.}
The Jacobian columns are the partial derivatives of the residual with
respect to each parameter. Writing $e_{i} = e^{-k t_{i}}$,
\[
\pdv{r_{i}}{A} = -e_{i}, \qquad
\pdv{r_{i}}{k} = A \, t_{i} \, e_{i}, \qquad
\pdv{r_{i}}{C} = -1,
\]
so the $i$-th row of $J$ is $(-e_{i}, \; A t_{i} e_{i}, \; -1)$. The
Jacobian is the entire input the method needs beyond the residual
itself; no second derivative of the model is ever formed.

\paragraph{Running the iteration.}
From a reasonable start $\boldsymbol{\theta}_{0} = (2.0, 0.3, 0.5)$,
Gauss-Newton drives the residual norm from $\norm{\vect{r}_{0}} \approx
1.97$ down through $1.85$, $0.43$, to the converged $0.216$ in four
iterations, the last two iterations agreeing to six figures, the
near-quadratic tail the small-residual theory predicts. The converged
parameters are $\boldsymbol{\theta}^{*} \approx (2.606, \; 0.571, \;
0.418)$, close to the true generating values $(2.6, 0.55, 0.40)$ within
the noise. Levenberg-Marquardt, started from the deliberately poor
point $\boldsymbol{\theta}_{0} = (1.0, 0.1, 0.0)$ where pure
Gauss-Newton produces a rank-deficient $J^{\transpose} J$ and a
singular solve, instead damps its way down ($4.27 \to 3.62 \to 1.79 \to
0.66 \to 0.216$) and reaches the identical optimum in five iterations,
which is the robustness the damping was added to provide. Plain
gradient descent on the same objective, at the largest stable step, is
still far from the optimum after thousands of iterations, the geometric
crawl of \cref{fig:vol02:ch15:gauss-newton-fit}(b).

\begin{figure}[ht]
\centering
\begin{tikzpicture}[x=1cm, y=1cm, font=\small]

% Left panel: the fitted exponential-decay curve through scattered data.
\begin{scope}
  \draw[->, thick] (-0.2, 0) -- (5.6, 0)
    node[right, font=\scriptsize] {$t\ [\si{\hour}]$};
  \draw[->, thick] (0, -0.2) -- (0, 3.6)
    node[above, font=\scriptsize] {$y$};

  % Model y = A exp(-k t) + C with A=2.6, k=0.55, C=0.4 (axis-scaled).
  % Draw the smooth fitted curve.
  \draw[engblue, very thick, samples=80, domain=0:5.2, smooth]
    plot (\x, {0.4 + 2.6*exp(-0.55*\x)});

  % Scattered data points (noisy samples around the curve).
  \fill[engblack] (0.2, 3.05) circle (1.3pt);
  \fill[engblack] (0.6, 2.50) circle (1.3pt);
  \fill[engblack] (1.0, 2.18) circle (1.3pt);
  \fill[engblack] (1.5, 1.70) circle (1.3pt);
  \fill[engblack] (2.0, 1.32) circle (1.3pt);
  \fill[engblack] (2.6, 1.08) circle (1.3pt);
  \fill[engblack] (3.2, 0.86) circle (1.3pt);
  \fill[engblack] (3.9, 0.66) circle (1.3pt);
  \fill[engblack] (4.6, 0.58) circle (1.3pt);
  \fill[engblack] (5.1, 0.50) circle (1.3pt);

  \node[engblue, anchor=west, font=\scriptsize] at (2.7, 1.95)
    {$\hat{y}(t) = A\,e^{-kt} + C$};
  \node[anchor=north, font=\scriptsize] at (2.6, -0.55) {(a) fitted decay curve};
\end{scope}

% Right panel: residual-norm convergence, Gauss-Newton vs gradient descent.
\begin{scope}[xshift=7.5cm]
  \draw[->, thick] (-0.2, 0) -- (4.6, 0)
    node[right, font=\scriptsize] {iteration $k$};
  \draw[->, thick] (0, -0.2) -- (0, 3.6)
    node[above, font=\scriptsize] {$\log_{10}\norm{\vect{r}_{k}}$};

  % Gauss-Newton / LM: fast (near-quadratic) drop.
  \draw[engred, very thick]
    (0, 3.3) -- (0.7, 2.4) -- (1.4, 1.2) -- (2.1, 0.35) -- (2.8, 0.10)
    -- (3.6, 0.06);
  \foreach \p in {(0,3.3),(0.7,2.4),(1.4,1.2),(2.1,0.35),(2.8,0.10),(3.6,0.06)}
    \fill[engred] \p circle (1.4pt);
  \node[engred, anchor=south west, font=\scriptsize] at (2.2, 0.45)
    {Levenberg-Marquardt};

  % Gradient descent: slow geometric drop.
  \draw[enggray, very thick]
    (0, 3.3) -- (0.7, 3.05) -- (1.4, 2.82) -- (2.1, 2.60) -- (2.8, 2.40)
    -- (3.6, 2.18) -- (4.3, 1.98);
  \node[enggray, anchor=south west, font=\scriptsize] at (2.0, 2.55)
    {gradient descent};

  \node[anchor=north, font=\scriptsize] at (2.2, -0.55)
    {(b) residual norm vs iteration};
\end{scope}

\end{tikzpicture}
\caption[Gauss-Newton calibration of a decay model]{Nonlinear
least-squares calibration of the exponential-decay model $\hat{y}(t)
= A\,e^{-kt} + C$. Panel (a): the fitted curve threads the scattered
measurements, each datum contributing a residual $r_{i} = y_{i} -
\hat{y}(t_{i})$ that the calibration drives down in the sum-of-squares
sense. Panel (b): the residual norm against iteration on a log scale.
Levenberg-Marquardt, which uses the Gauss-Newton approximation
$J^{\transpose}J$ to the Hessian, reaches the optimum in a handful of
iterations with the near-quadratic tail of a second-order method;
plain gradient descent on the same objective crawls down the
ill-conditioned valley at the geometric rate set by the condition
number of $J^{\transpose}J$. The contrast is the practical reason
every curve-fitting library defaults to Levenberg-Marquardt rather
than to a first-order method.}
\label{fig:vol02:ch15:gauss-newton-fit}
\end{figure}


\paragraph{Is the fit trustworthy?}
The condition number of the Gauss-Newton Hessian $J^{\transpose} J$ at
the optimum is about $73$, a well-conditioned problem: the three
parameters are jointly identifiable from the data, and a small change
in the measurements moves the estimate only modestly. The diagnostic
matters because nonlinear least squares can return a fit that matches
the data while leaving a parameter all but undetermined, which
shows up as a near-singular $J^{\transpose} J$ and a huge condition
number. The covariance of the estimated parameters is
$\hat{\sigma}^{2} (J^{\transpose} J)^{-1}$, where $\hat{\sigma}^{2} =
\norm{\vect{r}^{*}}^{2} / (N - p)$ is the residual variance, so the
same $J^{\transpose} J$ that drove the iteration also delivers the
error bars on the answer. The engineer reports the parameters with
their standard errors, not the parameters alone, because a calibration
without uncertainty is a number without a tolerance. The residual
variance here, $\hat{\sigma}^{2} = 0.216^{2} / 27 \approx 1.7 \times
10^{-3}\,\si{\volt\squared}$, gives a residual standard deviation of
about $42\,\si{\milli\volt}$, consistent with the
$50\,\si{\milli\volt}$ measurement noise the brief stated, which is the
final check that the model has not under- or over-fit: a residual
standard deviation far below the known noise floor would signal
overfitting, and one far above it would signal a model too simple for
the data. The full calibration, with all three solvers and the
condition-number and covariance diagnostics, is in
\repopath{volumes/02-form/ch15-optimisation/code/gauss_newton_calibration.py}.

\paragraph{What the case study teaches.}
The nonlinear-least-squares calibration is the second archetype of the
chapter, distinct from the truss LP in three ways that recur across
engineering. The objective is nonconvex, so the answer is a local
optimum and the engineer must consider the start point and, for a hard
fit, multiple starts. The structure (sum of squares) buys a
second-order method without a second derivative, the single most
useful structural recognition in calibration. And the same linear
algebra that computes the step also certifies the answer through the
condition number and the parameter covariance, so the optimum arrives
with its uncertainty attached. The truss case study showed how to read
an active set; this one shows how to read a Jacobian.

\begin{mastery}
\textbf{Gauss-Newton as a fixed-point map, and when it diverges.} The
Gauss-Newton iteration is a fixed-point map $\boldsymbol{\theta}_{k+1}
= G(\boldsymbol{\theta}_{k})$ whose convergence is governed by the
spectral radius of the iteration's Jacobian at the optimum. A careful
expansion shows that map's Jacobian is $(J^{\transpose} J)^{-1} S$,
where $S = \sum_{i} r_{i}^{*} \nabla^{2} r_{i}$ is exactly the
second-order Hessian term that Gauss-Newton dropped. The iteration
converges when the spectral radius $\rho\big((J^{\transpose} J)^{-1}
S\big) < 1$, and the convergence is linear with that radius as the
rate. Two regimes follow directly. In the small-residual regime the
residuals $r_{i}^{*}$ are small, $S$ is small, the spectral radius is
near zero, and convergence is fast and effectively quadratic. In the
large-residual regime, where the model cannot fit the data and the
$r_{i}^{*}$ stay large at the optimum, $S$ is large, the spectral
radius can exceed one, and Gauss-Newton diverges even from a good
start. The diagnosis is precise: Gauss-Newton failing to converge on a
problem with large residuals at the optimum is not a bug but the
signature of a model that does not describe the data, and the engineer
should question the model before reaching for a more forgiving solver.
Levenberg-Marquardt's damping rescues convergence by shrinking the
step, but a large-residual fit that needs heavy damping to converge is
telling the engineer the same thing: the model is wrong, not just the
solver.
\end{mastery}

\section{Conic optimisation: SOCP and SDP}\pathtag{mastery}

Between the linearly constrained convex programs of Section 15.3 and
the general convex programs of Section 15.2 sits a structured middle
layer that modern solvers handle in polynomial time: the conic
programs, whose constraints require a vector or a matrix to lie in a
convex cone. Two cones cover most engineering applications, and
recognising them in a problem statement unlocks a solver that a general
nonlinear method would never match.

\subsection{Second-order cone programs}

\begin{definition}[Second-order cone program]
A \engterm{second-order cone program} (SOCP) minimises a linear
objective subject to constraints of the form
\[
\norm{A_{i} \vect{x} + \vect{b}_{i}}_{2} \leq
\vect{c}_{i}^{\transpose} \vect{x} + d_{i}, \qquad i = 1, \ldots, m,
\]
each of which requires the pair $(A_{i} \vect{x} + \vect{b}_{i}, \;
\vect{c}_{i}^{\transpose} \vect{x} + d_{i})$ to lie in the
\engterm{second-order cone} $\mathcal{Q} = \set{(\vect{u}, t) :
\norm{\vect{u}}_{2} \leq t}$, the ice-cream cone in $\R^{n+1}$.
\end{definition}

The constraint $\norm{\vect{u}}_{2} \leq t$ is convex (a norm bounded
by an affine function), so an SOCP is a convex program, solvable by the
interior-point methods of the Section 15.3 mastery box. The class is
strictly larger than QP and strictly smaller than SDP. Engineering
instances are abundant: robust least squares (minimise the worst-case
residual over a bounded perturbation of $A$), the minimum-enclosing-
ball problem, antenna and loudspeaker beamforming (bound the sidelobe
level while steering the main lobe), and any QP, which is an SOCP with
a single cone constraint, since $\tfrac{1}{2} \vect{x}^{\transpose} P
\vect{x} \leq t$ rewrites as $\norm{P^{1/2} \vect{x}}_{2}^{2} \leq 2t$.

\paragraph{Worked SOCP: a robust linear constraint.}
A design must satisfy $\vect{a}^{\transpose} \vect{x} \leq b$, but the
coefficient vector $\vect{a}$ is known only to lie in a ball
$\set{\bar{\vect{a}} + P \vect{u} : \norm{\vect{u}}_{2} \leq 1}$ around
a nominal $\bar{\vect{a}}$ (the columns of $P$ are the uncertainty
directions). Requiring the constraint to hold for every $\vect{a}$ in
the ball is the robust constraint $\max_{\norm{\vect{u}} \leq 1}
(\bar{\vect{a}} + P \vect{u})^{\transpose} \vect{x} \leq b$. The inner
maximum is attained by aligning $\vect{u}$ with $P^{\transpose}
\vect{x}$, giving $\bar{\vect{a}}^{\transpose} \vect{x} +
\norm{P^{\transpose} \vect{x}}_{2} \leq b$, a single second-order cone
constraint. The robust version costs nothing beyond recognising the
cone: the worst-case-over-an-ellipsoid constraint is exactly an SOCP
constraint, and the same interior-point solver handles it. This is the
mechanism behind robust portfolio optimisation, robust beamforming, and
chance-constrained design under Gaussian uncertainty, all of which
reduce a worst-case or probabilistic guarantee to a norm-bounded affine
constraint.

\subsection{Semidefinite programs}

\begin{definition}[Semidefinite program]
A \engterm{semidefinite program} (SDP) minimises a linear function of
a symmetric matrix variable $X$ subject to linear matrix equalities and
the constraint that $X$ be positive semidefinite:
\[
\min_{X = X^{\transpose}} \; \tr(C X)
\quad \text{subject to} \quad
\tr(A_{i} X) = b_{i}, \;\; i = 1, \ldots, m, \quad X \succeq 0.
\]
The constraint $X \succeq 0$ requires $X$ to lie in the
\engterm{positive semidefinite cone}, the convex cone of symmetric
matrices with nonnegative eigenvalues.
\end{definition}

SDP is the most expressive convex class an engineer routinely meets.
Every SOCP is an SDP (a second-order cone constraint embeds in a
semidefinite one), so the hierarchy is LP $\subset$ QP $\subset$ SOCP
$\subset$ SDP, each class solvable in polynomial time by interior-point
methods and each strictly more expressive than the last. SDP
applications include control synthesis through linear matrix
inequalities (Volume IX), optimal experiment design, the
nearest-correlation-matrix problem, and, most famously, tight
relaxations of hard combinatorial problems.

\paragraph{Worked SDP: a small semidefinite relaxation.}
Consider the problem of finding the largest $\gamma$ such that the
matrix
\[
M(\gamma) = \begin{pmatrix} 2 & 1 \\ 1 & \gamma \end{pmatrix}
\succeq 0.
\]
The positive-semidefinite constraint on a symmetric $2 \times 2$ matrix
is, by the Sylvester criterion of Volume II Chapter 9, the pair of
conditions: the leading entry $2 \geq 0$ (satisfied) and the
determinant $2\gamma - 1 \geq 0$. The feasible set is therefore $\gamma
\geq 1/2$, and the smallest feasible $\gamma$ is $1/2$, at which $M$ is
singular with a zero eigenvalue. The example is a one-variable SDP, and
it shows the geometry the cone constraint encodes: the boundary of the
semidefinite cone is the locus where an eigenvalue hits zero, and an
optimum on the boundary is a matrix with a nontrivial null vector,
which in control synthesis is a marginally stable mode and in
relaxation is a tight bound. The full power of SDP appears when the
matrix variable is large and the eigenvalue constraint couples many
scalar decisions, but the two-by-two case already shows that the cone
constraint is the smooth convex surrogate for a determinant or
eigenvalue requirement that would otherwise be a nonconvex polynomial
condition.

\begin{mastery}
\textbf{The SDP relaxation of MAX-CUT.} The single most influential
use of semidefinite programming is the Goemans-Williamson relaxation
of MAX-CUT, the NP-hard problem of partitioning a graph's vertices into
two sets to maximise the number of edges crossing the partition.
Assign each vertex $i$ a label $x_{i} \in \set{-1, +1}$; an edge
$(i, j)$ is cut when $x_{i} \neq x_{j}$, i.e.\ when $(1 - x_{i}
x_{j})/2 = 1$. MAX-CUT is $\max_{x \in \set{-1,+1}^{n}} \tfrac{1}{2}
\sum_{(i,j) \in E} (1 - x_{i} x_{j})$, a quadratic objective over the
discrete cube, NP-hard. The relaxation replaces each scalar label
$x_{i}$ with a unit vector $\vect{v}_{i} \in \R^{n}$ and each product
$x_{i} x_{j}$ with the inner product $\vect{v}_{i}^{\transpose}
\vect{v}_{j}$. Collecting the inner products into the Gram matrix
$X_{ij} = \vect{v}_{i}^{\transpose} \vect{v}_{j}$, the constraints
become $X \succeq 0$ (a Gram matrix is positive semidefinite) and
$X_{ii} = 1$ (unit vectors), and the objective is linear in $X$. The
relaxation is therefore an SDP, solvable in polynomial time, and its
optimum upper-bounds the true MAX-CUT. Goemans and Williamson then
round the SDP solution back to a cut by slicing the unit vectors with a
random hyperplane, and they prove the expected cut is at least $0.878$
times the SDP optimum, hence at least $0.878$ times the true optimum: a
polynomial-time algorithm with a provable approximation ratio for an
NP-hard problem, obtained by lifting a discrete quadratic to a
continuous matrix program. The construction is the template for
semidefinite relaxation across combinatorial optimisation, and it is
the reason an engineer facing a hard quadratic Boolean problem checks
whether an SDP relaxation gives a usable bound before settling for a
heuristic.
\end{mastery}

\section{Multi-objective optimisation: the Pareto front}\pathtag{standard}

The chapter has so far assumed a single scalar objective. Real
engineering decisions trade several objectives that pull against each
other: a structure that is both light and cheap, a controller that is
both fast and energy-frugal, a portfolio that is both high-return and
low-risk. When the objectives genuinely conflict, no single design is
best on all of them at once, and the right output of the optimisation
is a frontier of trade-offs rather than a single point.

\subsection{Dominance and Pareto optimality}

\begin{definition}[Pareto dominance and Pareto optimality]
For a vector of objectives $\vect{f}(\vect{x}) = (f_{1}(\vect{x}),
\ldots, f_{q}(\vect{x}))$ to be minimised, a design $\vect{x}$
\engterm{dominates} a design $\vect{y}$ if $f_{k}(\vect{x}) \leq
f_{k}(\vect{y})$ for every objective $k$ and $f_{k}(\vect{x}) <
f_{k}(\vect{y})$ for at least one. A feasible design is
\engterm{Pareto optimal} if no feasible design dominates it. The
\engterm{Pareto front} is the image, in objective space, of the set of
Pareto-optimal designs.
\end{definition}

A dominated design is one a rational engineer never chooses: some other
feasible design is at least as good on every objective and strictly
better on one. The Pareto-optimal designs are the ones where every
remaining improvement in one objective forces a sacrifice in another,
and the Pareto front (\cref{fig:vol02:ch15:pareto-front}) is the set of
those irreducible trade-offs. The engineer's deliverable in a
multi-objective problem is the front, because choosing a single point
on it is a value judgment (how much mass is one dollar worth?) that the
optimisation cannot make and the decision-maker must.

\begin{figure}[ht]
\centering
\begin{tikzpicture}[x=1cm, y=1cm, font=\small]

% Bi-objective design space: cost on x, mass on y. The achievable
% region is a blob; its lower-left boundary is the Pareto front.
% Front curve: y = 0.6 + 6/(x+0.2) on the window x in [1.0, 5.6],
% which gives y from ~5.6 down to ~1.6, kept inside the canvas.
\draw[->, thick] (-0.2, 0) -- (6.4, 0)
  node[right, font=\scriptsize] {objective 1: cost};
\draw[->, thick] (0, -0.2) -- (0, 4.6)
  node[above, font=\scriptsize] {objective 2: mass};

% Clip the shaded achievable region to a tidy box so it does not
% collide with the vertical axis.
\begin{scope}
  \clip (1.0, 0.0) rectangle (6.0, 4.3);
  \fill[engblue!10]
    plot[smooth, domain=1.45:6.0, samples=60] (\x, {0.6 + 6/(\x+0.2)})
    -- (6.0, 4.3) -- (1.45, 4.3) -- cycle;
\end{scope}

% The Pareto front: the lower-left boundary (convex, decreasing).
\draw[engred, very thick, smooth, samples=80, domain=1.5:5.9]
  plot (\x, {0.6 + 6/(\x+0.2)});
\node[engred, anchor=north west, font=\scriptsize] at (2.0, 1.5)
  {Pareto front};

% A dominated interior point.
\fill[enggray] (3.7, 3.2) circle (2pt);
\node[enggray, anchor=south west, font=\scriptsize] at (3.75, 3.25)
  {dominated};
% Arrows: it is dominated by front points (less of both objectives).
\draw[->, enggray, thin] (3.7, 3.2) -- (3.7, 2.18);
\draw[->, enggray, thin] (3.7, 3.2) -- (2.45, 3.2);

% Three scalarisation optima sitting on the front.
% Front y at x=1.7: 0.6 + 6/1.9 = 3.76
\fill[engblack] (1.7, 3.76) circle (1.6pt);
\node[engblack, anchor=west, font=\scriptsize] at (1.78, 3.86)
  {$w$ small};
% Front y at x=2.8: 0.6 + 6/3.0 = 2.6
\fill[engblack] (2.8, 2.6) circle (1.6pt);
\node[engblack, anchor=south west, font=\scriptsize] at (2.88, 2.66)
  {$w$ mid};
% Front y at x=4.8: 0.6 + 6/5.0 = 1.8
\fill[engblack] (4.8, 1.8) circle (1.6pt);
\node[engblack, anchor=south west, font=\scriptsize] at (4.88, 1.86)
  {$w$ large};

\end{tikzpicture}
\caption[The Pareto front of a bi-objective design]{The Pareto front
of a two-objective design problem trading cost against mass. The shaded
region is the achievable set: every feasible design realises some
$(\text{cost}, \text{mass})$ pair inside it. The lower-left boundary,
drawn in red, is the Pareto front, the set of designs that cannot be
improved in one objective without worsening the other. The grey
interior point is dominated: a front design exists with both lower cost
and lower mass, so no rational engineer would choose the interior
point. Scalarising the two objectives into a single weighted sum and
minimising for several weights $w$ recovers individual front points,
black markers, with small $w$ favouring low cost and large $w$
favouring low mass. The weighted-sum sweep traces the convex part of
the front; non-convex pieces of a front require other scalarisations,
which is the reason the engineer reports the whole front rather than a
single scalar optimum.}
\label{fig:vol02:ch15:pareto-front}
\end{figure}


\subsection{Scalarisation: recovering the front}

The working method for tracing a Pareto front is \engterm{weighted-sum
scalarisation}: combine the objectives into a single weighted objective
$\sum_{k} w_{k} f_{k}(\vect{x})$ with weights $w_{k} > 0$, solve the
resulting single-objective problem, and sweep the weights. Each weight
vector returns one Pareto-optimal point, and varying the weights traces
the front.

\paragraph{Worked example: a two-objective sweep.}
Trade mass against cost for a design whose two objectives, as functions
of a single scalar design parameter $d \in [1, 5]$ (a wall thickness,
say), are $f_{1}(d) = d$ (cost, rising with material) and $f_{2}(d) =
9/d$ (mass, falling as the wall thickens and the structure needs less
reinforcing elsewhere). The scalarised objective is $w_{1} d + w_{2}
(9/d)$, minimised at $\dv{}{d}(w_{1} d + 9 w_{2}/d) = w_{1} - 9 w_{2}/
d^{2} = 0$, giving $d^{*}(w) = 3 \sqrt{w_{2}/w_{1}}$. Sweeping the
weight ratio $w_{2}/w_{1}$ over its range moves $d^{*}$ across the
feasible interval and traces the front $(f_{1}, f_{2}) = (d^{*},
9/d^{*})$, a hyperbola $f_{1} f_{2} = 9$. At equal weights $d^{*} = 3$,
giving $(f_{1}, f_{2}) = (3, 3)$; tilting toward cost ($w_{2}/w_{1} =
1/9$) gives $d^{*} = 1$, the cheapest, heaviest design; tilting toward
mass ($w_{2}/w_{1} = 9/4$) gives $d^{*} = 4.5$, the lightest, dearest
design. Every point on the hyperbola is Pareto optimal, and the
weighted sum recovered the whole front because the front is convex.

\paragraph{The limit of weighted sums.}
Weighted-sum scalarisation recovers only the convex part of a Pareto
front. A front with a non-convex (re-entrant) region has points that no
weighted sum can reach: the linear contours of $\sum_{k} w_{k} f_{k}$
touch a non-convex front only at its convex hull, skipping the dimpled
interior. The repair is a different scalarisation. The
\engterm{$\epsilon$-constraint method} optimises one objective subject
to upper bounds on the others, $\min f_{1}$ subject to $f_{k} \leq
\epsilon_{k}$, and sweeping the bounds $\epsilon_{k}$ recovers the
whole front, convex or not, at the cost of solving a constrained
problem at each sweep point. The engineer who reports a Pareto front
from a weighted-sum sweep alone, on a problem whose front might be
non-convex, has under-reported the achievable designs, and the
$\epsilon$-constraint sweep is the honest alternative.

\section{Failure: local minima, ill-posed objectives, optimisation
theatre}\pathtag{core}

A solved optimisation problem is not a solved engineering problem.
The optimiser solves the problem the engineer wrote down. The
engineer's job is to write down the right problem, recognise when
the optimum found is local and not global, and notice when the
objective has been gamed by the optimiser to produce a number
that is mathematically optimal and practically harmful. The
failure modes:

\subsection{Local minima look like global minima}

A nonconvex objective can have many local minima, with the
gradient zero and the Hessian positive definite at each, and no
local diagnostic distinguishes the local minimum the iteration
found from the global minimum it missed. Gradient descent and
Newton's method converge to a local minimum from any starting
point; the result depends on the initialisation. Multiple random
restarts (run the optimiser from many different initial points
and take the best result) is the working hedge against this
failure, accepting that even the best of $K$ random restarts can
miss a deep narrow basin.

The diagnostic: inspect the loss values at the converged
iterates from many restarts. If the spread of converged losses is
narrow, the objective is plausibly close to convex, or the
restarts cluster on a flat valley. If the spread is wide, the
objective is genuinely multimodal and the engineer must report
the result with appropriate caveats. The diagnostic does not
prove the global optimum has been found; it only reports what
multiple restarts found.

\paragraph{Worked example: a one-dimensional multimodal fit.}
Fit a single Gaussian bump $g(x; a, c) = a \exp(-(x - c)^{2})$ to
data generated by two well-separated bumps centred at $c = 2$ and
$c = 8$. The least-squares objective $\sum_{i}(y_{i} - g(x_{i}; a,
c))^{2}$, viewed as a function of the centre $c$, has two basins,
one around $c = 2$ and one around $c = 8$, with a ridge between
them where the single bump matches neither cluster. Gradient
descent on $c$ initialised at $c_{0} = 1$ slides into the left
basin and reports $c^{*} \approx 2$; initialised at $c_{0} = 9$ it
slides into the right basin and reports $c^{*} \approx 8$. Both are
local minima with zero gradient and positive second derivative;
no local test distinguishes them. Ten restarts spread uniformly
over $c_{0} \in [0, 10]$ recover both basins, and the engineer
reports the better of the two residuals along with the observation
that the objective is bimodal in $c$, which is itself a finding:
the data are not described by one bump, and the right model has
two. The multistart diagnostic did not just find the optimum, it
diagnosed the model misspecification, which is the more valuable
output.

\subsection{Ill-posed objectives: the wrong objective is optimised
exactly}

Diet-problem solutions optimise stated cost subject to stated
nutritional constraints, and produce diets no one would eat.
Curve fits optimise stated residuals on stated training data, and
produce models that overfit. Trading strategies optimise stated
backtest returns on stated historical data, and produce models
that will lose money in deployment. Each of these is the
optimiser doing its job correctly on a specification that did not
quite say what the engineer meant.

The working diagnostic: state the objective in plain English
before writing it as a function. Identify what the objective does
not capture. Ask what optimum the formulation rewards if the
engineer's intent is taken literally. If the literal optimum is
not the engineering intent, rewrite the objective. Common
reformulations: add a regularisation term that penalises model
complexity; replace training-set residual with cross-validated
residual; add a constraint that bounds the worst-case loss
instead of the average loss; multi-objective formulations that
report a Pareto frontier instead of a scalar optimum.

\begin{principle}[Goodhart's law]
When a measure becomes a target, it ceases to be a good measure.
An objective installed for optimisation is, by construction, a
target; the optimisation will move the system to maximise the
target, including by exploiting any gap between the target and
the engineer's intent. The engineer's defence is to anticipate
the gap before the optimiser exploits it.
\end{principle}

The principle, originally articulated by the British economist
\textcite{text:v2ch15-goodhart1975} in the context of monetary
policy, is the most important single warning in optimisation
practice. A
factory measured on units produced will produce defective units;
a hospital measured on average wait time will discharge patients
on the wait-time clock; a model measured on training accuracy
will memorise the training set; a recommendation system measured
on click-through rate will surface clickbait. None of these are
optimiser failures; each is a specification failure caught by
the optimiser doing its job.

\subsection{Worked example: a regression that overfits the
training set}

A polynomial regression of degree $d$ on $N = 12$ data points
$(x_{i}, y_{i})$, with $y_{i} = \sin(\pi x_{i}) + \epsilon_{i}$,
$\epsilon_{i} \sim \mathcal{N}(0, 0.1^{2})$, $x_{i} \in [0, 1]$.
The objective is the training mean-squared error:
\[
\mathcal{L}_{\text{train}}(\boldsymbol{\theta}) = \frac{1}{N}
\sum_{i=1}^{N} (y_{i} - p_{d}(x_{i}; \boldsymbol{\theta}))^{2},
\]
with $p_{d}$ a polynomial of degree $d$. The
$\boldsymbol{\theta}^{*}$ that minimises $\mathcal{L}_{\text{train}}$
is a least-squares fit, computable in closed form by the normal
equations of Volume II Chapter 9.

For $d = 11$ the polynomial passes through every data point:
$\mathcal{L}_{\text{train}} = 0$. The optimiser has done its job
exactly. The fitted polynomial oscillates wildly between the
data points. The mean-squared error on a held-out test set,
evaluated at intermediate $x$ values not seen during fitting, is
two orders of magnitude larger than for $d = 3$. The training-set
optimum is the wrong objective for the engineering goal of
predicting $\sin(\pi x)$ on unseen data.

The corrected formulation adds an $\ell_{2}$ regulariser:
\[
\mathcal{L}_{\text{ridge}}(\boldsymbol{\theta}) =
\mathcal{L}_{\text{train}}(\boldsymbol{\theta}) + \lambda
\norm{\boldsymbol{\theta}}_{2}^{2},
\]
with $\lambda > 0$ the regularisation strength. The same
optimiser applied to $\mathcal{L}_{\text{ridge}}$ produces a fit
whose test error is comparable to $d = 3$ even at $d = 11$, with
$\lambda$ chosen by cross-validation (Volume II Chapter 14). The
optimiser is the same. The objective is different. The objective
makes the difference.

\subsection{Optimisation theatre}

\engterm{Optimisation theatre} is the practice of running an
optimiser on a problem that was not the right problem in order
to produce a number that looks defensible. The factory floor
where every shift's output is reported to the third decimal but
the rejection rate is not; the school district where standardised
test scores improve without learning; the climate model whose
parameters are tuned to match historical records and then
extrapolated as if the parameters were physical; the pricing
optimiser that maximises short-term revenue while destroying
customer relationships measured nowhere in the loss. The pattern
is consistent: a real engineering or organisational question is
replaced by a measurable proxy, the proxy is optimised
aggressively, the underlying question is no better answered, and
the optimisation provides cover that the question has been taken
seriously.

The engineer's defence is to ask, before running the optimiser,
three questions. What does the optimum of this objective look
like in plain English? Is that what we want? What would a
malicious optimiser exploit in the gap between the formulation
and the intent? The questions do not eliminate Goodhart effects;
they make the engineer responsible for them.

\section{What to carry forward}\pathtag{core}

The optimisation archetype reaches its formal home in this
chapter, and the working engineer carries forward a short list of
decisions that recur in every subsequent volume.

The first question on any optimisation problem is the structural
class, and the line that matters is convex versus nonconvex. A
convex problem has a global theory: every local optimum is global,
the KKT conditions characterise the optimum, duality supplies a
certificate, and a catalogue of algorithms converges with stated
rates. A nonconvex problem forfeits all of this; the engineer gets
a local optimum, multiple restarts as a hedge, and a duty to
report the guarantee actually earned. Recognising convexity from
the convex-function library, without computing a Hessian, is the
single most valuable diagnostic skill the chapter teaches.

The second decision is method selection, and it is an instance of
the scaling archetype: multiply the per-iteration cost by the
expected iteration count and check the product against the compute
budget. Small smooth problems take Newton or interior-point;
medium ones take L-BFGS; large or stochastic ones take gradient
methods with momentum or Adam; constrained ones take a constrained
solver through a modelling layer. The condition number $\kappa$
sets the first-order convergence rate, and the engineer who knows
$\kappa$ knows in advance whether plain gradient descent will
finish in this decade.

The third habit is duality. Every constrained optimum comes with
multipliers, and the multipliers are the derivatives of the
optimal value with respect to the constraints. They say which
constraint to spend the next dollar relaxing, they certify
optimality through the duality gap, and they convert a raw number
into an engineering recommendation. An optimum reported without its
multipliers is half an answer.

The fourth, and the one the failure section exists to install, is
that a solved optimisation problem is not a solved engineering
problem. The optimiser solves the problem written down. Goodhart's
law guarantees that any gap between the stated objective and the
real intent will be found and exploited by a sufficiently capable
optimiser. The engineer's responsibility is to anticipate that gap
before the optimiser does, by stating the objective in plain
English and asking what its literal optimum rewards. This habit
carries into Volume VII's machine learning, Volume IX's control,
and every design decision that hides an optimisation inside it.

\section{Project}\pathtag{core}

\begin{project}[Gradient descent from scratch with at least one
acceleration scheme]

\textbf{Track.} Simulation with analysis. \textbf{Hazard class.}
None. \textbf{Tools.} Paper, pencil, programming environment of
the reader's choice (Python, MATLAB, Julia, or equivalent).

\textbf{The problem.} Implement gradient descent from scratch on
a chosen objective, plus at least one of momentum, Nesterov
accelerated gradient, RMSProp, or Adam. Compare the
implementations against each other and against a library
optimiser, and report on convergence rate, learning-rate
sensitivity, and qualitative behaviour on the chosen objective.

\textbf{Suggested objectives.}
\begin{enumerate}[leftmargin=*]
  \item The polynomial-fit problem from the Volume II Chapter 9
    project: minimise $\norm{A \vect{x} - \vect{b}}_{2}^{2}$ for
    a polynomial design matrix $A$ on a noisy dataset of the
    reader's choice, possibly with an $\ell_{2}$ regularisation
    term. Closed-form least-squares is available as the
    reference; the reader's iterative methods should converge to
    the same answer.
  \item The Rosenbrock function $f(x, y) = (1 - x)^{2} + 100
    (y - x^{2})^{2}$. Nonconvex, with a famously elongated
    valley. The classical hard test for first-order methods.
  \item A small neural-network loss: a two-layer fully connected
    network on MNIST, fashion-MNIST, or any small classification
    dataset. Nonconvex, large $n$, the regime that justifies
    Adam.
  \item A logistic regression on a real binary classification
    dataset of the reader's choice. Convex, a fair comparison
    target where the global optimum is known.
\end{enumerate}

\textbf{The deliverable.} A 6-to-10-page report and a 1500-word
reflection. The report must contain
\begin{enumerate}[leftmargin=*]
  \item A statement of the chosen objective $f$, including the
    parameter dimension $n$, the data, and the closed-form (or
    library-computed) optimum where available.
  \item A from-scratch implementation of vanilla gradient
    descent. From-scratch means the reader writes the iteration
    explicitly; library calls for matrix operations are allowed,
    but \verb|library.gradient_descent(...)| is not. The gradient
    is computed analytically from the objective formula or by
    automatic differentiation, with the choice documented.
  \item A from-scratch implementation of at least one
    acceleration: momentum, Nesterov, RMSProp, or Adam. The
    update rule, the state maintained, and the hyperparameter
    defaults must be stated explicitly.
  \item A learning-rate sensitivity sweep: at least three step
    sizes spanning two orders of magnitude (e.g.\ $\eta \in
    \set{10^{-1}, 10^{-2}, 10^{-3}}$). The reader plots the loss
    versus iteration count for each, and reports which step size
    gave the fastest convergence and which step size diverged.
  \item A comparison against a library optimiser (e.g.\
    \repopath{scipy.optimize.minimize}, PyTorch's optimisers,
    Julia's Optim.jl). The reader reports the iteration count,
    wall-clock time, and final loss for each method.
  \item Plotted convergence trajectories. For 2D objectives
    (Rosenbrock), plot the trajectory on a contour plot of $f$.
    For higher-dimensional objectives, plot the loss versus
    iteration on log scale.
\end{enumerate}

\textbf{The reflection.} The 1500-word reflection addresses
\begin{itemize}[leftmargin=*]
  \item Local minima encountered (or absent). Did the reader's
    method converge to the same point from different
    initialisations? Did multiple restarts help? On a convex
    objective, the answer is yes by Section 15.2's theorem; on
    a nonconvex objective, the answer is informative.
  \item Learning-rate sensitivity. How wide was the working
    range of step sizes? At what step size did the iteration
    diverge? Cite the smoothness constant $L$ if computable, or
    estimate it.
  \item When stochastic methods help. If the reader chose a
    large-data objective, compare full-batch gradient descent
    against minibatch SGD on iteration count and on wall-clock
    time. Identify the regime in which SGD is faster and the
    regime in which it is not.
  \item How the engineer should choose between methods given a
    real-world objective. State a working decision rule
    referencing the principle in Section 15.5.
  \item One source of error the reader noticed in their own
    working, and one heuristic the reader has acquired by
    completing the project.
\end{itemize}

\textbf{Time.} Eight to twelve hours of implementation, two to
three hours of experiments, three to four hours of writing.
Total approximately fifteen to twenty hours, distributed across
three or four sessions.

\textbf{Solutions.} A reference Python implementation of vanilla
gradient descent and Adam on the polynomial-fit and Rosenbrock
objectives, available from the general editor.

\end{project}

\section{Exercises}

The exercises train calculation, derivation, estimation,
simulation, design, diagnosis, failure analysis, and judgment.
Calculation and derivation problems have full solutions;
open-ended problems have rubrics. Exercises are tagged with their
reader-path tier.

\subsection*{Calculation}\pathtag{core}

\begin{exercise}
Solve the LP $\min_{\vect{x}} \vect{c}^{\transpose} \vect{x}$
subject to $\vect{x} \geq \vect{0}$ and $\vect{1}^{\transpose}
\vect{x} = 1$, with $\vect{c} = (3, 1, 4, 1, 5)$. Identify the
optimum and the active constraints, and write down the dual
variable.
\end{exercise}

\paragraph{Solution sketch.}
The feasible set is the probability simplex in $\R^{5}$. A linear
objective on a simplex is minimised at a vertex; the vertices are
the unit-coordinate points $\vect{e}_{i}$, with objective value
$c_{i}$. The smallest $c_{i}$ is $1$, achieved at $i = 2$ and $i =
4$. Any convex combination of $\vect{e}_{2}$ and $\vect{e}_{4}$
also attains $p^{*} = 1$; the optimum is the line segment
$\set{(0, t, 0, 1 - t, 0) : t \in [0, 1]}$. Active constraints at
any optimum: the equality $\vect{1}^{\transpose} \vect{x} = 1$ and
the nonnegativity bounds on the three components fixed at zero.
The dual variable on the equality is $\nu^{*} = 1$ (the optimal
cost per unit relaxation of the simplex sum), and the dual
variables on the active $x_{i} \geq 0$ constraints are $c_{i} -
\nu^{*}$, vanishing for $i \in \set{2, 4}$ as required by
complementary slackness.

\begin{exercise}
Compute the optimum of the QP $\min_{\vect{x}} \tfrac{1}{2}
\norm{\vect{x}}_{2}^{2}$ subject to $\vect{1}^{\transpose}
\vect{x} = 1$ in $\R^{n}$ for general $n$. Report
$\vect{x}^{*}$, $p^{*}$, and the multiplier $\nu^{*}$.
\end{exercise}

\paragraph{Solution sketch.}
Lagrangian $\mathcal{L} = \tfrac{1}{2} \vect{x}^{\transpose}
\vect{x} + \nu (\vect{1}^{\transpose} \vect{x} - 1)$. Stationarity:
$\vect{x}^{*} = -\nu^{*} \vect{1}$. Primal feasibility: $-\nu^{*}
n = 1$, so $\nu^{*} = -1/n$ and $\vect{x}^{*} = (1/n) \vect{1}$,
the uniform distribution. Optimal value $p^{*} = \tfrac{1}{2}
\norm{\vect{x}^{*}}^{2} = \tfrac{1}{2} n (1/n)^{2} = 1/(2n)$. The
multiplier $\nu^{*}$ is the rate at which $p^{*}$ changes when the
simplex level is shifted: $\dv{p^{*}}{b} = -1/n$ at $b = 1$,
matching $\nu^{*}$.

\begin{exercise}
For $f(x, y) = x^{2} + 4 y^{2} - 2 x + 8 y + 5$, compute the
gradient and the Hessian, identify the unique critical point, and
verify that it is a global minimum by the Hessian test. Report
$f(\vect{x}^{*})$.
\end{exercise}

\paragraph{Solution sketch.}
$\grad f = (2 x - 2, 8 y + 8)$, vanishing at $(x^{*}, y^{*}) =
(1, -1)$. Hessian $\nabla^{2} f = \diag(2, 8)$, with eigenvalues
$2$ and $8$, both positive, so $f$ is strictly convex on $\R^{2}$
and the critical point is the global minimum. $f(1, -1) = 1 + 4 -
2 - 8 + 5 = 0$.

\begin{exercise}
Run two iterations of gradient descent by hand on $f(x) = x^{2}$
from $x_{0} = 4$ with constant step size $\eta = 0.1$. Report
$x_{1}, x_{2}, x_{3}$ and verify the geometric convergence rate.
\end{exercise}

\paragraph{Solution sketch.}
$f'(x) = 2x$, so $x_{k+1} = x_{k} - 0.2 x_{k} = 0.8 x_{k}$. From
$x_{0} = 4$: $x_{1} = 3.2$, $x_{2} = 2.56$, $x_{3} = 2.048$. The
contraction factor per step is $0.8$, matching the theoretical
$1 - \eta \cdot 2 = 0.8$ for this quadratic. Geometric convergence
with rate $0.8$ per step; reaching $|x_{k}| < 10^{-6}$ takes
$\log(4/10^{-6}) / \log(1/0.8) \approx 15.2 / 0.223 \approx 68$
iterations.

\begin{exercise}
Run one Newton step by hand on $f(x, y) = x^{2} + 2 x y + 3 y^{2}
- 2 x - 4 y$ from $(x_{0}, y_{0}) = (0, 0)$. Verify that the
result is the global minimum.
\end{exercise}

\paragraph{Solution sketch.}
$\grad f = (2 x + 2 y - 2, \; 2 x + 6 y - 4)$, evaluating to
$(-2, -4)$ at $(0, 0)$. Hessian $H = \begin{pmatrix} 2 & 2 \\ 2 & 6
\end{pmatrix}$, with $\det H = 8$ and inverse $H^{-1} = \tfrac{1}{8}
\begin{pmatrix} 6 & -2 \\ -2 & 2 \end{pmatrix}$. Newton step:
$-H^{-1} \grad f = -\tfrac{1}{8} (6 \cdot (-2) + (-2)(-4), \;
(-2)(-2) + 2 \cdot (-4)) = -\tfrac{1}{8} (-4, -4) = (1/2, 1/2)$.
The new point is $(0, 0) + (1/2, 1/2) = (1/2, 1/2)$. Setting $\grad
f = \vect{0}$ directly: $x + y = 1$ and $x + 3 y = 2$, giving $y =
1/2, x = 1/2$. The Newton step reached the unique critical point
in one move, as expected for a strictly convex quadratic.

\begin{exercise}
Solve the KKT conditions for $\min_{\vect{x}} \tfrac{1}{2}
\norm{\vect{x}}_{2}^{2}$ subject to $\vect{a}^{\transpose}
\vect{x} \geq b$ with $\vect{a} = (1, 2, 2)$ and $b = 6$. Report
$\vect{x}^{*}$, $p^{*}$, and $\lambda^{*}$, and verify
complementary slackness.
\end{exercise}

\paragraph{Solution sketch.}
Rewrite as $f_{1}(\vect{x}) = b - \vect{a}^{\transpose} \vect{x}
\leq 0$. Stationarity: $\vect{x}^{*} = \lambda^{*} \vect{a}$. If
$\lambda^{*} = 0$ then $\vect{x}^{*} = \vect{0}$, which violates
$\vect{a}^{\transpose} \vect{x} \geq b$ for $b > 0$; so
$\lambda^{*} > 0$ and the constraint is active:
$\vect{a}^{\transpose} \vect{x}^{*} = b$, giving $\lambda^{*} =
b / \norm{\vect{a}}^{2} = 6 / 9 = 2/3$. Then $\vect{x}^{*} =
(2/3)(1, 2, 2) = (2/3, 4/3, 4/3)$ and $p^{*} = \tfrac{1}{2}
\norm{\vect{x}^{*}}^{2} = \tfrac{1}{2} \cdot (4/9) \cdot 9 = 2$.
Complementary slackness $\lambda^{*} (b - \vect{a}^{\transpose}
\vect{x}^{*}) = (2/3) \cdot 0 = 0$ holds because the constraint is
active. Shadow-price check: $\dv{p^{*}}{b} = b / \norm{\vect{a}}^{2}
= 2/3 = \lambda^{*}$ at $b = 6$.

\subsection*{Derivation}\pathtag{standard}

\begin{exercise}
Derive the dual of the LP $\min_{\vect{x}} \vect{c}^{\transpose}
\vect{x}$ subject to $A \vect{x} \leq \vect{b}$, $\vect{x} \geq
\vect{0}$ from the Lagrangian. Identify the dual variables and
their sign constraints.
\end{exercise}

\paragraph{Solution sketch.}
Two inequality blocks: $A \vect{x} - \vect{b} \leq \vect{0}$ with
multiplier $\boldsymbol{\lambda} \geq \vect{0}$, and
$-\vect{x} \leq \vect{0}$ with multiplier $\boldsymbol{\mu} \geq
\vect{0}$. Lagrangian
$\mathcal{L} = \vect{c}^{\transpose} \vect{x} +
\boldsymbol{\lambda}^{\transpose}(A \vect{x} - \vect{b})
- \boldsymbol{\mu}^{\transpose} \vect{x}
= (\vect{c} + A^{\transpose} \boldsymbol{\lambda}
- \boldsymbol{\mu})^{\transpose} \vect{x}
- \boldsymbol{\lambda}^{\transpose} \vect{b}$. The infimum over
$\vect{x}$ is $-\infty$ unless $\vect{c} + A^{\transpose}
\boldsymbol{\lambda} - \boldsymbol{\mu} = \vect{0}$, in which case
the value is $-\boldsymbol{\lambda}^{\transpose} \vect{b}$.
Eliminate $\boldsymbol{\mu} = \vect{c} + A^{\transpose}
\boldsymbol{\lambda} \geq \vect{0}$. The dual is
$\max_{\boldsymbol{\lambda} \geq \vect{0}} -\vect{b}^{\transpose}
\boldsymbol{\lambda}$ subject to $A^{\transpose} \boldsymbol{\lambda}
\geq -\vect{c}$, equivalently (substituting $\boldsymbol{\lambda}
\to -\boldsymbol{\lambda}$, conventional sign): the dual carries one
nonpositive multiplier per primal inequality and one nonnegative
multiplier per primal nonnegativity, and the dual feasibility set
$A^{\transpose} \boldsymbol{\lambda} \leq \vect{c}$ is the
constraint side of the original problem.

\begin{exercise}
Derive the dual of a general convex QP $\min_{\vect{x}}
\tfrac{1}{2} \vect{x}^{\transpose} P \vect{x} +
\vect{q}^{\transpose} \vect{x}$ subject to $A \vect{x} =
\vect{b}$. Solve the inner minimisation analytically using the
positive definiteness of $P$.
\end{exercise}

\paragraph{Solution sketch.}
Lagrangian $\mathcal{L} = \tfrac{1}{2} \vect{x}^{\transpose} P
\vect{x} + \vect{q}^{\transpose} \vect{x} +
\boldsymbol{\nu}^{\transpose}(A \vect{x} - \vect{b})$. Stationarity
in $\vect{x}$: $P \vect{x} + \vect{q} + A^{\transpose}
\boldsymbol{\nu} = \vect{0}$, so $\vect{x}^{*}(\boldsymbol{\nu}) =
-P^{-1}(\vect{q} + A^{\transpose} \boldsymbol{\nu})$ (using $P
\succ 0$). Substituting back,
$g(\boldsymbol{\nu}) = -\tfrac{1}{2}(\vect{q} + A^{\transpose}
\boldsymbol{\nu})^{\transpose} P^{-1} (\vect{q} + A^{\transpose}
\boldsymbol{\nu}) - \boldsymbol{\nu}^{\transpose} \vect{b}$. The
dual problem is $\max_{\boldsymbol{\nu}} g(\boldsymbol{\nu})$
(unconstrained in $\boldsymbol{\nu}$), a concave quadratic in
$\boldsymbol{\nu}$. Setting $\nabla g = \vect{0}$:
$A P^{-1}(\vect{q} + A^{\transpose} \boldsymbol{\nu}^{*}) =
-\vect{b}$, giving $\boldsymbol{\nu}^{*} =
-(A P^{-1} A^{\transpose})^{-1}(\vect{b} + A P^{-1} \vect{q})$,
which is the Schur-complement formula for the equality-constrained
QP.

\begin{exercise}
Derive the gradient-descent convergence rate for a quadratic
objective $f(\vect{x}) = \tfrac{1}{2} \vect{x}^{\transpose} P
\vect{x}$ with $P$ positive definite. Identify the optimal
constant step size in terms of the eigenvalues of $P$.
\end{exercise}

\paragraph{Solution sketch.}
$\grad f = P \vect{x}$, so the iteration is $\vect{x}_{k+1} =
(I - \eta P) \vect{x}_{k}$, hence $\vect{x}_{k} = (I - \eta P)^{k}
\vect{x}_{0}$. The spectral radius of $I - \eta P$ governs the
geometric rate. Let $\mu = \lambda_{\min}(P)$ and $L =
\lambda_{\max}(P)$; the eigenvalues of $I - \eta P$ are $1 - \eta
\lambda_{i}$. The worst-case contraction is $\max(|1 - \eta \mu|,
|1 - \eta L|)$, minimised when $1 - \eta \mu = -(1 - \eta L)$,
giving $\eta^{*} = 2 / (\mu + L)$ and worst-case rate $(L - \mu)/
(L + \mu) = (\kappa - 1)/(\kappa + 1)$. The familiar $1 - \mu/L$
bound is the loose version obtained from $\eta = 1/L$; the tight
bound uses the symmetric choice.

\begin{exercise}
Derive the Newton step for $f(\vect{x}) = \tfrac{1}{2}
\vect{x}^{\transpose} P \vect{x} + \vect{q}^{\transpose}
\vect{x} + r$ and verify that one Newton step from any
$\vect{x}_{0}$ reaches the optimum.
\end{exercise}

\paragraph{Solution sketch.}
$\grad f = P \vect{x} + \vect{q}$, Hessian $H = P$. Newton step
$\vect{p} = -H^{-1} \grad f(\vect{x}_{0}) = -P^{-1}(P \vect{x}_{0}
+ \vect{q}) = -\vect{x}_{0} - P^{-1} \vect{q}$. Updated point
$\vect{x}_{0} + \vect{p} = -P^{-1} \vect{q}$. Setting $\grad f =
\vect{0}$ directly gives $\vect{x}^{*} = -P^{-1} \vect{q}$, the
same point. One Newton step reaches the global optimum from any
$\vect{x}_{0}$ when $f$ is a strictly convex quadratic.

\begin{exercise}
Show that the pointwise maximum of a finite collection of convex
functions is convex. Use the definition of convexity directly.
\end{exercise}

\paragraph{Solution sketch.}
Let $f(\vect{x}) = \max_{i \in \set{1, \ldots, m}} f_{i}(\vect{x})$
with each $f_{i}$ convex on a common convex domain. For
$\vect{x}, \vect{y}$ in the domain and $\theta \in [0, 1]$, and any
$i$,
$f_{i}(\theta \vect{x} + (1 - \theta) \vect{y}) \leq
\theta f_{i}(\vect{x}) + (1 - \theta) f_{i}(\vect{y})
\leq \theta f(\vect{x}) + (1 - \theta) f(\vect{y})$. Taking the
maximum over $i$ on the left gives $f(\theta \vect{x} + (1 -
\theta) \vect{y}) \leq \theta f(\vect{x}) + (1 - \theta)
f(\vect{y})$, which is convexity of $f$.

\begin{exercise}
Show that the composition $g(A \vect{x} + \vect{b})$ of a convex
$g : \R^{m} \to \R$ with an affine map $\vect{x} \mapsto A \vect{x}
+ \vect{b}$, $A \in \R^{m \times n}$, is convex on $\R^{n}$.
\end{exercise}

\paragraph{Solution sketch.}
Let $h(\vect{x}) = g(A \vect{x} + \vect{b})$. For
$\vect{x}, \vect{y} \in \R^{n}$ and $\theta \in [0, 1]$, $A(\theta
\vect{x} + (1 - \theta) \vect{y}) + \vect{b} = \theta (A \vect{x}
+ \vect{b}) + (1 - \theta)(A \vect{y} + \vect{b})$ by linearity of
$A$. Applying convexity of $g$ to the right-hand side,
$h(\theta \vect{x} + (1 - \theta) \vect{y}) = g(\theta (A \vect{x}
+ \vect{b}) + (1 - \theta)(A \vect{y} + \vect{b})) \leq \theta g(A
\vect{x} + \vect{b}) + (1 - \theta) g(A \vect{y} + \vect{b}) =
\theta h(\vect{x}) + (1 - \theta) h(\vect{y})$, the defining
inequality for convexity of $h$.

\subsection*{Estimation}\pathtag{core}

\begin{exercise}
Estimate the optimal pipe diameter $d^{*}$ for the cost-of-pipe
problem in Section 15.1.2 with $c_{1} = 1000\,\text{USD}/\si{\meter\cubed}$
(material), $c_{2} = 5 \times 10^{8}\,\text{USD}\cdot\si{\second\cubed\per\meter\tothe{8}}$
(pumping). Use $L = 100\,\si{\meter}$, $Q = 0.05\,\si{\cubic\meter\per\second}$.
Verify $d^{*} \in [0.05, 0.30]\,\si{\meter}$ and report the
total optimal cost.
\end{exercise}

\paragraph{Solution sketch.}
The first-order condition gives $d^{*} = (5 c_{2} Q^{3} / (2
c_{1}))^{1/7} = (156)^{1/7} \approx 2.07\,\si{\meter}$, which lies
outside the feasible interval $[0.05, 0.30]$. The constrained
optimum is the upper endpoint $d = 0.30\,\si{\meter}$. Substituting
into $f_{0}$: material cost $c_{1} L d^{2} = 1000 \cdot 100 \cdot
0.09 = 9000$\,USD; pumping cost $c_{2} L Q^{3} d^{-5} = 5 \times
10^{8} \cdot 100 \cdot 1.25 \times 10^{-4} / 0.00243 \approx 2.57
\times 10^{9}$\,USD. The pumping cost dominates the bound, which
is the formal warning that the toy coefficients are not consistent
with a real engineering economy at this scale. The exercise serves
its purpose: it demonstrates that an out-of-range first-order
optimum forces a boundary solution and that the numerical answer
is dictated by the cost-coefficient ratio, which the engineer must
calibrate against actual market and pump-curve data before
reporting any number.

\begin{exercise}
Estimate the condition number $\kappa = L/\mu$ of the
least-squares objective $f(\vect{x}) = \tfrac{1}{2} \norm{A
\vect{x} - \vect{b}}_{2}^{2}$ when $A \in \R^{1000 \times 100}$
has singular values logarithmically distributed between $\sigma_{1}
= 100$ and $\sigma_{100} = 0.01$. Estimate the iteration count
gradient descent needs to reach error $\epsilon = 10^{-6}$ from a
unit-norm initial error.
\end{exercise}

\paragraph{Solution sketch.}
The Hessian is $A^{\transpose} A$, whose eigenvalues are
$\sigma_{i}^{2}$. With $\sigma_{\max} = 100$ and $\sigma_{\min} =
0.01$, $L = 10^{4}$, $\mu = 10^{-4}$, and $\kappa = L/\mu =
10^{8}$. The geometric-rate iteration count to reach error
$\epsilon$ scales as $\kappa \log(1/\epsilon) = 10^{8} \cdot \log
10^{6} \approx 1.4 \times 10^{9}$ iterations. The estimate is a
warning: a problem this ill-conditioned is not solvable by plain
gradient descent in any reasonable time, and the working response
is preconditioning (rescaling the variables so the Hessian
eigenvalues cluster), a second-order method (Newton or L-BFGS
which absorbs the conditioning), or reformulation (regularisation
that floors the small eigenvalues).

\begin{exercise}
Estimate the iteration count for SGD with step schedule $\eta_{k}
= \eta_{0} / \sqrt{k+1}$ to reach expected suboptimality
$10^{-3}$ on a smooth convex objective with per-sample gradient
variance $\sigma^{2} = 1$ and Lipschitz constant $L = 10$.
Compare to the iteration count of full gradient descent under
the same accuracy target.
\end{exercise}

\paragraph{Solution sketch.}
The SGD convergence bound is $\E[f] - f^{*} \leq O(L D^{2}/\sqrt{K}
+ \sigma D / \sqrt{K})$ for diameter $D$ of the search domain.
Solving for $10^{-3}$ at $D \sim 1$: $K \approx (L + \sigma)^{2} /
\epsilon^{2} \approx 121 / 10^{-6} \approx 10^{8}$ iterations. Full
gradient descent on the same objective has geometric rate, so it
takes $\log(1/\epsilon)/\log(1/(1 - 1/\kappa))$ iterations; with
the same $L$ and an unknown $\mu$, even pessimistic $\kappa = 100$
gives $\approx 700$ iterations. The trade is per-iteration cost:
each full-gradient step costs $N$ per-sample gradient
computations, so on a dataset of $N > 10^{5}$ the SGD count
becomes competitive in wall-clock time.

\begin{exercise}
Estimate the wall-clock cost of one full-gradient evaluation on
a dataset of $N = 10^{6}$ samples, assuming each per-sample
gradient takes $10\,\si{\micro\second}$. Compare to the cost of a
minibatch gradient with $b = 256$. Report the ratio.
\end{exercise}

\paragraph{Solution sketch.}
Full gradient: $10^{6} \cdot 10\,\si{\micro\second} = 10\,\si{\second}$
per evaluation. Minibatch ($b = 256$): $256 \cdot 10\,\si{\micro\second}
\approx 2.6\,\si{\milli\second}$ per minibatch step, a factor
$\approx 3900$ faster per iteration. For $K_{\text{SGD}}$ SGD
iterations to make equivalent progress to one full-gradient step,
SGD wins on wall-clock time whenever $K_{\text{SGD}} < 3900$. The
SGD convergence theorem says one full-gradient step's progress is
matched after $O(\kappa)$ to $O(\kappa^{2})$ SGD iterations
depending on the regime; for any well-conditioned problem (small
$\kappa$) and any large dataset (large $N$), SGD dominates on
wall-clock time per unit of progress.

\begin{exercise}
Estimate the size of the branch-and-bound tree in the worst case
for an integer program with $n = 30$ binary variables, assuming
no pruning. Compare to the realistic case in which an LP-relaxation
gap of $5\%$ allows pruning $99\%$ of the tree.
\end{exercise}

\paragraph{Solution sketch.}
Worst case: complete binary tree of depth $30$, so $2^{30} \approx
10^{9}$ leaves. At $10^{6}$ LP-relaxation solves per second
(a representative figure for a modern solver on a small LP,
current as of 2026), the unpruned search takes $\approx 10^{3}\,
\si{\second} \approx 17\,\si{\minute}$. With $99\%$ pruning the
effective tree has $\approx 10^{7}$ nodes, finishing in seconds. The
estimate illustrates the practical message of Section 15.7: the
runtime of a MIP solver lives or dies on the LP-relaxation
quality, and the engineer's effort is well spent on formulation
strength (tight constraints, cutting planes, valid inequalities)
before paying for more compute.

\subsection*{Simulation}\pathtag{mastery}

\begin{exercise}
Implement gradient descent on $f(\vect{x}) = \tfrac{1}{2}
\vect{x}^{\transpose} P \vect{x}$ in $\R^{2}$ with $P =
\diag(1, 100)$ from $\vect{x}_{0} = (1, 1)$. Plot the trajectory
on a contour plot for step sizes $\eta \in \set{10^{-3}, 10^{-2},
0.02}$. Identify the step size at which the iteration diverges.
\end{exercise}

\paragraph{Solution sketch.}
The stability threshold is $\eta L < 2$, i.e.\ $\eta < 0.02$ here.
At $\eta = 10^{-3}$ the iteration converges smoothly but slowly
(rate $\approx 1 - 10^{-3}$ in the $x_{1}$ direction); at $\eta =
10^{-2}$ the $x_{2}$ direction oscillates but contracts (rate $1 -
1 = 0$ for $x_{1}$, rate $|1 - 1| = 0$ then dominant rate from
$x_{2}$ at $|1 - 100 \cdot 10^{-2}| = 0$); at $\eta = 0.02$
exactly, the iteration is marginally stable (rate $|1 - 100 \cdot
0.02| = 1$ in $x_{2}$). At $\eta$ slightly above $0.02$ the
iteration diverges. The reference implementation
\repopath{volumes/02-form/ch15-optimisation/code/gradient_descent.py}
sweeps these values and prints the convergence status.

\begin{exercise}
Implement Newton's method on $f(x, y) = (x - 1)^{2} + 100
(y - x^{2})^{2}$ (Rosenbrock with $a = 1, b = 100$) from
$(x_{0}, y_{0}) = (-1, 1)$. Compare the iteration count against
gradient descent with the best step size you find in $\set{10^{-2},
10^{-3}, 10^{-4}}$.
\end{exercise}

\paragraph{Solution sketch.}
Newton's method, with a line search to guard against the Hessian
becoming indefinite in the curved valley, converges in $\sim 20$
to $\sim 50$ iterations from $(-1, 1)$. Vanilla gradient descent
at the best step size in the given set ($\eta = 10^{-3}$ is
typical, larger values diverge) takes $> 10^{4}$ iterations to
reach $f < 10^{-4}$, sometimes failing to reach that threshold
inside the iteration cap. The ratio is two to three orders of
magnitude in iterations, paid back by a Hessian solve at each
Newton step. The reference comparison is in
\repopath{volumes/02-form/ch15-optimisation/code/rosenbrock_methods.py}.

\begin{exercise}
Implement vanilla SGD and SGD with momentum (decay $\beta = 0.9$)
on a logistic regression objective for a binary classification
dataset of your choice (synthesised or real). Run for 20 epochs
with minibatch size $b = 32$ and report training-loss curves for
both methods.
\end{exercise}

\paragraph{Solution sketch.}
Logistic loss per sample $\ell(\boldsymbol{\theta}; \vect{x}_{i},
y_{i}) = \log(1 + \exp(-y_{i} \boldsymbol{\theta}^{\transpose}
\vect{x}_{i}))$, gradient $-y_{i} \vect{x}_{i} \sigma(-y_{i}
\boldsymbol{\theta}^{\transpose} \vect{x}_{i})$ where $\sigma$ is
the logistic. With a well-scaled dataset, vanilla SGD at $\eta =
0.1$ converges to a smooth loss curve; momentum at $\beta = 0.9$
typically reaches the same loss in roughly half the epoch count
and shows visibly smoother loss curves. Both should plateau within
20 epochs on a small dataset; if either oscillates, reduce
$\eta$ by a factor of $3$.

\begin{exercise}
Solve the LP $\min_{\vect{x}} -\vect{c}^{\transpose} \vect{x}$
subject to $A \vect{x} \leq \vect{b}, \vect{x} \geq \vect{0}$ for
a randomly generated $A \in \R^{50 \times 100}$, $\vect{b} \in
\R^{50}_{+}$, $\vect{c} \in \R^{100}$. Use a library LP solver
(scipy.optimize.linprog, CVXPY+HiGHS, or equivalent). Report the
optimal value, the active constraints, and the dual variables.
\end{exercise}

\paragraph{Solution sketch.}
With $50$ constraints and $100$ variables in a generic random
instance, the optimal vertex has roughly $50$ tight constraints
(the binding rows of $A \vect{x} \leq \vect{b}$ plus some active
nonnegativity bounds). The solver returns the optimal $\vect{x}^{*}$,
the objective value, and the dual vector $\boldsymbol{\lambda}^{*}
\in \R^{50}_{\geq 0}$. Inspect $\boldsymbol{\lambda}^{*}$: the
nonzero components index the binding constraints, by complementary
slackness. The template
\repopath{volumes/02-form/ch15-optimisation/code/lp_diet.py}
adapts directly to the random version by replacing the data
loading with \verb|numpy.random.default_rng(seed)| and generating
$A, \vect{b}, \vect{c}$ with the right signs.

\begin{exercise}
Compare gradient descent, momentum, and Adam on the Rosenbrock
function from $(x_{0}, y_{0}) = (-1, 1)$. For each method, sweep
the step size over five values and report the iteration count to
reach $f < 10^{-4}$. Identify the most robust method on this
nonconvex problem.
\end{exercise}

\paragraph{Solution sketch.}
Adam is the most robust to step-size choice on this benchmark:
working step sizes span roughly an order of magnitude with similar
convergence counts. Momentum at $\beta = 0.9$ accelerates over
plain gradient descent by roughly $4\times$ on the working step
size. Plain gradient descent has a narrow working range: at $\eta$
below $10^{-3}$ progress is glacial, above $\eta = 3 \times 10^{-3}$
the curved valley triggers oscillation. The reference run in
\repopath{volumes/02-form/ch15-optimisation/code/rosenbrock_methods.py}
reports the iteration counts in
\repopath{volumes/02-form/ch15-optimisation/data/rosenbrock_convergence.csv}.

\subsection*{Design}\pathtag{standard}

\begin{exercise}
Design a portfolio optimisation problem (Markowitz) for $n = 10$
hypothetical assets with given expected returns $\boldsymbol{\mu}$
and a given positive-definite covariance matrix $\Sigma$ (you
choose the values). Solve as a QP for three risk-aversion values
$\gamma \in \set{0.1, 1, 10}$ and plot the resulting portfolios on
a return-vs-risk diagram (the efficient frontier).
\end{exercise}

\paragraph{Discussion.}
Generate $\boldsymbol{\mu}$ as ten plausible annual returns
($1\%$ to $15\%$) and build $\Sigma$ as $A A^{\transpose} +
\sigma^{2} I$ for a random $A$ so that $\Sigma$ is strictly
positive definite. At $\gamma = 0.1$ the optimiser tolerates risk
and concentrates the portfolio in the highest-return assets; at
$\gamma = 10$ the optimiser is risk-averse and spreads weight over
many low-correlation assets to minimise variance. The trajectory
from $\gamma = 0.1$ down to $\gamma = 10$ traces the upper edge of
the efficient frontier. Discuss: the frontier is convex; every
point on it is Pareto-optimal in $(\text{risk}, \text{return})$;
no point above the frontier is achievable by any
$\boldsymbol{w}$ on the simplex.

\begin{exercise}
Design a production-planning LP for a hypothetical factory
producing $n = 5$ products from $m = 3$ raw materials, with
specified per-unit profits and per-unit resource consumption.
Solve and identify the binding resource constraints. Compute the
shadow prices and interpret them as the maximum the engineer
should pay for one additional unit of each binding resource.
\end{exercise}

\paragraph{Discussion.}
At a generic optimum, between zero and three of the resource
constraints are binding (active) and the rest are slack. The
shadow prices $\nu_{i}^{*}$ on binding constraints are the per-unit
profit gained by relaxing the corresponding resource by one unit;
the slack constraints have $\nu_{i}^{*} = 0$ by complementary
slackness. The engineering interpretation: the engineer should pay
up to $\nu_{i}^{*}$ for each additional unit of a binding
resource. Above that price the additional unit is not worth its
cost; below it the engineer is leaving money on the table. The
shadow-price reading is the load-bearing reason LP duality is
useful in practice, not just as theory.

\begin{exercise}
Design a regularised regression problem on a dataset of your
choice (Volume I project data, a Volume II Chapter 9 dataset, or
a public dataset). Sweep the regularisation strength $\lambda$
over five values and plot the training and cross-validated
test error against $\lambda$. Identify the $\lambda$ that
minimises the test error.
\end{exercise}

\paragraph{Discussion.}
The training-error curve is monotonically nondecreasing in
$\lambda$: more regularisation, worse fit on the training set. The
test-error curve is U-shaped: too little regularisation overfits,
too much underfits. The minimum test error sits at a
problem-specific $\lambda^{*}$ that depends on signal-to-noise
ratio and sample size. The standard sweep is logarithmic across
five orders of magnitude (e.g.\ $\lambda \in \set{10^{-4}, 10^{-3},
\ldots, 10^{0}}$); the working $\lambda^{*}$ is the minimum of the
test-error curve, with cross-validation supplying the confidence
interval. Discuss how the $\lambda^{*}$ would shift if the sample
size doubled (smaller $\lambda^{*}$ as noise averages out) or if
the noise variance doubled (larger $\lambda^{*}$ to control the
overfit).

\begin{exercise}
Design a model-predictive-control QP for a one-dimensional
double-integrator system $x_{k+1} = x_{k} + v_{k}$, $v_{k+1} =
v_{k} + u_{k}$ with horizon $H = 10$, target state
$(x^{*}, v^{*}) = (1, 0)$, and actuator limit $|u_{k}| \leq 0.1$.
Solve once for the planning sequence $\set{u_{k}}_{k=0}^{H-1}$
and apply only the first input. Iterate to convergence.
\end{exercise}

\paragraph{Discussion.}
The per-step QP has $10$ decision variables (the $u_{k}$) with box
constraints. The objective penalises the state deviation from the
target plus a small actuator-effort term. The receding-horizon
trajectory from $(x_{0}, v_{0}) = (0, 0)$ accelerates under
$u_{k} = 0.1$ for several steps, coasts, then decelerates to
arrive near the target. Discuss the trade-offs: a longer horizon
$H$ produces smoother actuation but more compute per step;
softening the actuator bound (relaxing to a penalty) speeds the
QP solve but degrades safety; in practice the engineer ships the
shortest horizon that delivers acceptable tracking and uses warm-
starts (initialising each QP from the previous one) for an
order-of-magnitude speedup.

\begin{exercise}
Design a knapsack instance with $n = 20$ items of randomly
generated weights and values, and a capacity of $50\%$ of the
total weight. Solve as an integer program (using a library
solver) and as the LP relaxation. Report the optimality gap
between the two solutions.
\end{exercise}

\paragraph{Discussion.}
The LP relaxation of the 0/1 knapsack is solved by sorting items
by value-to-weight ratio and filling the knapsack greedily, with a
fractional item at the end. The LP optimum upper-bounds the
integer optimum. Typical gaps on uniformly random instances are
$1\%$ to $10\%$; tighter gaps arise when one item dominates the
sort. Discuss: when the gap is small, the LP relaxation already
gives a tight bound and branch-and-bound terminates fast; when
the gap is large (e.g.\ the LP would pack many fractional items)
the search is slow, and cutting planes or branching strategies
that focus on the most fractional variables become decisive.

\subsection*{Diagnosis and reverse engineering}\pathtag{standard}

\begin{exercise}
A colleague reports that an SGD run has converged to a training
loss of $0.05$ but the test loss is $0.50$. Identify the most
likely failure mode and propose three diagnostic measurements to
distinguish overfitting, label noise, and a learning-rate
problem.
\end{exercise}

\paragraph{Solution sketch.}
Most likely overfitting (the optimiser succeeded on the wrong
objective, as in Section 15.8). Three diagnostics:
(i) train on a $50\%$ subset and report test loss versus training-
set size; if the gap shrinks with more data, the model is
overfitting; (ii) compute the noise floor by relabelling a
small fraction of the training set and retraining; if the
training loss falls below the noise floor, the model is
memorising; (iii) reduce model capacity (smaller architecture,
stronger weight decay, dropout) and check whether the gap closes.
If none of the three change the gap, suspect label noise or
distribution shift between train and test.

\begin{exercise}
A colleague reports that gradient descent on a smooth convex
objective has not converged after $10^{5}$ iterations and
proposes increasing the iteration budget. Identify three checks
the engineer should run before paying for more iterations
(condition number, step size, smoothness assumption).
\end{exercise}

\paragraph{Solution sketch.}
(i) Estimate $\kappa = L/\mu$ by power iteration on the Hessian
(or on $A^{\transpose} A$ for least squares). If $\kappa > 10^{6}$,
no realistic iteration budget will reach the optimum with plain
gradient descent; switch to L-BFGS or precondition. (ii) Check
that the step size $\eta$ satisfies $\eta L < 2$; if not, the
iteration is oscillating. Run a backtracking line search to verify
the Armijo condition holds at each step. (iii) Verify the
smoothness assumption by computing $\norm{\grad f(\vect{x}_{k+1})
- \grad f(\vect{x}_{k})} / \norm{\vect{x}_{k+1} - \vect{x}_{k}}$
along the iterates; if the empirical Lipschitz constant exceeds
the assumed $L$, the convergence theorem does not apply and the
step size is too aggressive. Any of the three findings reframes
the problem; only after all three pass should the budget be
increased.

\begin{exercise}
A library LP solver returns "infeasible" on a problem the
engineer believes is feasible. Identify three causes (numerical
ill-conditioning, a typo in a constraint, an inconsistency
between equality and inequality constraints) and propose a
diagnostic protocol.
\end{exercise}

\paragraph{Solution sketch.}
Diagnostic protocol: (i) Compute an irreducible infeasible subset
(IIS), supported by most modern solvers as \verb|model.computeIIS()|;
this returns a minimal set of constraints whose removal restores
feasibility. Inspect each. (ii) If IIS reveals two contradictory
constraints, look for a sign error in the data or a unit mismatch
(metres vs centimetres, dollars vs cents). (iii) If IIS reveals
one constraint that should be trivially satisfiable, inspect the
right-hand side and the variable bounds for typos. Two structural
causes: numerical ill-conditioning (the solver's working tolerance
exceeds the constraint magnitude; rescale the problem so each
constraint and each variable is $O(1)$); and a mismatch between
equality and inequality blocks (an equality $A \vect{x} = \vect{b}$
combined with an inequality $A \vect{x} \leq \vect{b} - \epsilon$
is infeasible by construction).

\begin{exercise}
An MIP solver reports a $5\%$ optimality gap after the time
budget expires. Identify what the gap means, what the engineer
can report on the basis of the partial solve, and the three
options (extend time, tighten formulation, accept the bound) for
proceeding.
\end{exercise}

\paragraph{Solution sketch.}
The $5\%$ gap is $(\text{incumbent} - \text{best LP bound}) /
\text{incumbent}$. The incumbent is the best integer-feasible
solution found; the best LP bound is a lower bound (for
minimisation) certified by the unexplored subtree's LP
relaxations. The engineer can report: "Best known feasible
solution has value $V$; the true optimum lies in $[V \cdot 0.95,
V]$." Three options: (i) extend the time budget if the gap is
closing steadily (observable by logging the bound over time);
(ii) tighten the formulation (add cutting planes, valid
inequalities, stronger branching rules) which often produces an
order-of-magnitude speedup; (iii) accept the $5\%$ bound as the
engineering answer if the downstream decision is not sensitive to
optimality at that level. The third option is often the correct
engineering call: $5\%$ on a manufacturing-cost objective is below
the noise in the input data.

\subsection*{Failure analysis}\pathtag{core}

\begin{exercise}
A polynomial regression of degree $d = 11$ on $N = 12$ data points
fits the training data exactly but extrapolates wildly between
points. Argue, in 200 to 300 words, what the optimiser did
correctly and what the formulation got wrong, and propose a
specific remediation referencing Section 15.8.
\end{exercise}

\paragraph{Solution sketch.}
The optimiser correctly minimised the stated training-set MSE to
zero; with $d + 1 = 12$ free coefficients and $N = 12$ data
points, the interpolating polynomial exists and is unique. The
formulation got wrong its choice of objective: training-set fit is
a proxy for unseen-data fit, and the two diverge in the
high-capacity regime. Remediation A: regularise. Replace the
objective with $\mathcal{L}_{\text{train}} + \lambda
\norm{\boldsymbol{\theta}}_{2}^{2}$ for $\lambda$ chosen by cross-
validation; the ridge regression of Section 15.8 worked example
demonstrates this. Remediation B: reduce capacity. A lower-degree
polynomial (say $d = 3$ for $N = 12$) has fewer parameters than
data points and cannot interpolate, hence cannot oscillate.
Remediation C: change the basis. B-splines or low-order piecewise
polynomials trade global oscillation for local smoothness. The
common thread: the optimiser is doing its job; the engineer
re-specifies the job. The diagnostic in Section 15.8 (state the
objective in plain English, ask what literal optimum the
formulation rewards) catches this failure before it ships.

\begin{exercise}
A trading strategy is optimised on $10$ years of historical
backtest data and shows a $40\%$ annualised return. In live
deployment over 12 months, the strategy returns $-5\%$. Argue,
in 200 to 300 words, the most likely failure mode, citing the
relevant Section 15.8 principle.
\end{exercise}

\paragraph{Solution sketch.}
Most likely failure: backtest overfit, the trading equivalent of
the polynomial-regression failure of Section 15.8. The optimiser
correctly maximised the stated backtest objective; the
formulation got wrong its assumption that backtest performance
generalises to live deployment. The proxy (historical returns
under the strategy) was optimised aggressively, and the
optimisation has moved the strategy to exploit any peculiarity of
the historical sample, including peculiarities that will not
recur. By Goodhart's law (\textcite{text:v2ch15-goodhart1975}), once
backtest return became the target, it ceased to be a measure of
strategy quality. Remediations: (i) walk-forward validation, in
which the strategy is fit on one window and evaluated on the next,
with the windows rolling forward; (ii) regularisation of the
strategy's free parameters; (iii) penalisation of strategy
complexity (number of parameters, number of rules); (iv) explicit
out-of-sample reserve before any tuning is allowed. None of these
guarantee live performance; all of them reduce the gap between
backtest and live. The engineer's responsibility is to anticipate
the gap before the optimiser exploits it.

\begin{exercise}
A neural network is trained for 1000 epochs and the training
loss is reported to seven decimal places, but the validation
loss has plateaued and slightly increased over the last 200
epochs. Argue, in 200 to 300 words, why early stopping is the
working remedy, and identify what the optimiser optimised that
the engineer did not want.
\end{exercise}

\paragraph{Solution sketch.}
The optimiser minimised training-set loss across the full
$1000$-epoch trajectory; the validation loss plateau and slight
rise over the last $200$ epochs is the empirical signature of
overfitting beyond the validation optimum. Early stopping is the
working remedy because it freezes the parameters at the iterate
that minimises the validation loss, which is the iterate that
generalises best by the held-out estimate of generalisation error.
The engineer wanted the model that minimises true unseen-data
loss; the optimiser delivered the model that minimises training-set
loss after $1000$ epochs. Reporting the training loss to seven
decimal places signals false precision: the meaningful quantity
is the validation loss at the early-stopping iterate, with its
sampling uncertainty stated. Implementation: track the validation
loss at each epoch, snapshot the model when the validation loss
improves, restore the snapshot at the end. The procedure adds
zero compute over the unguarded run and yields the
generalisation-optimal model the engineer wanted.

\begin{exercise}
A factory's quality program targets a defect rate of $\leq 1\%$
and ties bonuses to it. Defect rates fall to $0.3\%$ but customer
complaints about product reliability rise. Argue, in 200 to 300
words, how Goodhart's law applies and propose a corrected
incentive structure that rewards what the engineer actually wants.
\end{exercise}

\paragraph{Solution sketch.}
Goodhart's law applies directly. The measured defect rate, once
installed as a bonus target, ceased to be a measure of product
reliability; the workforce optimises the measured rate, which can
be moved (lower the inspection threshold, reclassify failures as
maintenance items, push marginal units through without flagging)
without improving the underlying reliability. The rise in customer
complaints is the symptom: the engineering goal (reliable
products) and the measured target (low in-factory defect rate)
have diverged. Corrected incentive structure: target a composite
of in-factory defect rate, in-warranty failure rate, and customer-
reported failure rate over the first six months, with each
component independently auditable and each tied to a portion of
the bonus. The composite is harder to game because gaming one
component (e.g.\ understating in-factory defects) leaves the
others (in-warranty failures, customer reports) detecting the same
underlying problem. A second-line defence: independent
quality-assurance audits, conducted by inspectors whose
compensation depends on detection rate rather than on the measured
defect rate of the factory under audit. The general lesson of
Section 15.8: an installed target is, by construction, a target,
and any gap between target and intent will be exploited unless
the engineer closes the gap by reformulation.

\subsection*{Judgment}\pathtag{standard}

\begin{exercise}
A colleague proposes solving a nonconvex problem by running Adam
once from a single initial point and reporting the result.
Argue, in 200 to 300 words, when this is defensible (large
problem, weak local minima, no time for restarts) and when it is
not (small enough for restarts, multimodal objective, downstream
decisions sensitive to local minima).
\end{exercise}

\paragraph{Discussion.}
Defensible when (i) the problem is too large for restarts to be
affordable (deep-learning models with millions of parameters where
a single training run costs hours of GPU time), (ii) empirical
practice in the relevant subfield has established that local
minima of the loss landscape are weakly distinguished by their
generalisation behaviour (a substantial body of evidence supports
this for modern neural networks, current as of 2026), or (iii) the
downstream decision is insensitive to which local minimum is
found because the engineering goal is "any reasonable model" not
"the best possible model." Not defensible when (i) the problem is
small enough that even ten random restarts is cheap (any problem
with $n < 10^{4}$ parameters and a fast loss evaluation), (ii) the
objective is known to be multimodal with structurally different
basins (mixture-model fits, phylogenetic trees, combinatorial
relaxations), or (iii) the downstream decision is sensitive to
the final value of the optimum (safety-critical inference, dollar-
denominated decisions in operations research, scientific
parameter estimation). The general rule: report the local-versus-
global guarantee that has actually been earned, with the number
of restarts and the spread of converged losses, and let the
downstream user judge whether the guarantee is sufficient.

\begin{exercise}
A vendor sells a "global optimisation" package that claims to
find the global optimum of any nonconvex problem in polynomial
time. Argue, in one paragraph, why the claim is incompatible
with NP-hardness for any plausible problem class, and propose
the questions an engineer should ask before evaluating the
package.
\end{exercise}

\paragraph{Discussion.}
A polynomial-time global optimiser for any nonconvex problem would
imply $P = NP$ (since the integer-programming feasibility problem
is a special case, and Cook 1971 and Karp 1972 established NP-
completeness for the family). The vendor either restricts the
problem class without saying so, runs in polynomial time on every
problem they tested but is exponential on at least some problem
class, uses a heuristic that returns local optima with no global
guarantee, or is misrepresenting the package. Questions to ask:
(i) On what problem class does the polynomial-time guarantee
hold? Get the formal statement. (ii) What is the runtime
complexity in the worst case? In the average case under what
distribution? (iii) What benchmarks does the package solve, and
what is its performance against state-of-the-art solvers
(Gurobi, Mosek, CPLEX) on those benchmarks? (iv) Will the vendor
provide a no-cost trial on the engineer's actual problem
instances? Any vendor who refuses (iv) has failed the
engineering acceptance test.

\begin{exercise}
\textcite{text:boyd-vandenberghe2004} state that "almost any
optimisation problem with linear constraints can be modelled as
a convex program by an appropriate change of variables." Argue,
in one paragraph, why this is a strong but not universal claim,
and identify a class of linearly constrained problems (the
integer ones) where the change of variables does not exist.
\end{exercise}

\paragraph{Discussion.}
The claim is strong because it covers a large fraction of the
practical optimisation problems an engineer encounters:
geometric programming, $\log$-$\sum$-$\exp$ objectives, fractional
linear objectives, generalised polynomial inequalities all admit
convex change-of-variables that preserve linearity of constraints.
The claim is not universal: any problem with an integer constraint
$x_{i} \in \Z$ is linearly constrained (the integer requirement is
a constraint set, not a nonconvex function in the objective), but
no change of variables turns the integer lattice into a convex
set; the lattice fails the segment test by construction. The
integer-programming class is therefore a structurally separate
category, and the working strategy is relaxation plus branch-and-
bound rather than a convex reformulation. The same line of
reasoning excludes mixed-integer programs, combinatorial problems
defined on discrete graphs, and problems with cardinality
constraints (at most $k$ variables nonzero).

\begin{exercise}
\textcite{text:strang2016}'s treatment of least squares
emphasises the geometric content (orthogonal projection onto the
column space) over the optimisation content (gradient zero of a
quadratic). Argue, in one paragraph, why both viewpoints are
essential and how each illuminates the other.
\end{exercise}

\paragraph{Discussion.}
The geometric viewpoint reveals what the optimum is: the
orthogonal projection of $\vect{b}$ onto the column space of $A$,
unique whenever $A$ has full column rank. The optimisation
viewpoint reveals how to compute it: set $\grad f = A^{\transpose}
(A \vect{x} - \vect{b}) = \vect{0}$, giving the normal equations
$A^{\transpose} A \vect{x} = A^{\transpose} \vect{b}$. The
geometry explains uniqueness (the column space is a linear
subspace, projection onto it is single-valued); the optimisation
explains computation (the normal equations are a linear system,
solvable by Cholesky or QR in $O(n^{3})$). Each viewpoint
illuminates the other: the geometric reading shows why the normal
equations have a solution whenever $A^{\transpose} A$ is invertible
(the column space is the same as the column space of
$A^{\transpose} A$), and the optimisation reading shows why
iterative methods (gradient descent on the residual-squared
objective) converge to the projection. The two viewpoints meet at
the singular value decomposition, where geometry, optimisation,
and numerical stability converge in a single object.

\begin{exercise}
A practitioner uses Adam as the default optimiser for every
problem they encounter. Argue, in 200 to 300 words, when this
default is reasonable (large nonconvex deep-learning losses with
unknown smoothness) and when it is irresponsible (small smooth
convex problems where L-BFGS is faster, or constrained problems
where a constrained solver is required).
\end{exercise}

\paragraph{Discussion.}
Reasonable when (i) the objective is a large nonconvex deep-
learning loss with millions to billions of parameters, where the
per-iteration cost of L-BFGS or Newton is prohibitive and the
per-coordinate adaptation of Adam handles the wide range of
gradient scales typical in deep networks; (ii) the smoothness
constant $L$ is unknown and would require a tuning sweep to
estimate, which Adam's adaptive step sizes amortise; or (iii)
the engineering goal is "first working baseline" rather than
"smallest possible loss." Irresponsible when (i) the problem is
small and smooth and convex, where L-BFGS converges
super-linearly with a few-line setup and Adam's $O(1/\sqrt{K})$
rate is two orders of magnitude slower; (ii) the problem is
constrained, where Adam has no constraint-handling mechanism and
the engineer should use a projected-gradient method, an interior-
point method, or a modelling layer such as CVXPY that calls a
constrained solver; (iii) the problem has known structure
(quadratic, conic, semidefinite, integer) that a specialised
solver exploits with $10\times$ to $10^{6}\times$ speedups over a
general first-order method. The professional rule: choose the
method that fits the problem class, document the choice, and
default to the method only after the class has been recognised.
The principle in Section 15.5 is the working selection rule.

